% !TEX program = lualatex
% ---------------------------------------------------------------------------------------------
%  Readback note 17: The exchange theorem, one square of lifts, and the triangle of lifts and descents
%
%  Build:  python3 tools/rb.py note <slug>
%
%  One unit of a Lean artifact, rewritten as ordinary mathematics. The prose is hand-written. Lean names,
%  line ranges and counts come from data.tex; quoted Lean statements are copied from the artifact by the build.
%
%  Rules of the text: standard terminology; every symbol is explained where it first appears and listed again
%  in the notation at the end; no letter has two meanings; a proof is complete or says what it takes for granted;
%  Lean names appear only in the LEAN lines and in the concordance.
%
%  The signature rule (the build enforces it): a function, relation or predicate that the artifact defines is
%  written only through its notation macro, which the build makes from the spec (`notation`) with one slot for
%  each explicit argument of the Lean definition. No argument is fixed and dropped, also not "for brevity" and
%  not for one section. Every argument and every hypothesis of a stated declaration stands inside its
%  statement, or in a Setting that the statement cites by number; only implicit and instance arguments (the
%  ambient objects) may stand in the conventions.
% ---------------------------------------------------------------------------------------------
\documentclass[11pt,a4paper]{article}
\ifdefined\pdfvariable\pdfvariable suppressoptionalinfo \numexpr 1+2+4+8\relax\fi
\usepackage[a4paper,hmargin=35mm,top=22mm,bottom=27mm,footskip=12mm,nohead]{geometry}

\usepackage{fontspec}
\directlua{
  luaotfload.add_fallback("gfallback", {"STIXTwoMath-Regular.otf:mode=node;", "DejaVuSans.ttf:mode=node;"})
  luaotfload.add_fallback("gmonofallback", {"DejaVuSansMono.ttf:mode=node;", "STIXTwoMath-Regular.otf:mode=node;"})
}
\setmainfont{NewCM10}[Extension=.otf,UprightFont=*-Book,ItalicFont=*-BookItalic,
  BoldFont=*-Bold,BoldItalicFont=*-BoldItalic,SmallCapsFont=*-Book,SmallCapsFeatures={Letters=SmallCaps},
  RawFeature={fallback=gfallback}]
\setsansfont{NewCMSans10}[Extension=.otf,UprightFont=*-Book,ItalicFont=*-BookOblique,
  BoldFont=*-Bold,BoldItalicFont=*-BoldOblique,RawFeature={fallback=gfallback}]
\setmonofont{JetBrainsMono}[Extension=.otf,UprightFont=*-Regular,ItalicFont=*-Italic,
  BoldFont=*-SemiBold,BoldItalicFont=*-SemiBoldItalic,Scale=MatchLowercase,
  RawFeature={-calt;-liga;fallback=gmonofallback}]
\usepackage{amsmath}
\usepackage{amsthm}
\usepackage[math-style=TeX]{unicode-math}
\setmathfont{NewCMMath-Book.otf}
\emergencystretch=2.5em\relax

\usepackage{xcolor}
\usepackage{tikz}
\usetikzlibrary{calc,arrows.meta,positioning,backgrounds}
\usepackage[protrusion=true,expansion=true]{microtype}
\usepackage{enumitem}
\usepackage{needspace}
\usepackage{booktabs}
\usepackage{array}
\usepackage{longtable}
\usepackage{fancyhdr}
\usepackage{lastpage}
\usepackage{multicol}
\usepackage[hidelinks,pdfusetitle]{hyperref}

\definecolor{ink}{HTML}{0F172A}
\definecolor{inksoft}{HTML}{475569}
\definecolor{inkfaint}{HTML}{94A3B8}
\definecolor{rulegrey}{HTML}{CBD5E1}
\definecolor{arith}{HTML}{C2410C}
\definecolor{est}{HTML}{7C3AED}
\definecolor{ana}{HTML}{0891B2}
\definecolor{ok}{HTML}{047857}
\definecolor{leancol}{HTML}{3730A3}

\title{The exchange theorem, one square of lifts, and the triangle of lifts and descents}
\author{grh-compact}
\date{}

% ------------------------------------------------------------------------------------- macros
\newcommand{\leanfont}{\ttfamily\fontsize{8.6}{10.5}\selectfont\color{leancol}}
\newcommand{\lean}[1]{\csname LN@#1\endcsname}
\newcommand{\fillin}[1]{{\color{arith}\itshape[#1]}}  % a place of the template that is not written yet; the build reports every one that is left
\newcommand{\leanquote}[1]{\csname LQ@#1\endcsname}
% generated by rb.py note; do not edit
\newcommand{\ConcDecls}{129}
\newcommand{\Pin}{72fead10}
\newcommand{\PinText}{that of commit {\leanfont 72fead10}}
\newcommand{\ArtifactName}{grh-compact}
\expandafter\newcommand\csname LN@maptrivialLift\endcsname{\href{https://github.com/lsernau/grh-lean/blob/72fead102f663d5c7a407964a13dc0cf9d977207/grh-compact/src/New/Exchange/AlgebraicPartConjugates.lean\#L13-L47}{{\leanfont LValueData.\allowbreak{}map\_trivialLift}}}
\expandafter\newcommand\csname LN@lognormmapnormsqtrivialLift\endcsname{\href{https://github.com/lsernau/grh-lean/blob/72fead102f663d5c7a407964a13dc0cf9d977207/grh-compact/src/New/Exchange/AlgebraicPartConjugates.lean\#L52-L81}{{\leanfont Data.\allowbreak{}Pair.\allowbreak{}log\_norm\_map\_norm\_sq\_trivialLift}}}
\expandafter\newcommand\csname LN@cardYle\endcsname{\href{https://github.com/lsernau/grh-lean/blob/72fead102f663d5c7a407964a13dc0cf9d977207/grh-compact/src/New/Exchange/AlgebraicPartDefs.lean\#L20-L21}{{\leanfont Data.\allowbreak{}card\_Y\_le}}}
\expandafter\newcommand\csname LN@isHCParamY\endcsname{\href{https://github.com/lsernau/grh-lean/blob/72fead102f663d5c7a407964a13dc0cf9d977207/grh-compact/src/New/Exchange/AlgebraicPartDefs.lean\#L25-L38}{{\leanfont Data.\allowbreak{}isHCParam\_Y}}}
\expandafter\newcommand\csname LN@cardV\endcsname{\href{https://github.com/lsernau/grh-lean/blob/72fead102f663d5c7a407964a13dc0cf9d977207/grh-compact/src/New/Exchange/AlgebraicPartDefs.lean\#L40-L40}{{\leanfont Data.\allowbreak{}card\_V}}}
\expandafter\newcommand\csname LN@cardfilterbasisNumber\endcsname{\href{https://github.com/lsernau/grh-lean/blob/72fead102f663d5c7a407964a13dc0cf9d977207/grh-compact/src/New/Exchange/AlgebraicPartDefs.lean\#L42-L44}{{\leanfont Data.\allowbreak{}card\_filter\_basisNumber}}}
\expandafter\newcommand\csname LN@varsigmastrictAnti\endcsname{\href{https://github.com/lsernau/grh-lean/blob/72fead102f663d5c7a407964a13dc0cf9d977207/grh-compact/src/New/Exchange/AlgebraicPartDefs.lean\#L46-L47}{{\leanfont Data.\allowbreak{}varsigma\_strictAnti}}}
\expandafter\newcommand\csname LN@paramWltiff\endcsname{\href{https://github.com/lsernau/grh-lean/blob/72fead102f663d5c7a407964a13dc0cf9d977207/grh-compact/src/New/Exchange/AlgebraicPartDefs.lean\#L49-L50}{{\leanfont Data.\allowbreak{}paramW\_lt\_iff}}}
\expandafter\newcommand\csname LN@cardnegsignsW\endcsname{\href{https://github.com/lsernau/grh-lean/blob/72fead102f663d5c7a407964a13dc0cf9d977207/grh-compact/src/New/Exchange/AlgebraicPartDefs.lean\#L53-L54}{{\leanfont Data.\allowbreak{}card\_neg\_signsW}}}
\expandafter\newcommand\csname LN@fitsParamsauxPair\endcsname{\href{https://github.com/lsernau/grh-lean/blob/72fead102f663d5c7a407964a13dc0cf9d977207/grh-compact/src/New/Exchange/AlgebraicPartDefs.lean\#L58-L71}{{\leanfont Data.\allowbreak{}fitsParams\_auxPair}}}
\expandafter\newcommand\csname LN@innereqpetersson\endcsname{\href{https://github.com/lsernau/grh-lean/blob/72fead102f663d5c7a407964a13dc0cf9d977207/grh-compact/src/New/Exchange/AlgebraicPartDefs.lean\#L76-L82}{{\leanfont LevelPlaces.\allowbreak{}LevelForm.\allowbreak{}inner\_eq\_petersson}}}
\expandafter\newcommand\csname LN@chosenIndexEmb\endcsname{\href{https://github.com/lsernau/grh-lean/blob/72fead102f663d5c7a407964a13dc0cf9d977207/grh-compact/src/New/Exchange/AlgebraicPartDefs.lean\#L86-L86}{{\leanfont chosenIndexEmb}}}
\expandafter\newcommand\csname LN@LValueData\endcsname{\href{https://github.com/lsernau/grh-lean/blob/72fead102f663d5c7a407964a13dc0cf9d977207/grh-compact/src/New/Exchange/AlgebraicPartDefs.lean\#L91-L98}{{\leanfont LValueData}}}
\expandafter\newcommand\csname LN@auxPair\endcsname{\href{https://github.com/lsernau/grh-lean/blob/72fead102f663d5c7a407964a13dc0cf9d977207/grh-compact/src/New/Exchange/AlgebraicPartDefs.lean\#L101-L103}{{\leanfont LValueData.\allowbreak{}auxPair}}}
\expandafter\newcommand\csname LN@kernelValue\endcsname{\href{https://github.com/lsernau/grh-lean/blob/72fead102f663d5c7a407964a13dc0cf9d977207/grh-compact/src/New/Exchange/AlgebraicPartDefs.lean\#L110-L111}{{\leanfont LValueData.\allowbreak{}kernelValue}}}
\expandafter\newcommand\csname LN@isLiftWeight\endcsname{\href{https://github.com/lsernau/grh-lean/blob/72fead102f663d5c7a407964a13dc0cf9d977207/grh-compact/src/New/Exchange/AlgebraicPartDefs.lean\#L117-L118}{{\leanfont LValueData.\allowbreak{}isLiftWeight}}}
\expandafter\newcommand\csname LN@ahasWeight\endcsname{\href{https://github.com/lsernau/grh-lean/blob/72fead102f663d5c7a407964a13dc0cf9d977207/grh-compact/src/New/Exchange/AlgebraicPartDefs.lean\#L120-L121}{{\leanfont LValueData.\allowbreak{}a\_hasWeight}}}
\expandafter\newcommand\csname LN@bhasWeight\endcsname{\href{https://github.com/lsernau/grh-lean/blob/72fead102f663d5c7a407964a13dc0cf9d977207/grh-compact/src/New/Exchange/AlgebraicPartDefs.lean\#L123-L124}{{\leanfont LValueData.\allowbreak{}b\_hasWeight}}}
\expandafter\newcommand\csname LN@liftKernel\endcsname{\href{https://github.com/lsernau/grh-lean/blob/72fead102f663d5c7a407964a13dc0cf9d977207/grh-compact/src/New/Exchange/AlgebraicPartDefs.lean\#L128-L129}{{\leanfont LValueData.\allowbreak{}liftKernel}}}
\expandafter\newcommand\csname LN@conjWeylRepeqone\endcsname{\href{https://github.com/lsernau/grh-lean/blob/72fead102f663d5c7a407964a13dc0cf9d977207/grh-compact/src/New/Exchange/AlgebraicPartDefs.lean\#L132-L133}{{\leanfont LValueData.\allowbreak{}conjWeylRep\_eq\_one}}}
\expandafter\newcommand\csname LN@kernelForm\endcsname{\href{https://github.com/lsernau/grh-lean/blob/72fead102f663d5c7a407964a13dc0cf9d977207/grh-compact/src/New/Exchange/AlgebraicPartDefs.lean\#L137-L142}{{\leanfont LValueData.\allowbreak{}kernelForm}}}
\expandafter\newcommand\csname LN@chosenLevel\endcsname{\href{https://github.com/lsernau/grh-lean/blob/72fead102f663d5c7a407964a13dc0cf9d977207/grh-compact/src/New/Exchange/AlgebraicPartDefs.lean\#L146-L146}{{\leanfont LValueData.\allowbreak{}chosenLevel}}}
\expandafter\newcommand\csname LN@trivialLift\endcsname{\href{https://github.com/lsernau/grh-lean/blob/72fead102f663d5c7a407964a13dc0cf9d977207/grh-compact/src/New/Exchange/AlgebraicPartDefs.lean\#L150-L150}{{\leanfont LValueData.\allowbreak{}trivialLift}}}
\expandafter\newcommand\csname LN@kernelForminvariant\endcsname{\href{https://github.com/lsernau/grh-lean/blob/72fead102f663d5c7a407964a13dc0cf9d977207/grh-compact/src/New/Exchange/AlgebraicPartDefs.lean\#L155-L159}{{\leanfont LValueData.\allowbreak{}kernelForm\_invariant}}}
\expandafter\newcommand\csname LN@trivialLifteqstdClassSumtranslate\endcsname{\href{https://github.com/lsernau/grh-lean/blob/72fead102f663d5c7a407964a13dc0cf9d977207/grh-compact/src/New/Exchange/AlgebraicPartDefs.lean\#L161-L164}{{\leanfont LValueData.\allowbreak{}trivialLift\_eq\_stdClassSum\_translate}}}
\expandafter\newcommand\csname LN@toPolyeqrenameharmonicPoly\endcsname{\href{https://github.com/lsernau/grh-lean/blob/72fead102f663d5c7a407964a13dc0cf9d977207/grh-compact/src/New/Exchange/AlgebraicPartDefs.lean\#L170-L175}{{\leanfont LValueData.\allowbreak{}toPoly\_eq\_rename\_harmonicPoly}}}
\expandafter\newcommand\csname LN@harmonicPolyeqschurCoeffPoly\endcsname{\href{https://github.com/lsernau/grh-lean/blob/72fead102f663d5c7a407964a13dc0cf9d977207/grh-compact/src/New/Exchange/AlgebraicPartDefs.lean\#L178-L183}{{\leanfont LValueData.\allowbreak{}harmonicPoly\_eq\_schurCoeffPoly}}}
\expandafter\newcommand\csname LN@thetapairWeilfockDatum\endcsname{\href{https://github.com/lsernau/grh-lean/blob/72fead102f663d5c7a407964a13dc0cf9d977207/grh-compact/src/New/Exchange/AlgebraicPartDefs.lean\#L185-L190}{{\leanfont LValueData.\allowbreak{}theta\_pairWeil\_fockDatum}}}
\expandafter\newcommand\csname LN@kernelFormofjetDegeqzero\endcsname{\href{https://github.com/lsernau/grh-lean/blob/72fead102f663d5c7a407964a13dc0cf9d977207/grh-compact/src/New/Exchange/AlgebraicPartDefs.lean\#L195-L202}{{\leanfont LValueData.\allowbreak{}kernelForm\_of\_jetDeg\_eq\_zero}}}
\expandafter\newcommand\csname LN@thetaeqinnerkernelForm\endcsname{\href{https://github.com/lsernau/grh-lean/blob/72fead102f663d5c7a407964a13dc0cf9d977207/grh-compact/src/New/Exchange/AlgebraicPartDefs.lean\#L204-L206}{{\leanfont LValueData.\allowbreak{}theta\_eq\_inner\_kernelForm}}}
\expandafter\newcommand\csname LN@innertrivialLift\endcsname{\href{https://github.com/lsernau/grh-lean/blob/72fead102f663d5c7a407964a13dc0cf9d977207/grh-compact/src/New/Exchange/AlgebraicPartDefs.lean\#L211-L213}{{\leanfont LValueData.\allowbreak{}inner\_trivialLift}}}
\expandafter\newcommand\csname LN@hookNumbereqprodfactorial\endcsname{\href{https://github.com/lsernau/grh-lean/blob/72fead102f663d5c7a407964a13dc0cf9d977207/grh-compact/src/New/Exchange/AlgebraicPartDefs.lean\#L218-L222}{{\leanfont LValueData.\allowbreak{}hookNumber\_eq\_prod\_factorial}}}
\expandafter\newcommand\csname LN@HasLiftRank\endcsname{\href{https://github.com/lsernau/grh-lean/blob/72fead102f663d5c7a407964a13dc0cf9d977207/grh-compact/src/New/Exchange/AlgebraicPartDefs.lean\#L227-L227}{{\leanfont Data.\allowbreak{}Pair.\allowbreak{}HasLiftRank}}}
\expandafter\newcommand\csname LN@IsFlat\endcsname{\href{https://github.com/lsernau/grh-lean/blob/72fead102f663d5c7a407964a13dc0cf9d977207/grh-compact/src/New/Exchange/AlgebraicPartDefs.lean\#L231-L231}{{\leanfont Data.\allowbreak{}Pair.\allowbreak{}IsFlat}}}
\expandafter\newcommand\csname LN@denomBound\endcsname{\href{https://github.com/lsernau/grh-lean/blob/72fead102f663d5c7a407964a13dc0cf9d977207/grh-compact/src/New/Exchange/AlgebraicPartDefs.lean\#L235-L236}{{\leanfont Data.\allowbreak{}Pair.\allowbreak{}denomBound}}}
\expandafter\newcommand\csname LN@lValueData\endcsname{\href{https://github.com/lsernau/grh-lean/blob/72fead102f663d5c7a407964a13dc0cf9d977207/grh-compact/src/New/Exchange/AlgebraicPartDefs.lean\#L242-L244}{{\leanfont Data.\allowbreak{}Pair.\allowbreak{}lValueData}}}
\expandafter\newcommand\csname LN@trivialLifta\endcsname{\href{https://github.com/lsernau/grh-lean/blob/72fead102f663d5c7a407964a13dc0cf9d977207/grh-compact/src/New/Exchange/AlgebraicPartDefs.lean\#L249-L250}{{\leanfont Data.\allowbreak{}Pair.\allowbreak{}trivialLift}}}
\expandafter\newcommand\csname LN@intAlgPart\endcsname{\href{https://github.com/lsernau/grh-lean/blob/72fead102f663d5c7a407964a13dc0cf9d977207/grh-compact/src/New/Exchange/AlgebraicPartDefs.lean\#L254-L255}{{\leanfont Data.\allowbreak{}Pair.\allowbreak{}intAlgPart}}}
\expandafter\newcommand\csname LN@stdLevelIndexlastResLeveleq\endcsname{\href{https://github.com/lsernau/grh-lean/blob/72fead102f663d5c7a407964a13dc0cf9d977207/grh-compact/src/New/Exchange/ChainNorms.lean\#L19-L37}{{\leanfont Data.\allowbreak{}stdLevelIndex\_lastResLevel\_eq}}}
\expandafter\newcommand\csname LN@inneriotaAtLevel\endcsname{\href{https://github.com/lsernau/grh-lean/blob/72fead102f663d5c7a407964a13dc0cf9d977207/grh-compact/src/New/Exchange/ChainNorms.lean\#L42-L47}{{\leanfont Data.\allowbreak{}inner\_iotaAtLevel}}}
\expandafter\newcommand\csname LN@iotaAtLevelinjective\endcsname{\href{https://github.com/lsernau/grh-lean/blob/72fead102f663d5c7a407964a13dc0cf9d977207/grh-compact/src/New/Exchange/ChainNorms.lean\#L53-L55}{{\leanfont Data.\allowbreak{}iotaAtLevel\_injective}}}
\expandafter\newcommand\csname LN@CostsSquares\endcsname{\href{https://github.com/lsernau/grh-lean/blob/72fead102f663d5c7a407964a13dc0cf9d977207/grh-compact/src/New/Exchange/CostOfSquares.lean\#L17-L18}{{\leanfont Data.\allowbreak{}CostsSquares}}}
\expandafter\newcommand\csname LN@mul\endcsname{\href{https://github.com/lsernau/grh-lean/blob/72fead102f663d5c7a407964a13dc0cf9d977207/grh-compact/src/New/Exchange/CostOfSquares.lean\#L22-L24}{{\leanfont Data.\allowbreak{}CostsSquares.\allowbreak{}mul}}}
\expandafter\newcommand\csname LN@card\endcsname{\href{https://github.com/lsernau/grh-lean/blob/72fead102f663d5c7a407964a13dc0cf9d977207/grh-compact/src/New/Exchange/CostOfSquares.lean\#L28-L29}{{\leanfont Data.\allowbreak{}card\_𝔪}}}
\expandafter\newcommand\csname LN@ofnormsqeqmulcardzpow\endcsname{\href{https://github.com/lsernau/grh-lean/blob/72fead102f663d5c7a407964a13dc0cf9d977207/grh-compact/src/New/Exchange/CostOfSquares.lean\#L33-L36}{{\leanfont Data.\allowbreak{}CostsSquares.\allowbreak{}of\_norm\_sq\_eq\_mul\_card\_zpow}}}
\expandafter\newcommand\csname LN@one\endcsname{\href{https://github.com/lsernau/grh-lean/blob/72fead102f663d5c7a407964a13dc0cf9d977207/grh-compact/src/New/Exchange/CostOfSquares.lean\#L38-L38}{{\leanfont Data.\allowbreak{}CostsSquares.\allowbreak{}one}}}
\expandafter\newcommand\csname LN@isPrimitiveMulCharofisTameFor\endcsname{\href{https://github.com/lsernau/grh-lean/blob/72fead102f663d5c7a407964a13dc0cf9d977207/grh-compact/src/New/Exchange/CostOfSquares.lean\#L42-L49}{{\leanfont Data.\allowbreak{}isPrimitiveMulChar\_of\_isTameFor}}}
\expandafter\newcommand\csname LN@existsisTameForcostsSquaresone\endcsname{\href{https://github.com/lsernau/grh-lean/blob/72fead102f663d5c7a407964a13dc0cf9d977207/grh-compact/src/New/Exchange/CostOfSquares.lean\#L54-L63}{{\leanfont Data.\allowbreak{}exists\_isTameFor\_costsSquares\_one}}}
\expandafter\newcommand\csname LN@isLiftWeightkappa\endcsname{\href{https://github.com/lsernau/grh-lean/blob/72fead102f663d5c7a407964a13dc0cf9d977207/grh-compact/src/New/Exchange/CrudeModularData.lean\#L28-L31}{{\leanfont Data.\allowbreak{}isLiftWeight\_kappa}}}
\expandafter\newcommand\csname LN@homX\endcsname{\href{https://github.com/lsernau/grh-lean/blob/72fead102f663d5c7a407964a13dc0cf9d977207/grh-compact/src/New/Exchange/CrudeModularData.lean\#L35-L35}{{\leanfont Data.\allowbreak{}homX}}}
\expandafter\newcommand\csname LN@liftKernela\endcsname{\href{https://github.com/lsernau/grh-lean/blob/72fead102f663d5c7a407964a13dc0cf9d977207/grh-compact/src/New/Exchange/CrudeModularData.lean\#L40-L44}{{\leanfont Data.\allowbreak{}liftKernel}}}
\expandafter\newcommand\csname LN@Translations\endcsname{\href{https://github.com/lsernau/grh-lean/blob/72fead102f663d5c7a407964a13dc0cf9d977207/grh-compact/src/New/Exchange/CrudeModularData.lean\#L48-L48}{{\leanfont Data.\allowbreak{}Translations}}}
\expandafter\newcommand\csname LN@translateDatum\endcsname{\href{https://github.com/lsernau/grh-lean/blob/72fead102f663d5c7a407964a13dc0cf9d977207/grh-compact/src/New/Exchange/CrudeModularData.lean\#L52-L53}{{\leanfont Data.\allowbreak{}translateDatum}}}
\expandafter\newcommand\csname LN@chainComb\endcsname{\href{https://github.com/lsernau/grh-lean/blob/72fead102f663d5c7a407964a13dc0cf9d977207/grh-compact/src/New/Exchange/CrudeModularData.lean\#L57-L58}{{\leanfont Data.\allowbreak{}chainComb}}}
\expandafter\newcommand\csname LN@B\endcsname{\href{https://github.com/lsernau/grh-lean/blob/72fead102f663d5c7a407964a13dc0cf9d977207/grh-compact/src/New/Exchange/CrudeModularData.lean\#L62-L62}{{\leanfont Data.\allowbreak{}Bμ}}}
\expandafter\newcommand\csname LN@spanpiGBar\endcsname{\href{https://github.com/lsernau/grh-lean/blob/72fead102f663d5c7a407964a13dc0cf9d977207/grh-compact/src/New/Exchange/CrudeModularData.lean\#L66-L70}{{\leanfont span\_piGBar}}}
\expandafter\newcommand\csname LN@inducedPlace\endcsname{\href{https://github.com/lsernau/grh-lean/blob/72fead102f663d5c7a407964a13dc0cf9d977207/grh-compact/src/New/Exchange/CrudeModularData.lean\#L73-L73}{{\leanfont Data.\allowbreak{}inducedPlace}}}
\expandafter\newcommand\csname LN@IsMarked\endcsname{\href{https://github.com/lsernau/grh-lean/blob/72fead102f663d5c7a407964a13dc0cf9d977207/grh-compact/src/New/Exchange/CrudeModularData.lean\#L77-L81}{{\leanfont Data.\allowbreak{}IsMarked}}}
\expandafter\newcommand\csname LN@IsFullMark\endcsname{\href{https://github.com/lsernau/grh-lean/blob/72fead102f663d5c7a407964a13dc0cf9d977207/grh-compact/src/New/Exchange/CrudeModularData.lean\#L85-L86}{{\leanfont Data.\allowbreak{}IsFullMark}}}
\expandafter\newcommand\csname LN@integralColumns\endcsname{\href{https://github.com/lsernau/grh-lean/blob/72fead102f663d5c7a407964a13dc0cf9d977207/grh-compact/src/New/Exchange/CrudeModularData.lean\#L90-L91}{{\leanfont Data.\allowbreak{}integralColumns}}}
\expandafter\newcommand\csname LN@evenPlaces\endcsname{\href{https://github.com/lsernau/grh-lean/blob/72fead102f663d5c7a407964a13dc0cf9d977207/grh-compact/src/New/Exchange/CrudeModularData.lean\#L95-L95}{{\leanfont Data.\allowbreak{}evenPlaces}}}
\expandafter\newcommand\csname LN@markExponent\endcsname{\href{https://github.com/lsernau/grh-lean/blob/72fead102f663d5c7a407964a13dc0cf9d977207/grh-compact/src/New/Exchange/CrudeModularData.lean\#L98-L98}{{\leanfont Data.\allowbreak{}markExponent}}}
\expandafter\newcommand\csname LN@lowColumns\endcsname{\href{https://github.com/lsernau/grh-lean/blob/72fead102f663d5c7a407964a13dc0cf9d977207/grh-compact/src/New/Exchange/CrudeModularData.lean\#L102-L102}{{\leanfont Data.\allowbreak{}lowColumns}}}
\expandafter\newcommand\csname LN@lowBoxes\endcsname{\href{https://github.com/lsernau/grh-lean/blob/72fead102f663d5c7a407964a13dc0cf9d977207/grh-compact/src/New/Exchange/CrudeModularData.lean\#L106-L106}{{\leanfont Data.\allowbreak{}lowBoxes}}}
\expandafter\newcommand\csname LN@cardresRingle\endcsname{\href{https://github.com/lsernau/grh-lean/blob/72fead102f663d5c7a407964a13dc0cf9d977207/grh-compact/src/New/Exchange/CrudeModularData.lean\#L111-L113}{{\leanfont Modulus.\allowbreak{}card\_resRing\_le}}}
\expandafter\newcommand\csname LN@smulvacuumDatummemstandardData\endcsname{\href{https://github.com/lsernau/grh-lean/blob/72fead102f663d5c7a407964a13dc0cf9d977207/grh-compact/src/New/Exchange/CrudeModularData.lean\#L121-L128}{{\leanfont Parameters.\allowbreak{}smul\_vacuumDatum\_mem\_standardData}}}
\expandafter\newcommand\csname LN@isAdmissibleCombvacuumDatum\endcsname{\href{https://github.com/lsernau/grh-lean/blob/72fead102f663d5c7a407964a13dc0cf9d977207/grh-compact/src/New/Exchange/CrudeModularData.lean\#L133-L139}{{\leanfont Parameters.\allowbreak{}isAdmissibleComb\_vacuumDatum}}}
\expandafter\newcommand\csname LN@existsisBaseAdapted\endcsname{\href{https://github.com/lsernau/grh-lean/blob/72fead102f663d5c7a407964a13dc0cf9d977207/grh-compact/src/New/Exchange/EqualRankValue.lean\#L24-L29}{{\leanfont exists\_isBaseAdapted}}}
\expandafter\newcommand\csname LN@slicezerocongr\endcsname{\href{https://github.com/lsernau/grh-lean/blob/72fead102f663d5c7a407964a13dc0cf9d977207/grh-compact/src/New/Exchange/EqualRankValue.lean\#L31-L33}{{\leanfont slice\_zero\_congr}}}
\expandafter\newcommand\csname LN@slicezerocomp\endcsname{\href{https://github.com/lsernau/grh-lean/blob/72fead102f663d5c7a407964a13dc0cf9d977207/grh-compact/src/New/Exchange/EqualRankValue.lean\#L35-L39}{{\leanfont slice\_zero\_comp}}}
\expandafter\newcommand\csname LN@existsstdExchangeremovingplanes\endcsname{\href{https://github.com/lsernau/grh-lean/blob/72fead102f663d5c7a407964a13dc0cf9d977207/grh-compact/src/New/Exchange/EqualRankValue.lean\#L44-L87}{{\leanfont exists\_stdExchange\_removing\_planes}}}
\expandafter\newcommand\csname LN@exchangeofpartners\endcsname{\href{https://github.com/lsernau/grh-lean/blob/72fead102f663d5c7a407964a13dc0cf9d977207/grh-compact/src/New/Exchange/EqualRankValue.lean\#L97-L141}{{\leanfont exchange\_of\_partners}}}
\expandafter\newcommand\csname LN@namevacElemofisEmpty\endcsname{\href{https://github.com/lsernau/grh-lean/blob/72fead102f663d5c7a407964a13dc0cf9d977207/grh-compact/src/New/Exchange/EqualRankValue.lean\#L146-L149}{{\leanfont name\_vacElem\_of\_isEmpty}}}
\expandafter\newcommand\csname LN@existsequalRankThetaRatiovalue\endcsname{\href{https://github.com/lsernau/grh-lean/blob/72fead102f663d5c7a407964a13dc0cf9d977207/grh-compact/src/New/Exchange/EqualRankValue.lean\#L160-L209}{{\leanfont exists\_equalRankThetaRatio\_value}}}
\expandafter\newcommand\csname LN@exchangetheorem\endcsname{\href{https://github.com/lsernau/grh-lean/blob/72fead102f663d5c7a407964a13dc0cf9d977207/grh-compact/src/New/Exchange/ExchangeTheorem.lean\#L16-L46}{{\leanfont exchange\_theorem}}}
\expandafter\newcommand\csname LQ@exchangetheorem\endcsname{\leanquoteline{0}{lemma exchange\_theorem σ (hs : negCount s = negCount t + negCount t') \{e\} (he : card a =}\leanquoteline{2}{card b + 2 * e) (hn : 2 * negCount s + 2 ≤ card n) \{μ : \_ → ℤ\} (hμ : ∀ v,}\leanquoteline{2}{card n ≤ 2 * μ v + 2 * negCount s) \{v\} (hv : IsRapidDatum v) : ∃ C x, IsRapidDatum x ∧ ∀ 𝔪 ν χ,}\leanquoteline{2}{IsTameFor 𝔪 ν μ χ → let γ := rankTwist ν⁻¹ n; ∃ c ≠ 0, ‖c‖ \textasciicircum{} 2 = C * lValueProd 𝔪.places ν e}\leanquoteline{4}{∧ ∀ Φ Ψ, IsPartner ⊤ σ χ Φ Ψ → ∀ g, thetaIntegral W V γ 1 (modDatum W V 𝔪 Φ x) g =}\leanquoteline{4}{c * ∫ h, pairTheta V' W γ γ (modDatum \_ W 𝔪 Ψ v) (h, g) ∂V'.quotMeasure}}
\expandafter\newcommand\csname LN@integralintegralsumeq\endcsname{\href{https://github.com/lsernau/grh-lean/blob/72fead102f663d5c7a407964a13dc0cf9d977207/grh-compact/src/New/Exchange/OneSquare.lean\#L22-L27}{{\leanfont integral\_integral\_sum\_eq}}}
\expandafter\newcommand\csname LN@applyeqmulofformRepone\endcsname{\href{https://github.com/lsernau/grh-lean/blob/72fead102f663d5c7a407964a13dc0cf9d977207/grh-compact/src/New/Exchange/OneSquare.lean\#L29-L31}{{\leanfont apply\_eq\_mul\_of\_formRep\_one}}}
\expandafter\newcommand\csname LN@existsintegrablekernelbound\endcsname{\href{https://github.com/lsernau/grh-lean/blob/72fead102f663d5c7a407964a13dc0cf9d977207/grh-compact/src/New/Exchange/OneSquare.lean\#L33-L44}{{\leanfont exists\_integrable\_kernel\_bound}}}
\expandafter\newcommand\csname LN@onesquare\endcsname{\href{https://github.com/lsernau/grh-lean/blob/72fead102f663d5c7a407964a13dc0cf9d977207/grh-compact/src/New/Exchange/OneSquare.lean\#L59-L134}{{\leanfont one\_square}}}
\expandafter\newcommand\csname LN@isFormnegiffandintegralneg\endcsname{\href{https://github.com/lsernau/grh-lean/blob/72fead102f663d5c7a407964a13dc0cf9d977207/grh-compact/src/New/Exchange/SquareLawFourier.lean\#L29-L40}{{\leanfont isForm\_neg\_iff\_and\_integral\_neg}}}
\expandafter\newcommand\csname LN@thetaLargemove\endcsname{\href{https://github.com/lsernau/grh-lean/blob/72fead102f663d5c7a407964a13dc0cf9d977207/grh-compact/src/New/Exchange/SquareLawFourier.lean\#L45-L62}{{\leanfont thetaLarge\_move}}}
\expandafter\newcommand\csname LN@SquareLaw\endcsname{\href{https://github.com/lsernau/grh-lean/blob/72fead102f663d5c7a407964a13dc0cf9d977207/grh-compact/src/New/Exchange/SquareLawFourier.lean\#L69-L72}{{\leanfont SquareLaw}}}
\expandafter\newcommand\csname LQ@SquareLaw\endcsname{\leanquoteline{0}{structure SquareLaw σ c (L : Matrix a n R → ℂ) (L' : Matrix n b R → ℂ) where}\leanquoteline{2}{left\_mul (k : GL a R) x : L (k.1 * x) = χ k.1.det * L x}\leanquoteline{2}{right\_mul (k : GL b R) y : L' (y * k.1) = χ k.1.det * L' y}\leanquoteline{2}{partners Φ Ψ : IsPartner ⊤ σ χ Φ Ψ → ∑ y, Ψ y * L' y = c * ∑ x, Φ x * L x}}
\expandafter\newcommand\csname LN@inclusionDatum\endcsname{\href{https://github.com/lsernau/grh-lean/blob/72fead102f663d5c7a407964a13dc0cf9d977207/grh-compact/src/New/Exchange/SquareLawFourier.lean\#L76-L77}{{\leanfont inclusionDatum}}}
\expandafter\newcommand\csname LN@nextDatum\endcsname{\href{https://github.com/lsernau/grh-lean/blob/72fead102f663d5c7a407964a13dc0cf9d977207/grh-compact/src/New/Exchange/SquareLawFourier.lean\#L81-L82}{{\leanfont nextDatum}}}
\expandafter\newcommand\csname LN@swapFirstTwo\endcsname{\href{https://github.com/lsernau/grh-lean/blob/72fead102f663d5c7a407964a13dc0cf9d977207/grh-compact/src/New/Exchange/SquareLawFourier.lean\#L85-L86}{{\leanfont swapFirstTwo}}}
\expandafter\newcommand\csname LN@sumeqsumfromRows\endcsname{\href{https://github.com/lsernau/grh-lean/blob/72fead102f663d5c7a407964a13dc0cf9d977207/grh-compact/src/New/Exchange/SquareLawFourier.lean\#L89-L90}{{\leanfont sum\_eq\_sum\_fromRows}}}
\expandafter\newcommand\csname LN@sumeqsumfromCols\endcsname{\href{https://github.com/lsernau/grh-lean/blob/72fead102f663d5c7a407964a13dc0cf9d977207/grh-compact/src/New/Exchange/SquareLawFourier.lean\#L93-L94}{{\leanfont sum\_eq\_sum\_fromCols}}}
\expandafter\newcommand\csname LN@injectiveglinvmulmul\endcsname{\href{https://github.com/lsernau/grh-lean/blob/72fead102f663d5c7a407964a13dc0cf9d977207/grh-compact/src/New/Exchange/SquareLawFourier.lean\#L97-L98}{{\leanfont injective\_gl\_inv\_mul\_mul}}}
\expandafter\newcommand\csname LN@inclMatrixinl\endcsname{\href{https://github.com/lsernau/grh-lean/blob/72fead102f663d5c7a407964a13dc0cf9d977207/grh-compact/src/New/Exchange/SquareLawFourier.lean\#L106-L106}{{\leanfont inclMatrix\_inl}}}
\expandafter\newcommand\csname LN@sumeqnextDatum\endcsname{\href{https://github.com/lsernau/grh-lean/blob/72fead102f663d5c7a407964a13dc0cf9d977207/grh-compact/src/New/Exchange/SquareLawFourier.lean\#L111-L127}{{\leanfont SquareLaw.\allowbreak{}sum\_eq\_nextDatum}}}
\expandafter\newcommand\csname LN@blockInlGL\endcsname{\href{https://github.com/lsernau/grh-lean/blob/72fead102f663d5c7a407964a13dc0cf9d977207/grh-compact/src/New/Exchange/SquareLawFourier.lean\#L130-L130}{{\leanfont blockInlGL}}}
\expandafter\newcommand\csname LN@graphKernel\endcsname{\href{https://github.com/lsernau/grh-lean/blob/72fead102f663d5c7a407964a13dc0cf9d977207/grh-compact/src/New/Exchange/SquareLawFourier.lean\#L134-L135}{{\leanfont graphKernel}}}
\expandafter\newcommand\csname LN@graphKernelinvLastRowsfirstCols\endcsname{\href{https://github.com/lsernau/grh-lean/blob/72fead102f663d5c7a407964a13dc0cf9d977207/grh-compact/src/New/Exchange/SquareLawFourier.lean\#L137-L139}{{\leanfont graphKernel\_invLastRows\_firstCols}}}
\expandafter\newcommand\csname LN@existseqofgraphKernelnezero\endcsname{\href{https://github.com/lsernau/grh-lean/blob/72fead102f663d5c7a407964a13dc0cf9d977207/grh-compact/src/New/Exchange/SquareLawFourier.lean\#L142-L143}{{\leanfont exists\_eq\_of\_graphKernel\_ne\_zero}}}
\expandafter\newcommand\csname LN@graphKernelequivariant\endcsname{\href{https://github.com/lsernau/grh-lean/blob/72fead102f663d5c7a407964a13dc0cf9d977207/grh-compact/src/New/Exchange/SquareLawFourier.lean\#L147-L153}{{\leanfont graphKernel\_equivariant}}}
\expandafter\newcommand\csname LN@graphKernelmulright\endcsname{\href{https://github.com/lsernau/grh-lean/blob/72fead102f663d5c7a407964a13dc0cf9d977207/grh-compact/src/New/Exchange/SquareLawFourier.lean\#L155-L156}{{\leanfont graphKernel\_mul\_right}}}
\expandafter\newcommand\csname LN@sumgraphKernelmul\endcsname{\href{https://github.com/lsernau/grh-lean/blob/72fead102f663d5c7a407964a13dc0cf9d977207/grh-compact/src/New/Exchange/SquareLawFourier.lean\#L160-L166}{{\leanfont sum\_graphKernel\_mul}}}
\expandafter\newcommand\csname LN@existspartnerkernel\endcsname{\href{https://github.com/lsernau/grh-lean/blob/72fead102f663d5c7a407964a13dc0cf9d977207/grh-compact/src/New/Exchange/SquareLawFourier.lean\#L171-L193}{{\leanfont exists\_partner\_kernel}}}
\expandafter\newcommand\csname LN@sumfouriereqnextDatum\endcsname{\href{https://github.com/lsernau/grh-lean/blob/72fead102f663d5c7a407964a13dc0cf9d977207/grh-compact/src/New/Exchange/SquareLawFourier.lean\#L202-L220}{{\leanfont SquareLaw.\allowbreak{}sum\_fourier\_eq\_nextDatum}}}
\expandafter\newcommand\csname LN@formNegEquiv\endcsname{\href{https://github.com/lsernau/grh-lean/blob/72fead102f663d5c7a407964a13dc0cf9d977207/grh-compact/src/New/Exchange/TriangleArrows.lean\#L27-L30}{{\leanfont formNegEquiv}}}
\expandafter\newcommand\csname LN@existsadjointfirstarrow\endcsname{\href{https://github.com/lsernau/grh-lean/blob/72fead102f663d5c7a407964a13dc0cf9d977207/grh-compact/src/New/Exchange/TriangleArrows.lean\#L35-L44}{{\leanfont IsEdgeFamily.\allowbreak{}exists\_adjoint\_first\_arrow}}}
\expandafter\newcommand\csname LN@TriForm\endcsname{\href{https://github.com/lsernau/grh-lean/blob/72fead102f663d5c7a407964a13dc0cf9d977207/grh-compact/src/New/Exchange/TriangleArrows.lean\#L52-L52}{{\leanfont Data.\allowbreak{}TriForm}}}
\expandafter\newcommand\csname LN@triLift\endcsname{\href{https://github.com/lsernau/grh-lean/blob/72fead102f663d5c7a407964a13dc0cf9d977207/grh-compact/src/New/Exchange/TriangleArrows.lean\#L57-L59}{{\leanfont Data.\allowbreak{}triLift}}}
\expandafter\newcommand\csname LN@triNegEquiv\endcsname{\href{https://github.com/lsernau/grh-lean/blob/72fead102f663d5c7a407964a13dc0cf9d977207/grh-compact/src/New/Exchange/TriangleArrows.lean\#L63-L64}{{\leanfont Data.\allowbreak{}triNegEquiv}}}
\expandafter\newcommand\csname LN@triBase\endcsname{\href{https://github.com/lsernau/grh-lean/blob/72fead102f663d5c7a407964a13dc0cf9d977207/grh-compact/src/New/Exchange/TriangleArrows.lean\#L68-L69}{{\leanfont Data.\allowbreak{}triBase}}}
\expandafter\newcommand\csname LN@triNegEquivformRep\endcsname{\href{https://github.com/lsernau/grh-lean/blob/72fead102f663d5c7a407964a13dc0cf9d977207/grh-compact/src/New/Exchange/TriangleArrows.lean\#L71-L72}{{\leanfont Data.\allowbreak{}triNegEquiv\_formRep}}}
\expandafter\newcommand\csname LN@hasSuminnertriLift\endcsname{\href{https://github.com/lsernau/grh-lean/blob/72fead102f663d5c7a407964a13dc0cf9d977207/grh-compact/src/New/Exchange/TriangleArrows.lean\#L77-L106}{{\leanfont Data.\allowbreak{}hasSum\_inner\_triLift}}}
\expandafter\newcommand\csname LN@existsisTriangle\endcsname{\href{https://github.com/lsernau/grh-lean/blob/72fead102f663d5c7a407964a13dc0cf9d977207/grh-compact/src/New/Exchange/TriangleArrows.lean\#L113-L188}{{\leanfont Data.\allowbreak{}exists\_isTriangle}}}
\expandafter\newcommand\csname LQ@existsisTriangle\endcsname{\leanquoteline{0}{theorem exists\_isTriangle : ∃ d, IsTriangle D.m (D.triangleOrder z) D.CostsSquares (D.triLift z) d}\leanquoteline{4}{(fun i ↦ (D.triNegEquiv z i).toLinearIsometry.comp (D.triBase z).subtypeₗᵢ) (z = false)}}
\expandafter\newcommand\csname LN@IsIntegralMeasure\endcsname{\href{https://github.com/lsernau/grh-lean/blob/72fead102f663d5c7a407964a13dc0cf9d977207/grh-compact/src/New/Exchange/TriangleLevelExponent.lean\#L25-L27}{{\leanfont Data.\allowbreak{}TameChar.\allowbreak{}IsIntegralMeasure}}}
\expandafter\newcommand\csname LN@normsq\endcsname{\href{https://github.com/lsernau/grh-lean/blob/72fead102f663d5c7a407964a13dc0cf9d977207/grh-compact/src/New/Exchange/TriangleLevelExponent.lean\#L31-L34}{{\leanfont Data.\allowbreak{}TameChar.\allowbreak{}IsIntegralMeasure.\allowbreak{}norm\_sq}}}
\expandafter\newcommand\csname LN@normapplyonele\endcsname{\href{https://github.com/lsernau/grh-lean/blob/72fead102f663d5c7a407964a13dc0cf9d977207/grh-compact/src/New/Exchange/TriangleLevelExponent.lean\#L40-L44}{{\leanfont Data.\allowbreak{}TameChar.\allowbreak{}IsIntegralMeasure.\allowbreak{}norm\_apply\_one\_le}}}
\expandafter\newcommand\csname LN@pointDatum\endcsname{\href{https://github.com/lsernau/grh-lean/blob/72fead102f663d5c7a407964a13dc0cf9d977207/grh-compact/src/New/Exchange/TriangleLevelExponent.lean\#L55-L56}{{\leanfont Data.\allowbreak{}Pair.\allowbreak{}pointDatum}}}
\expandafter\newcommand\csname LN@pointKernel\endcsname{\href{https://github.com/lsernau/grh-lean/blob/72fead102f663d5c7a407964a13dc0cf9d977207/grh-compact/src/New/Exchange/TriangleLevelExponent.lean\#L61-L62}{{\leanfont Data.\allowbreak{}Pair.\allowbreak{}pointKernel}}}
\expandafter\newcommand\csname LN@trivialLiftapply\endcsname{\href{https://github.com/lsernau/grh-lean/blob/72fead102f663d5c7a407964a13dc0cf9d977207/grh-compact/src/New/Exchange/TriangleLevelExponent.lean\#L67-L74}{{\leanfont Data.\allowbreak{}Pair.\allowbreak{}trivialLift\_apply}}}
\expandafter\newcommand\csname LN@weylGaugepointKernelleone\endcsname{\href{https://github.com/lsernau/grh-lean/blob/72fead102f663d5c7a407964a13dc0cf9d977207/grh-compact/src/New/Exchange/TriangleLevelExponent.lean\#L80-L86}{{\leanfont Data.\allowbreak{}Pair.\allowbreak{}weylGauge\_pointKernel\_le\_one}}}
\expandafter\newcommand\csname LN@isIntegralMeasureintAlgPart\endcsname{\href{https://github.com/lsernau/grh-lean/blob/72fead102f663d5c7a407964a13dc0cf9d977207/grh-compact/src/New/Exchange/TriangleLevelExponent.lean\#L94-L107}{{\leanfont Data.\allowbreak{}Pair.\allowbreak{}isIntegralMeasure\_intAlgPart}}}
\expandafter\newcommand\csname LN@intAlgParteqmulL\endcsname{\href{https://github.com/lsernau/grh-lean/blob/72fead102f663d5c7a407964a13dc0cf9d977207/grh-compact/src/New/Exchange/TriangleLevelExponent.lean\#L115-L121}{{\leanfont Data.\allowbreak{}Pair.\allowbreak{}intAlgPart\_eq\_mul\_L}}}
\expandafter\newcommand\csname LN@rationezeroandreal\endcsname{\href{https://github.com/lsernau/grh-lean/blob/72fead102f663d5c7a407964a13dc0cf9d977207/grh-compact/src/New/Exchange/TriangleLevelExponent.lean\#L125-L127}{{\leanfont Data.\allowbreak{}Pair.\allowbreak{}ratio\_ne\_zero\_and\_real}}}
\expandafter\newcommand\csname LN@ratioeqintAlgPartdiv\endcsname{\href{https://github.com/lsernau/grh-lean/blob/72fead102f663d5c7a407964a13dc0cf9d977207/grh-compact/src/New/Exchange/TriangleLevelExponent.lean\#L131-L133}{{\leanfont Data.\allowbreak{}Pair.\allowbreak{}ratio\_eq\_intAlgPart\_div}}}
\expandafter\newcommand\csname LN@intAlgPartnezero\endcsname{\href{https://github.com/lsernau/grh-lean/blob/72fead102f663d5c7a407964a13dc0cf9d977207/grh-compact/src/New/Exchange/TriangleLevelExponent.lean\#L137-L138}{{\leanfont Data.\allowbreak{}Pair.\allowbreak{}intAlgPart\_ne\_zero}}}
\expandafter\newcommand\csname LN@trivialLiftnezero\endcsname{\href{https://github.com/lsernau/grh-lean/blob/72fead102f663d5c7a407964a13dc0cf9d977207/grh-compact/src/New/Exchange/TriangleLevelExponent.lean\#L142-L142}{{\leanfont Data.\allowbreak{}Pair.\allowbreak{}trivialLift\_ne\_zero}}}
\expandafter\newcommand\csname LN@heightLeofeqzpowmul\endcsname{\href{https://github.com/lsernau/grh-lean/blob/72fead102f663d5c7a407964a13dc0cf9d977207/grh-compact/src/New/Exchange/TriangleLevelExponent.lean\#L151-L154}{{\leanfont Data.\allowbreak{}IsFitFamily.\allowbreak{}heightLe\_of\_eq\_zpow\_mul}}}
\expandafter\newcommand\csname LN@lowerPoint\endcsname{\href{https://github.com/lsernau/grh-lean/blob/72fead102f663d5c7a407964a13dc0cf9d977207/grh-compact/src/New/Exchange/TriangleLevelExponent.lean\#L160-L161}{{\leanfont Data.\allowbreak{}Pair.\allowbreak{}lowerPoint}}}
\expandafter\newcommand\csname LN@lValueProdeq\endcsname{\href{https://github.com/lsernau/grh-lean/blob/72fead102f663d5c7a407964a13dc0cf9d977207/grh-compact/src/New/Exchange/TriangleLevelExponent.lean\#L166-L173}{{\leanfont Data.\allowbreak{}Pair.\allowbreak{}lValueProd\_eq}}}
\expandafter\newcommand\csname LN@IsTriangle\endcsname{\href{https://github.com/lsernau/grh-lean/blob/72fead102f663d5c7a407964a13dc0cf9d977207/grh-compact/src/New/Exchange/TriangleNorms.lean\#L20-L27}{{\leanfont IsTriangle}}}
\expandafter\newcommand\csname LN@weightedChain\endcsname{\href{https://github.com/lsernau/grh-lean/blob/72fead102f663d5c7a407964a13dc0cf9d977207/grh-compact/src/New/Exchange/TriangleNorms.lean\#L31-L31}{{\leanfont weightedChain}}}
\expandafter\newcommand\csname LN@weightedChaindescent\endcsname{\href{https://github.com/lsernau/grh-lean/blob/72fead102f663d5c7a407964a13dc0cf9d977207/grh-compact/src/New/Exchange/TriangleNorms.lean\#L36-L43}{{\leanfont IsTriangle.\allowbreak{}weightedChain\_descent}}}
\expandafter\newcommand\csname LN@eqzeroofweightedChain\endcsname{\href{https://github.com/lsernau/grh-lean/blob/72fead102f663d5c7a407964a13dc0cf9d977207/grh-compact/src/New/Exchange/TriangleNorms.lean\#L47-L49}{{\leanfont IsTriangle.\allowbreak{}eq\_zero\_of\_weightedChain}}}
\expandafter\newcommand\csname LN@innerweightedChaineq\endcsname{\href{https://github.com/lsernau/grh-lean/blob/72fead102f663d5c7a407964a13dc0cf9d977207/grh-compact/src/New/Exchange/TriangleNorms.lean\#L54-L66}{{\leanfont IsTriangle.\allowbreak{}inner\_weightedChain\_eq}}}
\newcommand{\ConcordanceRows}{\crow{maptrivialLift}{\href{https://github.com/lsernau/grh-lean/blob/72fead102f663d5c7a407964a13dc0cf9d977207/grh-compact/src/New/Exchange/AlgebraicPartConjugates.lean\#L13-L47}{{\leanfont LValueData.\allowbreak{}map\_trivialLift}}}{AlgebraicPartConjugates.lean}{13--47}
\crow{lognormmapnormsqtrivialLift}{\href{https://github.com/lsernau/grh-lean/blob/72fead102f663d5c7a407964a13dc0cf9d977207/grh-compact/src/New/Exchange/AlgebraicPartConjugates.lean\#L52-L81}{{\leanfont Data.\allowbreak{}Pair.\allowbreak{}log\_norm\_map\_norm\_sq\_trivialLift}}}{AlgebraicPartConjugates.lean}{52--81}
\crow{cardYle}{\href{https://github.com/lsernau/grh-lean/blob/72fead102f663d5c7a407964a13dc0cf9d977207/grh-compact/src/New/Exchange/AlgebraicPartDefs.lean\#L20-L21}{{\leanfont Data.\allowbreak{}card\_Y\_le}}}{AlgebraicPartDefs.lean}{20--21}
\crow{isHCParamY}{\href{https://github.com/lsernau/grh-lean/blob/72fead102f663d5c7a407964a13dc0cf9d977207/grh-compact/src/New/Exchange/AlgebraicPartDefs.lean\#L25-L38}{{\leanfont Data.\allowbreak{}isHCParam\_Y}}}{AlgebraicPartDefs.lean}{25--38}
\crow{cardV}{\href{https://github.com/lsernau/grh-lean/blob/72fead102f663d5c7a407964a13dc0cf9d977207/grh-compact/src/New/Exchange/AlgebraicPartDefs.lean\#L40-L40}{{\leanfont Data.\allowbreak{}card\_V}}}{AlgebraicPartDefs.lean}{40--40}
\crow{cardfilterbasisNumber}{\href{https://github.com/lsernau/grh-lean/blob/72fead102f663d5c7a407964a13dc0cf9d977207/grh-compact/src/New/Exchange/AlgebraicPartDefs.lean\#L42-L44}{{\leanfont Data.\allowbreak{}card\_filter\_basisNumber}}}{AlgebraicPartDefs.lean}{42--44}
\crow{varsigmastrictAnti}{\href{https://github.com/lsernau/grh-lean/blob/72fead102f663d5c7a407964a13dc0cf9d977207/grh-compact/src/New/Exchange/AlgebraicPartDefs.lean\#L46-L47}{{\leanfont Data.\allowbreak{}varsigma\_strictAnti}}}{AlgebraicPartDefs.lean}{46--47}
\crow{paramWltiff}{\href{https://github.com/lsernau/grh-lean/blob/72fead102f663d5c7a407964a13dc0cf9d977207/grh-compact/src/New/Exchange/AlgebraicPartDefs.lean\#L49-L50}{{\leanfont Data.\allowbreak{}paramW\_lt\_iff}}}{AlgebraicPartDefs.lean}{49--50}
\crow{cardnegsignsW}{\href{https://github.com/lsernau/grh-lean/blob/72fead102f663d5c7a407964a13dc0cf9d977207/grh-compact/src/New/Exchange/AlgebraicPartDefs.lean\#L53-L54}{{\leanfont Data.\allowbreak{}card\_neg\_signsW}}}{AlgebraicPartDefs.lean}{53--54}
\crow{fitsParamsauxPair}{\href{https://github.com/lsernau/grh-lean/blob/72fead102f663d5c7a407964a13dc0cf9d977207/grh-compact/src/New/Exchange/AlgebraicPartDefs.lean\#L58-L71}{{\leanfont Data.\allowbreak{}fitsParams\_auxPair}}}{AlgebraicPartDefs.lean}{58--71}
\crow{innereqpetersson}{\href{https://github.com/lsernau/grh-lean/blob/72fead102f663d5c7a407964a13dc0cf9d977207/grh-compact/src/New/Exchange/AlgebraicPartDefs.lean\#L76-L82}{{\leanfont LevelPlaces.\allowbreak{}LevelForm.\allowbreak{}inner\_eq\_petersson}}}{AlgebraicPartDefs.lean}{76--82}
\crow{chosenIndexEmb}{\href{https://github.com/lsernau/grh-lean/blob/72fead102f663d5c7a407964a13dc0cf9d977207/grh-compact/src/New/Exchange/AlgebraicPartDefs.lean\#L86-L86}{{\leanfont chosenIndexEmb}}}{AlgebraicPartDefs.lean}{86--86}
\crow{LValueData}{\href{https://github.com/lsernau/grh-lean/blob/72fead102f663d5c7a407964a13dc0cf9d977207/grh-compact/src/New/Exchange/AlgebraicPartDefs.lean\#L91-L98}{{\leanfont LValueData}}}{AlgebraicPartDefs.lean}{91--98}
\crow{auxPair}{\href{https://github.com/lsernau/grh-lean/blob/72fead102f663d5c7a407964a13dc0cf9d977207/grh-compact/src/New/Exchange/AlgebraicPartDefs.lean\#L101-L103}{{\leanfont LValueData.\allowbreak{}auxPair}}}{AlgebraicPartDefs.lean}{101--103}
\crow{kernelValue}{\href{https://github.com/lsernau/grh-lean/blob/72fead102f663d5c7a407964a13dc0cf9d977207/grh-compact/src/New/Exchange/AlgebraicPartDefs.lean\#L110-L111}{{\leanfont LValueData.\allowbreak{}kernelValue}}}{AlgebraicPartDefs.lean}{110--111}
\crow{isLiftWeight}{\href{https://github.com/lsernau/grh-lean/blob/72fead102f663d5c7a407964a13dc0cf9d977207/grh-compact/src/New/Exchange/AlgebraicPartDefs.lean\#L117-L118}{{\leanfont LValueData.\allowbreak{}isLiftWeight}}}{AlgebraicPartDefs.lean}{117--118}
\crow{ahasWeight}{\href{https://github.com/lsernau/grh-lean/blob/72fead102f663d5c7a407964a13dc0cf9d977207/grh-compact/src/New/Exchange/AlgebraicPartDefs.lean\#L120-L121}{{\leanfont LValueData.\allowbreak{}a\_hasWeight}}}{AlgebraicPartDefs.lean}{120--121}
\crow{bhasWeight}{\href{https://github.com/lsernau/grh-lean/blob/72fead102f663d5c7a407964a13dc0cf9d977207/grh-compact/src/New/Exchange/AlgebraicPartDefs.lean\#L123-L124}{{\leanfont LValueData.\allowbreak{}b\_hasWeight}}}{AlgebraicPartDefs.lean}{123--124}
\crow{liftKernel}{\href{https://github.com/lsernau/grh-lean/blob/72fead102f663d5c7a407964a13dc0cf9d977207/grh-compact/src/New/Exchange/AlgebraicPartDefs.lean\#L128-L129}{{\leanfont LValueData.\allowbreak{}liftKernel}}}{AlgebraicPartDefs.lean}{128--129}
\crow{conjWeylRepeqone}{\href{https://github.com/lsernau/grh-lean/blob/72fead102f663d5c7a407964a13dc0cf9d977207/grh-compact/src/New/Exchange/AlgebraicPartDefs.lean\#L132-L133}{{\leanfont LValueData.\allowbreak{}conjWeylRep\_eq\_one}}}{AlgebraicPartDefs.lean}{132--133}
\crow{kernelForm}{\href{https://github.com/lsernau/grh-lean/blob/72fead102f663d5c7a407964a13dc0cf9d977207/grh-compact/src/New/Exchange/AlgebraicPartDefs.lean\#L137-L142}{{\leanfont LValueData.\allowbreak{}kernelForm}}}{AlgebraicPartDefs.lean}{137--142}
\crow{chosenLevel}{\href{https://github.com/lsernau/grh-lean/blob/72fead102f663d5c7a407964a13dc0cf9d977207/grh-compact/src/New/Exchange/AlgebraicPartDefs.lean\#L146-L146}{{\leanfont LValueData.\allowbreak{}chosenLevel}}}{AlgebraicPartDefs.lean}{146--146}
\crow{trivialLift}{\href{https://github.com/lsernau/grh-lean/blob/72fead102f663d5c7a407964a13dc0cf9d977207/grh-compact/src/New/Exchange/AlgebraicPartDefs.lean\#L150-L150}{{\leanfont LValueData.\allowbreak{}trivialLift}}}{AlgebraicPartDefs.lean}{150--150}
\crow{kernelForminvariant}{\href{https://github.com/lsernau/grh-lean/blob/72fead102f663d5c7a407964a13dc0cf9d977207/grh-compact/src/New/Exchange/AlgebraicPartDefs.lean\#L155-L159}{{\leanfont LValueData.\allowbreak{}kernelForm\_invariant}}}{AlgebraicPartDefs.lean}{155--159}
\crow{trivialLifteqstdClassSumtranslate}{\href{https://github.com/lsernau/grh-lean/blob/72fead102f663d5c7a407964a13dc0cf9d977207/grh-compact/src/New/Exchange/AlgebraicPartDefs.lean\#L161-L164}{{\leanfont LValueData.\allowbreak{}trivialLift\_eq\_stdClassSum\_translate}}}{AlgebraicPartDefs.lean}{161--164}
\crow{toPolyeqrenameharmonicPoly}{\href{https://github.com/lsernau/grh-lean/blob/72fead102f663d5c7a407964a13dc0cf9d977207/grh-compact/src/New/Exchange/AlgebraicPartDefs.lean\#L170-L175}{{\leanfont LValueData.\allowbreak{}toPoly\_eq\_rename\_harmonicPoly}}}{AlgebraicPartDefs.lean}{170--175}
\crow{harmonicPolyeqschurCoeffPoly}{\href{https://github.com/lsernau/grh-lean/blob/72fead102f663d5c7a407964a13dc0cf9d977207/grh-compact/src/New/Exchange/AlgebraicPartDefs.lean\#L178-L183}{{\leanfont LValueData.\allowbreak{}harmonicPoly\_eq\_schurCoeffPoly}}}{AlgebraicPartDefs.lean}{178--183}
\crow{thetapairWeilfockDatum}{\href{https://github.com/lsernau/grh-lean/blob/72fead102f663d5c7a407964a13dc0cf9d977207/grh-compact/src/New/Exchange/AlgebraicPartDefs.lean\#L185-L190}{{\leanfont LValueData.\allowbreak{}theta\_pairWeil\_fockDatum}}}{AlgebraicPartDefs.lean}{185--190}
\crow{kernelFormofjetDegeqzero}{\href{https://github.com/lsernau/grh-lean/blob/72fead102f663d5c7a407964a13dc0cf9d977207/grh-compact/src/New/Exchange/AlgebraicPartDefs.lean\#L195-L202}{{\leanfont LValueData.\allowbreak{}kernelForm\_of\_jetDeg\_eq\_zero}}}{AlgebraicPartDefs.lean}{195--202}
\crow{thetaeqinnerkernelForm}{\href{https://github.com/lsernau/grh-lean/blob/72fead102f663d5c7a407964a13dc0cf9d977207/grh-compact/src/New/Exchange/AlgebraicPartDefs.lean\#L204-L206}{{\leanfont LValueData.\allowbreak{}theta\_eq\_inner\_kernelForm}}}{AlgebraicPartDefs.lean}{204--206}
\crow{innertrivialLift}{\href{https://github.com/lsernau/grh-lean/blob/72fead102f663d5c7a407964a13dc0cf9d977207/grh-compact/src/New/Exchange/AlgebraicPartDefs.lean\#L211-L213}{{\leanfont LValueData.\allowbreak{}inner\_trivialLift}}}{AlgebraicPartDefs.lean}{211--213}
\crow{hookNumbereqprodfactorial}{\href{https://github.com/lsernau/grh-lean/blob/72fead102f663d5c7a407964a13dc0cf9d977207/grh-compact/src/New/Exchange/AlgebraicPartDefs.lean\#L218-L222}{{\leanfont LValueData.\allowbreak{}hookNumber\_eq\_prod\_factorial}}}{AlgebraicPartDefs.lean}{218--222}
\crow{HasLiftRank}{\href{https://github.com/lsernau/grh-lean/blob/72fead102f663d5c7a407964a13dc0cf9d977207/grh-compact/src/New/Exchange/AlgebraicPartDefs.lean\#L227-L227}{{\leanfont Data.\allowbreak{}Pair.\allowbreak{}HasLiftRank}}}{AlgebraicPartDefs.lean}{227--227}
\crow{IsFlat}{\href{https://github.com/lsernau/grh-lean/blob/72fead102f663d5c7a407964a13dc0cf9d977207/grh-compact/src/New/Exchange/AlgebraicPartDefs.lean\#L231-L231}{{\leanfont Data.\allowbreak{}Pair.\allowbreak{}IsFlat}}}{AlgebraicPartDefs.lean}{231--231}
\crow{denomBound}{\href{https://github.com/lsernau/grh-lean/blob/72fead102f663d5c7a407964a13dc0cf9d977207/grh-compact/src/New/Exchange/AlgebraicPartDefs.lean\#L235-L236}{{\leanfont Data.\allowbreak{}Pair.\allowbreak{}denomBound}}}{AlgebraicPartDefs.lean}{235--236}
\crow{lValueData}{\href{https://github.com/lsernau/grh-lean/blob/72fead102f663d5c7a407964a13dc0cf9d977207/grh-compact/src/New/Exchange/AlgebraicPartDefs.lean\#L242-L244}{{\leanfont Data.\allowbreak{}Pair.\allowbreak{}lValueData}}}{AlgebraicPartDefs.lean}{242--244}
\crow{trivialLifta}{\href{https://github.com/lsernau/grh-lean/blob/72fead102f663d5c7a407964a13dc0cf9d977207/grh-compact/src/New/Exchange/AlgebraicPartDefs.lean\#L249-L250}{{\leanfont Data.\allowbreak{}Pair.\allowbreak{}trivialLift}}}{AlgebraicPartDefs.lean}{249--250}
\crow{intAlgPart}{\href{https://github.com/lsernau/grh-lean/blob/72fead102f663d5c7a407964a13dc0cf9d977207/grh-compact/src/New/Exchange/AlgebraicPartDefs.lean\#L254-L255}{{\leanfont Data.\allowbreak{}Pair.\allowbreak{}intAlgPart}}}{AlgebraicPartDefs.lean}{254--255}
\crow{stdLevelIndexlastResLeveleq}{\href{https://github.com/lsernau/grh-lean/blob/72fead102f663d5c7a407964a13dc0cf9d977207/grh-compact/src/New/Exchange/ChainNorms.lean\#L19-L37}{{\leanfont Data.\allowbreak{}stdLevelIndex\_lastResLevel\_eq}}}{ChainNorms.lean}{19--37}
\crow{inneriotaAtLevel}{\href{https://github.com/lsernau/grh-lean/blob/72fead102f663d5c7a407964a13dc0cf9d977207/grh-compact/src/New/Exchange/ChainNorms.lean\#L42-L47}{{\leanfont Data.\allowbreak{}inner\_iotaAtLevel}}}{ChainNorms.lean}{42--47}
\crow{iotaAtLevelinjective}{\href{https://github.com/lsernau/grh-lean/blob/72fead102f663d5c7a407964a13dc0cf9d977207/grh-compact/src/New/Exchange/ChainNorms.lean\#L53-L55}{{\leanfont Data.\allowbreak{}iotaAtLevel\_injective}}}{ChainNorms.lean}{53--55}
\crow{CostsSquares}{\href{https://github.com/lsernau/grh-lean/blob/72fead102f663d5c7a407964a13dc0cf9d977207/grh-compact/src/New/Exchange/CostOfSquares.lean\#L17-L18}{{\leanfont Data.\allowbreak{}CostsSquares}}}{CostOfSquares.lean}{17--18}
\crow{mul}{\href{https://github.com/lsernau/grh-lean/blob/72fead102f663d5c7a407964a13dc0cf9d977207/grh-compact/src/New/Exchange/CostOfSquares.lean\#L22-L24}{{\leanfont Data.\allowbreak{}CostsSquares.\allowbreak{}mul}}}{CostOfSquares.lean}{22--24}
\crow{card}{\href{https://github.com/lsernau/grh-lean/blob/72fead102f663d5c7a407964a13dc0cf9d977207/grh-compact/src/New/Exchange/CostOfSquares.lean\#L28-L29}{{\leanfont Data.\allowbreak{}card\_𝔪}}}{CostOfSquares.lean}{28--29}
\crow{ofnormsqeqmulcardzpow}{\href{https://github.com/lsernau/grh-lean/blob/72fead102f663d5c7a407964a13dc0cf9d977207/grh-compact/src/New/Exchange/CostOfSquares.lean\#L33-L36}{{\leanfont Data.\allowbreak{}CostsSquares.\allowbreak{}of\_norm\_sq\_eq\_mul\_card\_zpow}}}{CostOfSquares.lean}{33--36}
\crow{one}{\href{https://github.com/lsernau/grh-lean/blob/72fead102f663d5c7a407964a13dc0cf9d977207/grh-compact/src/New/Exchange/CostOfSquares.lean\#L38-L38}{{\leanfont Data.\allowbreak{}CostsSquares.\allowbreak{}one}}}{CostOfSquares.lean}{38--38}
\crow{isPrimitiveMulCharofisTameFor}{\href{https://github.com/lsernau/grh-lean/blob/72fead102f663d5c7a407964a13dc0cf9d977207/grh-compact/src/New/Exchange/CostOfSquares.lean\#L42-L49}{{\leanfont Data.\allowbreak{}isPrimitiveMulChar\_of\_isTameFor}}}{CostOfSquares.lean}{42--49}
\crow{existsisTameForcostsSquaresone}{\href{https://github.com/lsernau/grh-lean/blob/72fead102f663d5c7a407964a13dc0cf9d977207/grh-compact/src/New/Exchange/CostOfSquares.lean\#L54-L63}{{\leanfont Data.\allowbreak{}exists\_isTameFor\_costsSquares\_one}}}{CostOfSquares.lean}{54--63}
\crow{isLiftWeightkappa}{\href{https://github.com/lsernau/grh-lean/blob/72fead102f663d5c7a407964a13dc0cf9d977207/grh-compact/src/New/Exchange/CrudeModularData.lean\#L28-L31}{{\leanfont Data.\allowbreak{}isLiftWeight\_kappa}}}{CrudeModularData.lean}{28--31}
\crow{homX}{\href{https://github.com/lsernau/grh-lean/blob/72fead102f663d5c7a407964a13dc0cf9d977207/grh-compact/src/New/Exchange/CrudeModularData.lean\#L35-L35}{{\leanfont Data.\allowbreak{}homX}}}{CrudeModularData.lean}{35--35}
\crow{liftKernela}{\href{https://github.com/lsernau/grh-lean/blob/72fead102f663d5c7a407964a13dc0cf9d977207/grh-compact/src/New/Exchange/CrudeModularData.lean\#L40-L44}{{\leanfont Data.\allowbreak{}liftKernel}}}{CrudeModularData.lean}{40--44}
\crow{Translations}{\href{https://github.com/lsernau/grh-lean/blob/72fead102f663d5c7a407964a13dc0cf9d977207/grh-compact/src/New/Exchange/CrudeModularData.lean\#L48-L48}{{\leanfont Data.\allowbreak{}Translations}}}{CrudeModularData.lean}{48--48}
\crow{translateDatum}{\href{https://github.com/lsernau/grh-lean/blob/72fead102f663d5c7a407964a13dc0cf9d977207/grh-compact/src/New/Exchange/CrudeModularData.lean\#L52-L53}{{\leanfont Data.\allowbreak{}translateDatum}}}{CrudeModularData.lean}{52--53}
\crow{chainComb}{\href{https://github.com/lsernau/grh-lean/blob/72fead102f663d5c7a407964a13dc0cf9d977207/grh-compact/src/New/Exchange/CrudeModularData.lean\#L57-L58}{{\leanfont Data.\allowbreak{}chainComb}}}{CrudeModularData.lean}{57--58}
\crow{B}{\href{https://github.com/lsernau/grh-lean/blob/72fead102f663d5c7a407964a13dc0cf9d977207/grh-compact/src/New/Exchange/CrudeModularData.lean\#L62-L62}{{\leanfont Data.\allowbreak{}Bμ}}}{CrudeModularData.lean}{62--62}
\crow{spanpiGBar}{\href{https://github.com/lsernau/grh-lean/blob/72fead102f663d5c7a407964a13dc0cf9d977207/grh-compact/src/New/Exchange/CrudeModularData.lean\#L66-L70}{{\leanfont span\_piGBar}}}{CrudeModularData.lean}{66--70}
\crow{inducedPlace}{\href{https://github.com/lsernau/grh-lean/blob/72fead102f663d5c7a407964a13dc0cf9d977207/grh-compact/src/New/Exchange/CrudeModularData.lean\#L73-L73}{{\leanfont Data.\allowbreak{}inducedPlace}}}{CrudeModularData.lean}{73--73}
\crow{IsMarked}{\href{https://github.com/lsernau/grh-lean/blob/72fead102f663d5c7a407964a13dc0cf9d977207/grh-compact/src/New/Exchange/CrudeModularData.lean\#L77-L81}{{\leanfont Data.\allowbreak{}IsMarked}}}{CrudeModularData.lean}{77--81}
\crow{IsFullMark}{\href{https://github.com/lsernau/grh-lean/blob/72fead102f663d5c7a407964a13dc0cf9d977207/grh-compact/src/New/Exchange/CrudeModularData.lean\#L85-L86}{{\leanfont Data.\allowbreak{}IsFullMark}}}{CrudeModularData.lean}{85--86}
\crow{integralColumns}{\href{https://github.com/lsernau/grh-lean/blob/72fead102f663d5c7a407964a13dc0cf9d977207/grh-compact/src/New/Exchange/CrudeModularData.lean\#L90-L91}{{\leanfont Data.\allowbreak{}integralColumns}}}{CrudeModularData.lean}{90--91}
\crow{evenPlaces}{\href{https://github.com/lsernau/grh-lean/blob/72fead102f663d5c7a407964a13dc0cf9d977207/grh-compact/src/New/Exchange/CrudeModularData.lean\#L95-L95}{{\leanfont Data.\allowbreak{}evenPlaces}}}{CrudeModularData.lean}{95--95}
\crow{markExponent}{\href{https://github.com/lsernau/grh-lean/blob/72fead102f663d5c7a407964a13dc0cf9d977207/grh-compact/src/New/Exchange/CrudeModularData.lean\#L98-L98}{{\leanfont Data.\allowbreak{}markExponent}}}{CrudeModularData.lean}{98--98}
\crow{lowColumns}{\href{https://github.com/lsernau/grh-lean/blob/72fead102f663d5c7a407964a13dc0cf9d977207/grh-compact/src/New/Exchange/CrudeModularData.lean\#L102-L102}{{\leanfont Data.\allowbreak{}lowColumns}}}{CrudeModularData.lean}{102--102}
\crow{lowBoxes}{\href{https://github.com/lsernau/grh-lean/blob/72fead102f663d5c7a407964a13dc0cf9d977207/grh-compact/src/New/Exchange/CrudeModularData.lean\#L106-L106}{{\leanfont Data.\allowbreak{}lowBoxes}}}{CrudeModularData.lean}{106--106}
\crow{cardresRingle}{\href{https://github.com/lsernau/grh-lean/blob/72fead102f663d5c7a407964a13dc0cf9d977207/grh-compact/src/New/Exchange/CrudeModularData.lean\#L111-L113}{{\leanfont Modulus.\allowbreak{}card\_resRing\_le}}}{CrudeModularData.lean}{111--113}
\crow{smulvacuumDatummemstandardData}{\href{https://github.com/lsernau/grh-lean/blob/72fead102f663d5c7a407964a13dc0cf9d977207/grh-compact/src/New/Exchange/CrudeModularData.lean\#L121-L128}{{\leanfont Parameters.\allowbreak{}smul\_vacuumDatum\_mem\_standardData}}}{CrudeModularData.lean}{121--128}
\crow{isAdmissibleCombvacuumDatum}{\href{https://github.com/lsernau/grh-lean/blob/72fead102f663d5c7a407964a13dc0cf9d977207/grh-compact/src/New/Exchange/CrudeModularData.lean\#L133-L139}{{\leanfont Parameters.\allowbreak{}isAdmissibleComb\_vacuumDatum}}}{CrudeModularData.lean}{133--139}
\crow{existsisBaseAdapted}{\href{https://github.com/lsernau/grh-lean/blob/72fead102f663d5c7a407964a13dc0cf9d977207/grh-compact/src/New/Exchange/EqualRankValue.lean\#L24-L29}{{\leanfont exists\_isBaseAdapted}}}{EqualRankValue.lean}{24--29}
\crow{slicezerocongr}{\href{https://github.com/lsernau/grh-lean/blob/72fead102f663d5c7a407964a13dc0cf9d977207/grh-compact/src/New/Exchange/EqualRankValue.lean\#L31-L33}{{\leanfont slice\_zero\_congr}}}{EqualRankValue.lean}{31--33}
\crow{slicezerocomp}{\href{https://github.com/lsernau/grh-lean/blob/72fead102f663d5c7a407964a13dc0cf9d977207/grh-compact/src/New/Exchange/EqualRankValue.lean\#L35-L39}{{\leanfont slice\_zero\_comp}}}{EqualRankValue.lean}{35--39}
\crow{existsstdExchangeremovingplanes}{\href{https://github.com/lsernau/grh-lean/blob/72fead102f663d5c7a407964a13dc0cf9d977207/grh-compact/src/New/Exchange/EqualRankValue.lean\#L44-L87}{{\leanfont exists\_stdExchange\_removing\_planes}}}{EqualRankValue.lean}{44--87}
\crow{exchangeofpartners}{\href{https://github.com/lsernau/grh-lean/blob/72fead102f663d5c7a407964a13dc0cf9d977207/grh-compact/src/New/Exchange/EqualRankValue.lean\#L97-L141}{{\leanfont exchange\_of\_partners}}}{EqualRankValue.lean}{97--141}
\crow{namevacElemofisEmpty}{\href{https://github.com/lsernau/grh-lean/blob/72fead102f663d5c7a407964a13dc0cf9d977207/grh-compact/src/New/Exchange/EqualRankValue.lean\#L146-L149}{{\leanfont name\_vacElem\_of\_isEmpty}}}{EqualRankValue.lean}{146--149}
\crow{existsequalRankThetaRatiovalue}{\href{https://github.com/lsernau/grh-lean/blob/72fead102f663d5c7a407964a13dc0cf9d977207/grh-compact/src/New/Exchange/EqualRankValue.lean\#L160-L209}{{\leanfont exists\_equalRankThetaRatio\_value}}}{EqualRankValue.lean}{160--209}
\crow{exchangetheorem}{\href{https://github.com/lsernau/grh-lean/blob/72fead102f663d5c7a407964a13dc0cf9d977207/grh-compact/src/New/Exchange/ExchangeTheorem.lean\#L16-L46}{{\leanfont exchange\_theorem}}}{ExchangeTheorem.lean}{16--46}
\crow{integralintegralsumeq}{\href{https://github.com/lsernau/grh-lean/blob/72fead102f663d5c7a407964a13dc0cf9d977207/grh-compact/src/New/Exchange/OneSquare.lean\#L22-L27}{{\leanfont integral\_integral\_sum\_eq}}}{OneSquare.lean}{22--27}
\crow{applyeqmulofformRepone}{\href{https://github.com/lsernau/grh-lean/blob/72fead102f663d5c7a407964a13dc0cf9d977207/grh-compact/src/New/Exchange/OneSquare.lean\#L29-L31}{{\leanfont apply\_eq\_mul\_of\_formRep\_one}}}{OneSquare.lean}{29--31}
\crow{existsintegrablekernelbound}{\href{https://github.com/lsernau/grh-lean/blob/72fead102f663d5c7a407964a13dc0cf9d977207/grh-compact/src/New/Exchange/OneSquare.lean\#L33-L44}{{\leanfont exists\_integrable\_kernel\_bound}}}{OneSquare.lean}{33--44}
\crow{onesquare}{\href{https://github.com/lsernau/grh-lean/blob/72fead102f663d5c7a407964a13dc0cf9d977207/grh-compact/src/New/Exchange/OneSquare.lean\#L59-L134}{{\leanfont one\_square}}}{OneSquare.lean}{59--134}
\crow{isFormnegiffandintegralneg}{\href{https://github.com/lsernau/grh-lean/blob/72fead102f663d5c7a407964a13dc0cf9d977207/grh-compact/src/New/Exchange/SquareLawFourier.lean\#L29-L40}{{\leanfont isForm\_neg\_iff\_and\_integral\_neg}}}{SquareLawFourier.lean}{29--40}
\crow{thetaLargemove}{\href{https://github.com/lsernau/grh-lean/blob/72fead102f663d5c7a407964a13dc0cf9d977207/grh-compact/src/New/Exchange/SquareLawFourier.lean\#L45-L62}{{\leanfont thetaLarge\_move}}}{SquareLawFourier.lean}{45--62}
\crow{SquareLaw}{\href{https://github.com/lsernau/grh-lean/blob/72fead102f663d5c7a407964a13dc0cf9d977207/grh-compact/src/New/Exchange/SquareLawFourier.lean\#L69-L72}{{\leanfont SquareLaw}}}{SquareLawFourier.lean}{69--72}
\crow{inclusionDatum}{\href{https://github.com/lsernau/grh-lean/blob/72fead102f663d5c7a407964a13dc0cf9d977207/grh-compact/src/New/Exchange/SquareLawFourier.lean\#L76-L77}{{\leanfont inclusionDatum}}}{SquareLawFourier.lean}{76--77}
\crow{nextDatum}{\href{https://github.com/lsernau/grh-lean/blob/72fead102f663d5c7a407964a13dc0cf9d977207/grh-compact/src/New/Exchange/SquareLawFourier.lean\#L81-L82}{{\leanfont nextDatum}}}{SquareLawFourier.lean}{81--82}
\crow{swapFirstTwo}{\href{https://github.com/lsernau/grh-lean/blob/72fead102f663d5c7a407964a13dc0cf9d977207/grh-compact/src/New/Exchange/SquareLawFourier.lean\#L85-L86}{{\leanfont swapFirstTwo}}}{SquareLawFourier.lean}{85--86}
\crow{sumeqsumfromRows}{\href{https://github.com/lsernau/grh-lean/blob/72fead102f663d5c7a407964a13dc0cf9d977207/grh-compact/src/New/Exchange/SquareLawFourier.lean\#L89-L90}{{\leanfont sum\_eq\_sum\_fromRows}}}{SquareLawFourier.lean}{89--90}
\crow{sumeqsumfromCols}{\href{https://github.com/lsernau/grh-lean/blob/72fead102f663d5c7a407964a13dc0cf9d977207/grh-compact/src/New/Exchange/SquareLawFourier.lean\#L93-L94}{{\leanfont sum\_eq\_sum\_fromCols}}}{SquareLawFourier.lean}{93--94}
\crow{injectiveglinvmulmul}{\href{https://github.com/lsernau/grh-lean/blob/72fead102f663d5c7a407964a13dc0cf9d977207/grh-compact/src/New/Exchange/SquareLawFourier.lean\#L97-L98}{{\leanfont injective\_gl\_inv\_mul\_mul}}}{SquareLawFourier.lean}{97--98}
\crow{inclMatrixinl}{\href{https://github.com/lsernau/grh-lean/blob/72fead102f663d5c7a407964a13dc0cf9d977207/grh-compact/src/New/Exchange/SquareLawFourier.lean\#L106-L106}{{\leanfont inclMatrix\_inl}}}{SquareLawFourier.lean}{106--106}
\crow{sumeqnextDatum}{\href{https://github.com/lsernau/grh-lean/blob/72fead102f663d5c7a407964a13dc0cf9d977207/grh-compact/src/New/Exchange/SquareLawFourier.lean\#L111-L127}{{\leanfont SquareLaw.\allowbreak{}sum\_eq\_nextDatum}}}{SquareLawFourier.lean}{111--127}
\crow{blockInlGL}{\href{https://github.com/lsernau/grh-lean/blob/72fead102f663d5c7a407964a13dc0cf9d977207/grh-compact/src/New/Exchange/SquareLawFourier.lean\#L130-L130}{{\leanfont blockInlGL}}}{SquareLawFourier.lean}{130--130}
\crow{graphKernel}{\href{https://github.com/lsernau/grh-lean/blob/72fead102f663d5c7a407964a13dc0cf9d977207/grh-compact/src/New/Exchange/SquareLawFourier.lean\#L134-L135}{{\leanfont graphKernel}}}{SquareLawFourier.lean}{134--135}
\crow{graphKernelinvLastRowsfirstCols}{\href{https://github.com/lsernau/grh-lean/blob/72fead102f663d5c7a407964a13dc0cf9d977207/grh-compact/src/New/Exchange/SquareLawFourier.lean\#L137-L139}{{\leanfont graphKernel\_invLastRows\_firstCols}}}{SquareLawFourier.lean}{137--139}
\crow{existseqofgraphKernelnezero}{\href{https://github.com/lsernau/grh-lean/blob/72fead102f663d5c7a407964a13dc0cf9d977207/grh-compact/src/New/Exchange/SquareLawFourier.lean\#L142-L143}{{\leanfont exists\_eq\_of\_graphKernel\_ne\_zero}}}{SquareLawFourier.lean}{142--143}
\crow{graphKernelequivariant}{\href{https://github.com/lsernau/grh-lean/blob/72fead102f663d5c7a407964a13dc0cf9d977207/grh-compact/src/New/Exchange/SquareLawFourier.lean\#L147-L153}{{\leanfont graphKernel\_equivariant}}}{SquareLawFourier.lean}{147--153}
\crow{graphKernelmulright}{\href{https://github.com/lsernau/grh-lean/blob/72fead102f663d5c7a407964a13dc0cf9d977207/grh-compact/src/New/Exchange/SquareLawFourier.lean\#L155-L156}{{\leanfont graphKernel\_mul\_right}}}{SquareLawFourier.lean}{155--156}
\crow{sumgraphKernelmul}{\href{https://github.com/lsernau/grh-lean/blob/72fead102f663d5c7a407964a13dc0cf9d977207/grh-compact/src/New/Exchange/SquareLawFourier.lean\#L160-L166}{{\leanfont sum\_graphKernel\_mul}}}{SquareLawFourier.lean}{160--166}
\crow{existspartnerkernel}{\href{https://github.com/lsernau/grh-lean/blob/72fead102f663d5c7a407964a13dc0cf9d977207/grh-compact/src/New/Exchange/SquareLawFourier.lean\#L171-L193}{{\leanfont exists\_partner\_kernel}}}{SquareLawFourier.lean}{171--193}
\crow{sumfouriereqnextDatum}{\href{https://github.com/lsernau/grh-lean/blob/72fead102f663d5c7a407964a13dc0cf9d977207/grh-compact/src/New/Exchange/SquareLawFourier.lean\#L202-L220}{{\leanfont SquareLaw.\allowbreak{}sum\_fourier\_eq\_nextDatum}}}{SquareLawFourier.lean}{202--220}
\crow{formNegEquiv}{\href{https://github.com/lsernau/grh-lean/blob/72fead102f663d5c7a407964a13dc0cf9d977207/grh-compact/src/New/Exchange/TriangleArrows.lean\#L27-L30}{{\leanfont formNegEquiv}}}{TriangleArrows.lean}{27--30}
\crow{existsadjointfirstarrow}{\href{https://github.com/lsernau/grh-lean/blob/72fead102f663d5c7a407964a13dc0cf9d977207/grh-compact/src/New/Exchange/TriangleArrows.lean\#L35-L44}{{\leanfont IsEdgeFamily.\allowbreak{}exists\_adjoint\_first\_arrow}}}{TriangleArrows.lean}{35--44}
\crow{TriForm}{\href{https://github.com/lsernau/grh-lean/blob/72fead102f663d5c7a407964a13dc0cf9d977207/grh-compact/src/New/Exchange/TriangleArrows.lean\#L52-L52}{{\leanfont Data.\allowbreak{}TriForm}}}{TriangleArrows.lean}{52--52}
\crow{triLift}{\href{https://github.com/lsernau/grh-lean/blob/72fead102f663d5c7a407964a13dc0cf9d977207/grh-compact/src/New/Exchange/TriangleArrows.lean\#L57-L59}{{\leanfont Data.\allowbreak{}triLift}}}{TriangleArrows.lean}{57--59}
\crow{triNegEquiv}{\href{https://github.com/lsernau/grh-lean/blob/72fead102f663d5c7a407964a13dc0cf9d977207/grh-compact/src/New/Exchange/TriangleArrows.lean\#L63-L64}{{\leanfont Data.\allowbreak{}triNegEquiv}}}{TriangleArrows.lean}{63--64}
\crow{triBase}{\href{https://github.com/lsernau/grh-lean/blob/72fead102f663d5c7a407964a13dc0cf9d977207/grh-compact/src/New/Exchange/TriangleArrows.lean\#L68-L69}{{\leanfont Data.\allowbreak{}triBase}}}{TriangleArrows.lean}{68--69}
\crow{triNegEquivformRep}{\href{https://github.com/lsernau/grh-lean/blob/72fead102f663d5c7a407964a13dc0cf9d977207/grh-compact/src/New/Exchange/TriangleArrows.lean\#L71-L72}{{\leanfont Data.\allowbreak{}triNegEquiv\_formRep}}}{TriangleArrows.lean}{71--72}
\crow{hasSuminnertriLift}{\href{https://github.com/lsernau/grh-lean/blob/72fead102f663d5c7a407964a13dc0cf9d977207/grh-compact/src/New/Exchange/TriangleArrows.lean\#L77-L106}{{\leanfont Data.\allowbreak{}hasSum\_inner\_triLift}}}{TriangleArrows.lean}{77--106}
\crow{existsisTriangle}{\href{https://github.com/lsernau/grh-lean/blob/72fead102f663d5c7a407964a13dc0cf9d977207/grh-compact/src/New/Exchange/TriangleArrows.lean\#L113-L188}{{\leanfont Data.\allowbreak{}exists\_isTriangle}}}{TriangleArrows.lean}{113--188}
\crow{IsIntegralMeasure}{\href{https://github.com/lsernau/grh-lean/blob/72fead102f663d5c7a407964a13dc0cf9d977207/grh-compact/src/New/Exchange/TriangleLevelExponent.lean\#L25-L27}{{\leanfont Data.\allowbreak{}TameChar.\allowbreak{}IsIntegralMeasure}}}{TriangleLevelExponent.lean}{25--27}
\crow{normsq}{\href{https://github.com/lsernau/grh-lean/blob/72fead102f663d5c7a407964a13dc0cf9d977207/grh-compact/src/New/Exchange/TriangleLevelExponent.lean\#L31-L34}{{\leanfont Data.\allowbreak{}TameChar.\allowbreak{}IsIntegralMeasure.\allowbreak{}norm\_sq}}}{TriangleLevelExponent.lean}{31--34}
\crow{normapplyonele}{\href{https://github.com/lsernau/grh-lean/blob/72fead102f663d5c7a407964a13dc0cf9d977207/grh-compact/src/New/Exchange/TriangleLevelExponent.lean\#L40-L44}{{\leanfont Data.\allowbreak{}TameChar.\allowbreak{}IsIntegralMeasure.\allowbreak{}norm\_apply\_one\_le}}}{TriangleLevelExponent.lean}{40--44}
\crow{pointDatum}{\href{https://github.com/lsernau/grh-lean/blob/72fead102f663d5c7a407964a13dc0cf9d977207/grh-compact/src/New/Exchange/TriangleLevelExponent.lean\#L55-L56}{{\leanfont Data.\allowbreak{}Pair.\allowbreak{}pointDatum}}}{TriangleLevelExponent.lean}{55--56}
\crow{pointKernel}{\href{https://github.com/lsernau/grh-lean/blob/72fead102f663d5c7a407964a13dc0cf9d977207/grh-compact/src/New/Exchange/TriangleLevelExponent.lean\#L61-L62}{{\leanfont Data.\allowbreak{}Pair.\allowbreak{}pointKernel}}}{TriangleLevelExponent.lean}{61--62}
\crow{trivialLiftapply}{\href{https://github.com/lsernau/grh-lean/blob/72fead102f663d5c7a407964a13dc0cf9d977207/grh-compact/src/New/Exchange/TriangleLevelExponent.lean\#L67-L74}{{\leanfont Data.\allowbreak{}Pair.\allowbreak{}trivialLift\_apply}}}{TriangleLevelExponent.lean}{67--74}
\crow{weylGaugepointKernelleone}{\href{https://github.com/lsernau/grh-lean/blob/72fead102f663d5c7a407964a13dc0cf9d977207/grh-compact/src/New/Exchange/TriangleLevelExponent.lean\#L80-L86}{{\leanfont Data.\allowbreak{}Pair.\allowbreak{}weylGauge\_pointKernel\_le\_one}}}{TriangleLevelExponent.lean}{80--86}
\crow{isIntegralMeasureintAlgPart}{\href{https://github.com/lsernau/grh-lean/blob/72fead102f663d5c7a407964a13dc0cf9d977207/grh-compact/src/New/Exchange/TriangleLevelExponent.lean\#L94-L107}{{\leanfont Data.\allowbreak{}Pair.\allowbreak{}isIntegralMeasure\_intAlgPart}}}{TriangleLevelExponent.lean}{94--107}
\crow{intAlgParteqmulL}{\href{https://github.com/lsernau/grh-lean/blob/72fead102f663d5c7a407964a13dc0cf9d977207/grh-compact/src/New/Exchange/TriangleLevelExponent.lean\#L115-L121}{{\leanfont Data.\allowbreak{}Pair.\allowbreak{}intAlgPart\_eq\_mul\_L}}}{TriangleLevelExponent.lean}{115--121}
\crow{rationezeroandreal}{\href{https://github.com/lsernau/grh-lean/blob/72fead102f663d5c7a407964a13dc0cf9d977207/grh-compact/src/New/Exchange/TriangleLevelExponent.lean\#L125-L127}{{\leanfont Data.\allowbreak{}Pair.\allowbreak{}ratio\_ne\_zero\_and\_real}}}{TriangleLevelExponent.lean}{125--127}
\crow{ratioeqintAlgPartdiv}{\href{https://github.com/lsernau/grh-lean/blob/72fead102f663d5c7a407964a13dc0cf9d977207/grh-compact/src/New/Exchange/TriangleLevelExponent.lean\#L131-L133}{{\leanfont Data.\allowbreak{}Pair.\allowbreak{}ratio\_eq\_intAlgPart\_div}}}{TriangleLevelExponent.lean}{131--133}
\crow{intAlgPartnezero}{\href{https://github.com/lsernau/grh-lean/blob/72fead102f663d5c7a407964a13dc0cf9d977207/grh-compact/src/New/Exchange/TriangleLevelExponent.lean\#L137-L138}{{\leanfont Data.\allowbreak{}Pair.\allowbreak{}intAlgPart\_ne\_zero}}}{TriangleLevelExponent.lean}{137--138}
\crow{trivialLiftnezero}{\href{https://github.com/lsernau/grh-lean/blob/72fead102f663d5c7a407964a13dc0cf9d977207/grh-compact/src/New/Exchange/TriangleLevelExponent.lean\#L142-L142}{{\leanfont Data.\allowbreak{}Pair.\allowbreak{}trivialLift\_ne\_zero}}}{TriangleLevelExponent.lean}{142--142}
\crow{heightLeofeqzpowmul}{\href{https://github.com/lsernau/grh-lean/blob/72fead102f663d5c7a407964a13dc0cf9d977207/grh-compact/src/New/Exchange/TriangleLevelExponent.lean\#L151-L154}{{\leanfont Data.\allowbreak{}IsFitFamily.\allowbreak{}heightLe\_of\_eq\_zpow\_mul}}}{TriangleLevelExponent.lean}{151--154}
\crow{lowerPoint}{\href{https://github.com/lsernau/grh-lean/blob/72fead102f663d5c7a407964a13dc0cf9d977207/grh-compact/src/New/Exchange/TriangleLevelExponent.lean\#L160-L161}{{\leanfont Data.\allowbreak{}Pair.\allowbreak{}lowerPoint}}}{TriangleLevelExponent.lean}{160--161}
\crow{lValueProdeq}{\href{https://github.com/lsernau/grh-lean/blob/72fead102f663d5c7a407964a13dc0cf9d977207/grh-compact/src/New/Exchange/TriangleLevelExponent.lean\#L166-L173}{{\leanfont Data.\allowbreak{}Pair.\allowbreak{}lValueProd\_eq}}}{TriangleLevelExponent.lean}{166--173}
\crow{IsTriangle}{\href{https://github.com/lsernau/grh-lean/blob/72fead102f663d5c7a407964a13dc0cf9d977207/grh-compact/src/New/Exchange/TriangleNorms.lean\#L20-L27}{{\leanfont IsTriangle}}}{TriangleNorms.lean}{20--27}
\crow{weightedChain}{\href{https://github.com/lsernau/grh-lean/blob/72fead102f663d5c7a407964a13dc0cf9d977207/grh-compact/src/New/Exchange/TriangleNorms.lean\#L31-L31}{{\leanfont weightedChain}}}{TriangleNorms.lean}{31--31}
\crow{weightedChaindescent}{\href{https://github.com/lsernau/grh-lean/blob/72fead102f663d5c7a407964a13dc0cf9d977207/grh-compact/src/New/Exchange/TriangleNorms.lean\#L36-L43}{{\leanfont IsTriangle.\allowbreak{}weightedChain\_descent}}}{TriangleNorms.lean}{36--43}
\crow{eqzeroofweightedChain}{\href{https://github.com/lsernau/grh-lean/blob/72fead102f663d5c7a407964a13dc0cf9d977207/grh-compact/src/New/Exchange/TriangleNorms.lean\#L47-L49}{{\leanfont IsTriangle.\allowbreak{}eq\_zero\_of\_weightedChain}}}{TriangleNorms.lean}{47--49}
\crow{innerweightedChaineq}{\href{https://github.com/lsernau/grh-lean/blob/72fead102f663d5c7a407964a13dc0cf9d977207/grh-compact/src/New/Exchange/TriangleNorms.lean\#L54-L66}{{\leanfont IsTriangle.\allowbreak{}inner\_weightedChain\_eq}}}{TriangleNorms.lean}{54--66}}
% the notations of the defined functions: one slot for each explicit argument of the Lean declaration
\newcommand{\LValueData}[2]{\operatorname{LD}(#1,#2)}
\newcommand{\auxPair}[1]{\operatorname{aux}(#1)}
\newcommand{\kernelValue}[5]{\operatorname{kv}_{#1}(#2;#3)(#4,#5)}
\newcommand{\liftKernel}[1]{\operatorname{lk}(#1)}
\newcommand{\kernelForm}[2]{\operatorname{kf}_{#1}(#2)}
\newcommand{\chosenLevel}[1]{\operatorname{chl}(#1)}
\newcommand{\trivialLift}[1]{\Xi(#1)}
\newcommand{\HasLiftRank}[2]{\operatorname{LiftRank}(#1,#2)}
\newcommand{\IsFlat}[2]{\operatorname{Flat}(#1,#2)}
\newcommand{\denomBound}[1]{\operatorname{den}(#1)}
\newcommand{\lValueData}[2]{\operatorname{ld}(#1;#2)}
\newcommand{\trivialLifta}[2]{\Xi^{\mathrm{sc}}(#1;#2)}
\newcommand{\intAlgPart}[2]{\operatorname{alg}(#1;#2)}
\newcommand{\CostsSquares}[3]{\operatorname{Cost}_{#1}(#2;#3)}
\newcommand{\homX}[1]{\operatorname{hom}^{X}(#1)}
\newcommand{\liftKernela}[2]{\operatorname{lk}^{\mathrm{col}}(#1,#2)}
\newcommand{\translateDatum}[2]{\operatorname{td}(#1,#2)}
\newcommand{\chainComb}[4]{\operatorname{cc}(#1;#2,#3;#4)}
\newcommand{\Bmu}[2]{\operatorname{B}_{#1}(#2)}
\newcommand{\inducedPlace}[1]{\operatorname{ind}(#1)}
\newcommand{\IsMarked}[1]{\operatorname{Marked}(#1)}
\newcommand{\IsFullMark}[1]{\operatorname{FullMark}(#1)}
\newcommand{\evenPlaces}[1]{\operatorname{ev}(#1)}
\newcommand{\markExponent}[1]{\operatorname{mexp}(#1)}
\newcommand{\lowColumns}[2]{\operatorname{lc}(#1,#2)}
\newcommand{\lowBoxes}[2]{\operatorname{lb}(#1,#2)}
\newcommand{\SquareLaw}[5]{\operatorname{SqLaw}(#1,#2;#3;#4,#5)}
\newcommand{\inclusionDatum}[2]{\mathbb{1}^{\mathrm{incl}}_{#1}(#2)}
\newcommand{\nextDatum}[2]{\operatorname{next}(#1)(#2)}
\newcommand{\blockInlGL}[1]{\operatorname{bl}_{#1}}
\newcommand{\graphKernel}[4]{\operatorname{gr}_{#1,#2}(#3,#4)}
\newcommand{\formNegEquiv}[2]{\operatorname{opp}_{#1,#2}}
\newcommand{\TriForm}[4]{\operatorname{TF}_{#1}(#2;#3,#4)}
\newcommand{\triLift}[6]{\operatorname{tl}_{#1,#2,#3}(#4,#5)(#6)}
\newcommand{\triNegEquiv}[3]{\operatorname{tn}_{#1,#2}(#3)}
\newcommand{\triBase}[2]{\operatorname{TB}_{#1}(#2)}
\newcommand{\IsIntegralMeasure}[2]{\operatorname{IntMeas}(#1;#2)}
\newcommand{\pointDatum}[2]{\operatorname{pd}(#1;#2)}
\newcommand{\pointKernel}[2]{\operatorname{pk}(#1;#2)}
\newcommand{\lowerPoint}[2]{\operatorname{low}(#1,#2)}
\newcommand{\IsTriangle}[7]{\operatorname{Tri}(#1;#2,#3;#4,#5;#6;#7)}
\newcommand{\weightedChain}[4]{\operatorname{wc}(#1,#2;#3,#4)}
\newcommand{\UnitMembers}{305}
\newcommand{\UnitInNote}{129}
\newcommand{\UnitProofs}{122}
\newcommand{\UnitSteps}{879}
\newcommand{\UnitAllSteps}{933}


\linespread{1.08}                       % 11 pt on 14.7 pt
\setlength{\parindent}{0pt}
\setlength{\parskip}{5.5pt plus 1.5pt minus 0.5pt}
\setlength{\abovedisplayskip}{7.5pt plus 2pt minus 3pt}
\setlength{\belowdisplayskip}{7.5pt plus 2pt minus 3pt}
\setlength{\abovedisplayshortskip}{3pt plus 2pt}
\setlength{\belowdisplayshortskip}{6pt plus 2pt minus 2pt}
\emergencystretch=1.5em
\clubpenalty=10000 \widowpenalty=10000 \displaywidowpenalty=10000
\relpenalty=7000 \binoppenalty=9000
\raggedbottom
\pagestyle{fancy}\fancyhf{}\renewcommand{\headrulewidth}{0pt}
\fancyfoot[C]{\sffamily\fontsize{7.6}{9}\selectfont\color{inkfaint}\thepage\,/\,\pageref*{LastPage}}
\setlist{topsep=2pt,partopsep=0pt,parsep=0pt,itemsep=4pt,leftmargin=2.1em,labelsep=0.55em}

\newcommand{\Q}{\mathbb{Q}}
\newcommand{\Z}{\mathbb{Z}}
\newcommand{\F}{\mathbb{F}}
\newcommand{\len}{\operatorname{len}}
\newcommand{\cok}{\operatorname{cok}}
\newcommand{\End}{\operatorname{End}}
\newcommand{\Hom}{\operatorname{Hom}}
\newcommand{\Gal}{\operatorname{Gal}}
\newcommand{\diag}{\operatorname{diag}}
\newcommand{\kw}[1]{\emph{#1}}
\newcommand{\m}[1]{\mbox{$#1$}}        % a short formula that must not be broken across lines
\newcommand{\head}[1]{\par\addvspace{17pt plus 4pt minus 2pt}\Needspace{6\baselineskip}%
  \noindent{\fontsize{13.2}{16}\selectfont\bfseries\color{ink}#1}\par\nopagebreak\vspace{-1pt}}
\newcommand{\sect}[2]{\refstepcounter{section}\label{#2}\par\addvspace{19pt plus 4pt minus 2pt}\Needspace{7\baselineskip}%
  \noindent\llap{{\fontsize{13.2}{16}\selectfont\bfseries\color{inkfaint}\thesection}\hspace{9pt}}%
  {\fontsize{13.2}{16}\selectfont\bfseries\color{ink}#1}\par\nopagebreak\vspace{-1pt}}
\newcommand{\subhead}[1]{\par\addvspace{9pt plus 2pt}\Needspace{4\baselineskip}\noindent\textbf{#1}\enspace\ignorespaces}
% where a statement lives in the Lean archive
\newcommand{\leantag}{{\sffamily\bfseries\fontsize{6.3}{8}\selectfont\color{leancol}\addfontfeatures{LetterSpace=7}LEAN}}
\newcommand{\sep}{\unskip\nobreak\hspace{5pt}{\color{inkfaint}\textperiodcentered}\hspace{5pt plus 2pt}\ignorespaces}
\newcommand{\leanline}[1]{\par\nopagebreak\vspace{-2.5pt}{\fontsize{8.9}{12.4}\selectfont\color{inksoft}\raggedright
  \noindent\llap{\leantag\hspace{9pt}}#1\par}\vspace{1pt}}
% a displayed statement in a tinted box: #1 colour, #2 text
\newcommand{\statement}[2]{\par\addvspace{7pt}\begingroup\setlength{\fboxsep}{8.5pt}\setlength{\fboxrule}{0.6pt}%
  \noindent\fcolorbox{#1!30!white}{#1!3!white}{\parbox{\dimexpr\linewidth-2\fboxsep-2\fboxrule\relax}{#2}}%
  \endgroup\par\addvspace{7pt}}
\newcommand{\nt}[2]{\par\noindent\hangindent=1.1em\hangafter=1 #1\enspace{\color{inksoft}#2}\par\vspace{1pt}}
\newcommand{\legdot}[1]{\tikz[baseline=-0.62ex]\fill[#1] (0,0) circle (2.5pt);\hspace{3pt}}
\newcommand{\crow}[4]{\csname CS@#1\endcsname & #2 & \textcolor{inksoft}{#3} & \textcolor{inksoft}{#4}\\}
\newcommand{\leanquoteline}[2]{\par\noindent\hangindent=\dimexpr#1\fontdimen2\font+2.2em\relax\hangafter=1
  \hspace*{#1\fontdimen2\font}#2}
\newcommand{\quoted}[1]{\par\vspace{3pt}{\ttfamily\fontsize{7.5}{10.4}\selectfont\color{leancol}\raggedright
  \setlength{\parskip}{0pt}\leanquote{#1}\par}\vspace{4pt}}
\newcommand{\bx}[3]{{\sffamily\bfseries\fontsize{7.3}{9.2}\selectfont\color{#1!82!black}#2}\\[1.5pt]#3}
\newcommand{\cs}[2]{\expandafter\newcommand\csname CS@#1\endcsname{#2}}

% theorems
\newtheoremstyle{plainx}{7pt plus 2pt}{7pt plus 2pt}{\itshape}{}{\bfseries}{.}{0.5em}{}
\newtheoremstyle{uprightx}{7pt plus 2pt}{7pt plus 2pt}{}{}{\bfseries}{.}{0.5em}{}
\theoremstyle{plainx}
\newtheorem{theorem}{Theorem}[section]
\newtheorem{lemma}[theorem]{Lemma}
\newtheorem{proposition}[theorem]{Proposition}
\newtheorem{corollary}[theorem]{Corollary}
\theoremstyle{uprightx}
\newtheorem{definition}[theorem]{Definition}
\newtheorem{setting}[theorem]{Setting}          % hypotheses shared by several statements; each of them cites it by number
\newtheorem{remark}[theorem]{Remark}
\newtheorem{example}[theorem]{Example}
\renewcommand{\qedsymbol}{{\setlength{\fboxsep}{0pt}\setlength{\fboxrule}{0.5pt}\fbox{\rule{0pt}{0.46em}\hspace{0.46em}}}}
% what holds throughout the note: the ambient objects only (implicit and instance arguments), never a hypothesis
\newenvironment{conventions}{\par\addvspace{5pt}\noindent\textbf{Conventions.}\enspace\ignorespaces}{\par\addvspace{3pt}}
\makeatletter
\renewenvironment{proof}[1][Proof]{\par\pushQED{\qed}\normalfont\topsep0pt\partopsep0pt
  \trivlist\item[\hskip\labelsep\itshape #1.]\ignorespaces}{\popQED\endtrivlist\@endpefalse}
\makeatother

% where each declaration of the concordance is treated: \cs{key}{place}, the place naming a \label or a lettered item
\cs{exchangetheorem}{Theorem 1}
\cs{onesquare}{Theorem 2}
\cs{existsisTriangle}{Theorem 3}
\cs{iotaAtLevelinjective}{Theorem 4}
\cs{intAlgParteqmulL}{Theorem 5}
\cs{isIntegralMeasureintAlgPart}{Theorem 6}
\cs{SquareLaw}{Definition~\ref{d:square}}
\cs{inclusionDatum}{Definition~\ref{d:square}}
\cs{nextDatum}{Definition~\ref{d:square}}
\cs{swapFirstTwo}{Definition~\ref{d:square}}
\cs{blockInlGL}{Definition~\ref{d:square}}
\cs{graphKernel}{Definition~\ref{d:square}}
\cs{sumeqsumfromRows}{Lemma~\ref{l:elem}}
\cs{sumeqsumfromCols}{Lemma~\ref{l:elem}}
\cs{injectiveglinvmulmul}{Lemma~\ref{l:elem}}
\cs{inclMatrixinl}{Lemma~\ref{l:elem}}
\cs{graphKernelinvLastRowsfirstCols}{Lemma~\ref{l:elem}}
\cs{existseqofgraphKernelnezero}{Lemma~\ref{l:elem}}
\cs{graphKernelmulright}{Lemma~\ref{l:elem}}
\cs{sumeqnextDatum}{Lemma~\ref{l:next}}
\cs{graphKernelequivariant}{Lemma~\ref{l:greq}}
\cs{sumgraphKernelmul}{Lemma~\ref{l:greq}}
\cs{existspartnerkernel}{Lemma~\ref{l:pk}}
\cs{sumfouriereqnextDatum}{Lemma~\ref{l:fourier}}
\cs{isFormnegiffandintegralneg}{Lemma~\ref{l:neg}}
\cs{thetaLargemove}{Lemma~\ref{l:move}}
\cs{existsisBaseAdapted}{Lemma~\ref{l:planes}}
\cs{slicezerocongr}{Lemma~\ref{l:planes}}
\cs{slicezerocomp}{Lemma~\ref{l:planes}}
\cs{namevacElemofisEmpty}{Lemma~\ref{l:planes}}
\cs{existsstdExchangeremovingplanes}{Proposition~\ref{p:descent}}
\cs{exchangeofpartners}{Proposition~\ref{p:partners}}
\cs{existsequalRankThetaRatiovalue}{Proposition~\ref{p:equal}}
\cs{integralintegralsumeq}{Lemma~\ref{l:tools}}
\cs{applyeqmulofformRepone}{Lemma~\ref{l:tools}}
\cs{existsintegrablekernelbound}{Lemma~\ref{l:tools}}
\cs{IsTriangle}{Definition~\ref{d:tri}}
\cs{weightedChain}{Definition~\ref{d:tri}}
\cs{weightedChaindescent}{Lemma~\ref{l:chain}}
\cs{eqzeroofweightedChain}{Lemma~\ref{l:chain}}
\cs{innerweightedChaineq}{Lemma~\ref{l:chain}}
\cs{CostsSquares}{Definition~\ref{d:cost}}
\cs{mul}{Lemma~\ref{l:cost}}
\cs{card}{Lemma~\ref{l:cost}}
\cs{ofnormsqeqmulcardzpow}{Lemma~\ref{l:cost}}
\cs{one}{Lemma~\ref{l:cost}}
\cs{isPrimitiveMulCharofisTameFor}{Lemma~\ref{l:cost}}
\cs{existsisTameForcostsSquaresone}{Lemma~\ref{l:cost}}
\cs{chosenIndexEmb}{Definition~\ref{d:ldata}}
\cs{LValueData}{Definition~\ref{d:ldata}}
\cs{auxPair}{Definition~\ref{d:ldata}}
\cs{kernelValue}{Definition~\ref{d:ldata}}
\cs{liftKernel}{Definition~\ref{d:ldata}}
\cs{kernelForm}{Definition~\ref{d:ldata}}
\cs{chosenLevel}{Definition~\ref{d:ldata}}
\cs{trivialLift}{Definition~\ref{d:ldata}}
\cs{ahasWeight}{Definition~\ref{d:ldata}}
\cs{bhasWeight}{Definition~\ref{d:ldata}}
\cs{HasLiftRank}{Definition~\ref{d:pairs}}
\cs{IsFlat}{Definition~\ref{d:pairs}}
\cs{denomBound}{Definition~\ref{d:pairs}}
\cs{lValueData}{Definition~\ref{d:pairs}}
\cs{trivialLifta}{Definition~\ref{d:pairs}}
\cs{intAlgPart}{Definition~\ref{d:pairs}}
\cs{pointDatum}{Definition~\ref{d:pairs}}
\cs{pointKernel}{Definition~\ref{d:pairs}}
\cs{lowerPoint}{Definition~\ref{d:pairs}}
\cs{IsIntegralMeasure}{Definition~\ref{d:intmeas}}
\cs{cardYle}{Definition~\ref{d:intmeas}}
\cs{cardV}{Definition~\ref{d:intmeas}}
\cs{cardfilterbasisNumber}{Definition~\ref{d:intmeas}}
\cs{varsigmastrictAnti}{Definition~\ref{d:intmeas}}
\cs{paramWltiff}{Definition~\ref{d:intmeas}}
\cs{cardnegsignsW}{Definition~\ref{d:intmeas}}
\cs{isLiftWeight}{Definition~\ref{d:intmeas}}
\cs{conjWeylRepeqone}{Definition~\ref{d:intmeas}}
\cs{thetaeqinnerkernelForm}{Definition~\ref{d:intmeas}}
\cs{innertrivialLift}{Definition~\ref{d:intmeas}}
\cs{ratioeqintAlgPartdiv}{Definition~\ref{d:intmeas}}
\cs{intAlgPartnezero}{Definition~\ref{d:intmeas}}
\cs{trivialLiftnezero}{Definition~\ref{d:intmeas}}
\cs{triNegEquivformRep}{Definition~\ref{d:intmeas}}
\cs{cardresRingle}{Definition~\ref{d:intmeas}}
\cs{stdLevelIndexlastResLeveleq}{Definition~\ref{d:intmeas}}
\cs{isHCParamY}{Lemma~\ref{l:params}}
\cs{fitsParamsauxPair}{Lemma~\ref{l:params}}
\cs{innereqpetersson}{Lemma~\ref{l:petersson}}
\cs{kernelForminvariant}{Lemma~\ref{l:kform}}
\cs{trivialLifteqstdClassSumtranslate}{Lemma~\ref{l:kform}}
\cs{thetapairWeilfockDatum}{Lemma~\ref{l:kform}}
\cs{kernelFormofjetDegeqzero}{Lemma~\ref{l:kform}}
\cs{toPolyeqrenameharmonicPoly}{Lemma~\ref{l:poly}}
\cs{harmonicPolyeqschurCoeffPoly}{Lemma~\ref{l:poly}}
\cs{hookNumbereqprodfactorial}{Lemma~\ref{l:poly}}
\cs{trivialLiftapply}{Lemma~\ref{l:apply}}
\cs{weylGaugepointKernelleone}{Lemma~\ref{l:apply}}
\cs{normsq}{Lemma~\ref{l:meas}}
\cs{normapplyonele}{Lemma~\ref{l:meas}}
\cs{heightLeofeqzpowmul}{Lemma~\ref{l:meas}}
\cs{rationezeroandreal}{Lemma~\ref{l:ratio}}
\cs{lValueProdeq}{Lemma~\ref{l:ratio}}
\cs{maptrivialLift}{Proposition~\ref{p:galois}}
\cs{lognormmapnormsqtrivialLift}{Proposition~\ref{p:galois}}
\cs{formNegEquiv}{Definition~\ref{d:triforms}}
\cs{TriForm}{Definition~\ref{d:triforms}}
\cs{triLift}{Definition~\ref{d:triforms}}
\cs{triNegEquiv}{Definition~\ref{d:triforms}}
\cs{triBase}{Definition~\ref{d:triforms}}
\cs{existsadjointfirstarrow}{Lemma~\ref{l:first}}
\cs{hasSuminnertriLift}{Lemma~\ref{l:first}}
\cs{inneriotaAtLevel}{Lemma~\ref{l:chainnorm}}
\cs{homX}{Definition~\ref{d:crude}}
\cs{liftKernela}{Definition~\ref{d:crude}}
\cs{Translations}{Definition~\ref{d:crude}}
\cs{translateDatum}{Definition~\ref{d:crude}}
\cs{chainComb}{Definition~\ref{d:crude}}
\cs{B}{Definition~\ref{d:crude}}
\cs{inducedPlace}{Definition~\ref{d:crude}}
\cs{IsMarked}{Definition~\ref{d:crude}}
\cs{IsFullMark}{Definition~\ref{d:crude}}
\cs{evenPlaces}{Definition~\ref{d:crude}}
\cs{markExponent}{Definition~\ref{d:crude}}
\cs{integralColumns}{Definition~\ref{d:crude}}
\cs{lowColumns}{Definition~\ref{d:crude}}
\cs{lowBoxes}{Definition~\ref{d:crude}}
\cs{isLiftWeightkappa}{Lemma~\ref{l:std}}
\cs{spanpiGBar}{Lemma~\ref{l:std}}
\cs{smulvacuumDatummemstandardData}{Lemma~\ref{l:std}}
\cs{isAdmissibleCombvacuumDatum}{Lemma~\ref{l:std}}
% letters that the coordinator may change in one place
\newcommand{\dL}{\check{\Psi}}         % a function on the large side (Lean Φ; the letter Φ is taken by Φ_{d,M})
\newcommand{\dS}{\Psi}
\newcommand{\sR}{\mathsf{R}}
\newcommand{\sF}{\mathsf{F}}
\newcommand{\ba}{\mathbf{a}}
\newcommand{\bb}{\mathbf{b}}
\newcommand{\bm}{\mathbf{m}}
\newcommand{\bn}{\mathbf{n}}
\newcommand{\br}{\mathbf{r}}
\newcommand{\bp}{\mathbf{p}}
\newcommand{\bN}{\mathbf{N}}
\newcommand{\bT}{\mathbf{T}}
\newcommand{\sW}{\mathsf{W}}
\newcommand{\sV}{\mathsf{V}}
\newcommand{\sX}{\mathsf{X}}
\newcommand{\sU}{\mathsf{U}}
\newcommand{\sA}{\mathsf{A}}
\newcommand{\sS}{\mathsf{S}}
\newcommand{\sT}{\mathsf{T}}
\newcommand{\sI}{\mathsf{I}}
\newcommand{\sY}{\mathsf{Y}}
\newcommand{\se}{\mathsf{e}}
\newcommand{\sg}{\mathsf{g}}
\newcommand{\sh}{\mathsf{h}}
\newcommand{\sk}{\mathsf{k}}
\newcommand{\sx}{\mathsf{x}}
\newcommand{\sy}{\mathsf{y}}
\newcommand{\sz}{\mathsf{z}}
\newcommand{\sff}{\mathsf{f}}
\newcommand{\sq}{\mathsf{q}}
\newcommand{\pp}{\mathsf{p}}
\newcommand{\fc}{\mathfrak{c}}
\newcommand{\fp}{\mathfrak{p}}
\newcommand{\fm}{\mathfrak{m}}
\newcommand{\fv}{\mathfrak{v}}
\newcommand{\fx}{\mathfrak{x}}
\newcommand{\fy}{\mathfrak{y}}
\newcommand{\fS}{\mathfrak{S}}
\newcommand{\fa}{\mathfrak{a}}
\newcommand{\fl}{\mathfrak{l}}
\newcommand{\rc}{\mathrm{c}}
\newcommand{\rC}{\mathrm{C}}
\newcommand{\A}{\mathbb{A}}
\newcommand{\CC}{\mathbb{C}}
\newcommand{\GL}{\operatorname{GL}}
\newcommand{\U}{\operatorname{U}}
\newcommand{\tr}{\operatorname{tr}}
\newcommand{\cL}{\mathcal{L}}
\newcommand{\cI}{\mathcal{I}}
\newcommand{\cF}{\mathcal{F}}
\newcommand{\cV}{\mathcal{V}}
\newcommand{\cM}{\mathcal{M}}
\newcommand{\cO}{\mathcal{O}}
\newcommand{\Mat}[2]{\operatorname{M}_{#1\times#2}}
\newcommand{\rowk}[2]{\operatorname{row}_{#1}(#2)}
\newcommand{\colk}[2]{\operatorname{col}_{#1}(#2)}
\newcommand{\stc}[2]{\dS^{\mathrm{st}}_{#1,#2}}
\newcommand{\str}[3]{\dL^{\mathrm{st}}_{#1,#2,#3}}
\newcommand{\thL}[5]{\theta^{\mathrm{lg}}_{#1,#2}(#3;#4,#5)}
\newcommand{\thS}[5]{\theta^{\mathrm{sm}}_{#1,#2}(#3;#4,#5)}
\newcommand{\thI}[5]{\operatorname{I}_{#1,#2}(#3,#4;#5)}
\newcommand{\eqr}[3]{\varrho(#1;#2,#3)}
\newcommand{\Exch}[8]{\operatorname{Exch}(#1,#2,#3;#4,#5;#6,#7;#8)}
\newcommand{\LP}[2]{\Pi(#1,#2)}
\newcommand{\LPP}[3]{\Pi(#1,#2,#3)}
\newcommand{\negc}{\operatorname{neg}}
\newcommand{\Np}{\operatorname{N}}
\newcommand{\pet}[2]{\langle #1,#2\rangle_{\mathrm{Pet}}}
\newcommand{\ip}[2]{\langle #1,#2\rangle}
\newcommand{\md}[3]{#1\otimes[#2]_{#3}}

\begin{document}
\color{ink}

{\raggedright
{\sffamily\bfseries\fontsize{7.2}{9}\selectfont\color{leancol}\addfontfeatures{LetterSpace=8}READBACK NOTE 17\par}
\vspace{5pt}
{\fontsize{21.5}{25.5}\selectfont\bfseries The exchange theorem, one square of lifts, and the triangle of lifts and
descents\par}
\vspace{5pt}
{\fontsize{12.4}{16}\selectfont\itshape\color{inksoft} One unit of the Lean artifact \ArtifactName, read back into
mathematics\par}}
\vspace{8pt}
{\color{rulegrey}\hrule height 0.6pt}
\vspace{7pt}
\thispagestyle{empty}

\textbf{Summary.} The exchange theorem (Theorem~1) says: for hermitian spaces $\sV$, $\sV'$ whose ranks add up to
the rank of $\sW$, the theta integrals over $[\U(\sV)]$ and $[\U(\sV')]$ of data that are partners modulo a
modulus are proportional, and the square of the absolute value of the constant is a fixed number times a product
of partial Hecke $L$-values. It is obtained from the identity at the top of a tower (first note of the unit) by
removing hyperbolic planes. Applied to a square of spaces it shows that a theta lift followed by a descent is
proportional to a descent followed by a lift (Theorem~2). The squares fit into a triangle of lifts and descents
attached to the data (Theorem~3), from which the composite of the $m$ lifts is injective (Theorem~4). The norm of
the lift of the trivial character from $\U(1)$ is a Hecke $L$-value, interpolated by one integral measure
(Theorems 5 and~6).

\statement{leancol}{\fontsize{9.6}{13.2}\selectfont\textbf{What this note is.} The artifact \ArtifactName{} is a
body of Lean code: the compact edition of a formal proof, attributed by its authors to L.~Alpöge with the AI system
Claude, whose final theorems are the Riemann hypothesis and its generalisation to Dirichlet $L$-functions. That
headline is a claim. The authors report that every file compiles and that Lean's kernel accepted the proofs; no
mathematician has read the argument through, and nothing here says that the Riemann hypothesis is proved. This
note is the second of three on one unit of the artifact, the unit of the exchange identity, which belongs to the
proof of Theorem~C of the handout. It treats \UnitInNote{} of the unit's \UnitMembers{} declarations; the notes
\emph{Partner data, vanishing on standard data, and the exchange identity at the top of the tower} (the first
note) and \emph{The grid of lifts and descents, and the crude rows} (the third note) treat the others. The proofs
are written from the structured proofs of the unit's dossier, which a program compared with the Lean text; what
they take from other units and from the two other notes is listed as Inputs in §\ref{s:setting} and is not proved
here. The writer compared with their Lean statements: Theorems 1, 2, 3, 5 and~6; Lemma~\ref{l:neg},
Lemma~\ref{l:first}, Lemma~\ref{l:cost}, Lemma~\ref{l:meas}\,(c), Lemma~\ref{l:pk} and
Propositions \ref{p:descent}, \ref{p:partners} and~\ref{p:equal}; and the definitions of the square law, of
triangles and weighted chains, of the forms of the triangle, of tameness for a modulus, of costs, of pairs, of the
algebraic part, of integral measures and of marks. The other statements, Theorem~4 among them, are written from
the dossier's readings; §\ref{s:notes} says so again. A second reader compared all statements of the unit and a passage of every structured proof with the Lean text after the dossier was made; where her reading differs from the dossier's, this note follows hers (§\ref{s:notes}). The Lean text read is
\PinText; Appendix~A is the concordance of \ConcDecls{} declarations.}

\sect{The statements}{s:intro}
The notions are recalled in §\ref{s:setting}. In short: a hermitian space $\sW$ over a CM field $E$ has an index
set and signs; partners $\dL$, $\dS$ for a character $\chi$ are functions on $\ba\times\bn$ and on $\bn\times\bb$
matrices that correspond under the duality between kernels and images; $\md{\fx}{\dL}{\fm}$ is the datum with real
part $\fx$ and finite part given by the function $\dL$ modulo a modulus $\fm$; and
$\LPP{\Sigma}{\omega}{e}=(\prod_{\fp\in\Sigma}\Np\fp)^{e^2}\prod_{i=1}^{e}|L^{\Sigma\cup\bar\Sigma}(i,\omega^{\mathrm{bc}})|^2$,
where $L^{\Sigma\cup\bar\Sigma}(s,\omega^{\mathrm{bc}})$ is, as in the first note, the Hecke $L$-function of the base
change $\omega^{\mathrm{bc}}(\mathsf{x})=\omega(\mathsf{x}/\bar{\mathsf{x}})$ of $\omega$ to $E$ without the Euler factors
at the places of $\Sigma$ and at their conjugates $\bar\Sigma$ (for real $s>1$ the Euler product, at $s=1$ its limit
from the right); for $e=0$ both products are $1$. This is the only $L$-function of the note and it is always
written so.

\statement{arith}{\phantomsection\label{t:1}\textbf{Theorem 1} (the exchange theorem).\enspace \emph{Let $E$ be a CM
field with a CM type and Weil data, with splitting character $\upsilon$. Let $\sW$, $\sV$, $\sV'$ be hermitian
spaces over $E$ with index sets $\bn$, $\ba$, $\bb$, and $\se\colon\bb\sqcup\ba\simeq\bn$ an identification. Suppose
that $\negc(\sW)=\negc(\sV)+\negc(\sV')$, that $|\ba|=|\bb|+2e$ for a natural number $e$, and that
$2\negc(\sW)+2\le|\bn|$. Let $\mu$ be integer weights with $|\bn|\le2\mu(\tau)+2\negc(\sW)$ for every infinite
place $\tau$, and $\fv$ a rapid real datum for the pair $(\sV',\sW)$. Then there are a real number $\rC$ and a
rapid real datum $\fx$ for $(\sW,\sV)$ such that, for every modulus $\fm$, every character $\omega$ of $\U(1)(\A)$
and every multiplicative character $\chi$ of $\cO_E/\fm$ for which $\omega$ is tame for $\fm$ with weights $\mu$
and character $\chi$, writing $\omega^{\circ}=\omega^{-1}\upsilon^{\lfloor|\bn|/2\rfloor}$, there is a complex number
$\rc\ne0$ with}
\[|\rc|^2=\rC\cdot\LPP{\Sigma(\fm)}{\omega}{e}=\rC\,\Bigl(\prod_{\fp\in\Sigma(\fm)}\Np\fp\Bigr)^{e^2}\,\prod_{i=1}^{e}\bigl|L^{\Sigma(\fm)\cup\overline{\Sigma(\fm)}}(i,\omega^{\mathrm{bc}})\bigr|^2,\]
\emph{$\Sigma(\fm)$ the set of places of $\fm$ and $\overline{\Sigma(\fm)}$ the set of their conjugates, such that for all partners $\dL$, $\dS$ for $\chi$ over $\cO_E/\fm$
with respect to $\se$ and every point $\sg$ of $\U(\sW)(\A)$}
\[\thI{\sW}{\sV}{\omega^{\circ}}{1}{\md{\fx}{\dL}{\fm}}(\sg)
=\rc\int_{[\U(\sV')]}\theta_{\sV',\sW}\bigl(\omega^{\circ},\omega^{\circ};\md{\fv}{\dS}{\fm}\bigr)(\sh,\sg)\,d\sh .\]}
\leanline{\lean{exchangetheorem}}

\statement{arith}{\phantomsection\label{t:2}\textbf{Theorem 2} (one square).\enspace \emph{Let $E$ be a CM field with
a CM type and Weil data. Let $\sA'$, $\sA$, $\sU_1$, $\sU_2$ be hermitian spaces over $E$ with index sets $\bn$,
$\bn'$, $\bn_1$, $\bn_2$. Let $\cI$ be an index set with families of characters $\omega^{\mathrm{a}}_i$,
$\omega^{\mathrm{b}}_i$ (each a character for every rank), characters $\omega_i$, moduli $\fm_i$ and
multiplicative characters $\chi_i$ of $\cO_E/\fm_i$. Suppose given four edge families: for $\omega^{\mathrm b}$
from the type $\mathrm{t}_2$ to the type $\mathrm{t}$ in a direction $\mathrm{d}$; for $\omega^{\mathrm a}$ from $\mathrm{t}_1$ to $\mathrm{t}$ in
the opposite direction; for $\omega^{\mathrm a}$ from $\mathrm{t}'$ to $\mathrm{t}_2$ in the opposite direction; and for
$\omega^{\mathrm b}$ from $\mathrm{t}'$ to $\mathrm{t}_1$ in the direction $\mathrm{d}$. Let $\se$ be an identification
$\bn\sqcup\bn'\simeq\bn_1\sqcup\bn_2$, and assume: $\negc(\sU_1\oplus\sU_2)=\negc(\sA)+\negc(\sA')$; $|\bn'|=|\bn|+2e$ for a
natural number $e$; $2\negc(\sU_1\oplus\sU_2)+2\le|\bn_1|+|\bn_2|$; $\mu$ are integer weights with
$|\bn_1|+|\bn_2|\le2\mu(\tau)+2\negc(\sU_1\oplus\sU_2)$ for every infinite place $\tau$; every $\omega_i$ is tame
for $\fm_i$ with weights $\mu$ and character $\chi_i$; and
$\omega^{\mathrm a}_i\cdot\omega^{\mathrm b}_i=\omega_i^{-1}\upsilon^{\lfloor(|\bn_1|+|\bn_2|)/2\rfloor}$ for every
$i$, the two characters being taken in ranks that add up to $|\bn_1|+|\bn_2|$. For a matrix $\sx$ with column
blocks $\sx_1$, $\sx_2$ and a matrix $\sy$ with row blocks $\sy_1$, $\sy_2$ let}
\begin{align*}
\mathsf{u}_i(\sx)&=(\text{descent from }\sA\text{ to }\sU_2\text{ at }\sx_2)\circ(\text{lift from }\sU_1\text{ to }\sA\text{ at }\sx_1),\\
\mathsf{v}_i(\sy)&=(\text{lift from }\sA'\text{ to }\sU_2\text{ at }\sy_2)\circ(\text{descent from }\sU_1\text{ to }\sA'\text{ at }\sy_1)
\end{align*}
\emph{be the composites of the theta lifts and descents of the four edges, with the characters of the index $i$
and the finite data that are the indicators of the blocks modulo $\fm_i$. Then there are a real number $\rC$ and
a complex number $\varkappa$ such that every $i\in\cI$ has a complex number $\rc\ne0$ with
$|\rc|^2=\rC\cdot\LPP{\Sigma(\fm_i)}{\omega_i}{e}$ and, for every form $F$ on $\sU_1$ and every linear functional
$\mathsf{l}$ on the forms on $\sU_2$,}
\[\SquareLaw{\chi_i}{\se}{\varkappa/\rc}{\sx\mapsto\mathsf{l}(\mathsf{u}_i(\sx)F)}{\sy\mapsto\mathsf{l}(\mathsf{v}_i(\sy)F)}.\]
\emph{That is, the two functions have the determinant covariances of Definition~\ref{d:square}, and for all
partners $\dL$, $\dS$ for $\chi_i$ over $\cO_E/\fm_i$ with respect to $\se$:
$\sum_{\sy}\dS(\sy)\,\mathsf{l}(\mathsf{v}_i(\sy)F)=(\varkappa/\rc)\sum_{\sx}\dL(\sx)\,\mathsf{l}(\mathsf{u}_i(\sx)F)$.
The number $\varkappa$ is the same for all $i$ and may be zero.}}
\leanline{\lean{onesquare}}

\statement{arith}{\phantomsection\label{t:3}\textbf{Theorem 3} (the triangle of the data).\enspace \emph{Let $P$ be
parameters, $D$ admissible data for $P$, and $o$ a flag (false or true) that names one of the two grids. Then
there are descents $\mathsf{d}$, linear maps
$\mathsf{d}_{\vartheta,i,n}\colon\TriForm{D}{o}{i}{n+1}\to\TriForm{D}{o}{i+1}{n}$ for all $m$-fold twists
$\vartheta$ of $D$ and natural numbers $i$, $n$, such that}
\[\IsTriangle{m}{\mathrm{ord}_{D,o}}{\CostsSquares{D}{\cdot}{\cdot}}{\triLift{D}{o}{\cdot}{\cdot}{\cdot}{\cdot}}{\mathsf{d}}{i\mapsto\triNegEquiv{D}{o}{i}|_{\triBase{D}{o}}}{o=\mathrm{false}}\]
\emph{holds: indexed by the $m$-fold twists $\vartheta$ of $D$, the spaces $\TriForm{D}{o}{i}{n}$ with the lifts
$\triLift{D}{o}{\vartheta}{i}{n}{\cdot}$, the descents $\mathsf{d}$, the weights given by the triangle orders
$\mathrm{ord}_{D,o}(\vartheta,n)$ of Definition~\ref{d:triforms}\,\textup{(f)}, the grading ``costs $n$ squares''
on the complex functions of $m$-fold twists, and the embeddings of $\triBase{D}{o}$ form a triangle of size $m$ in
the sense of Definition~\ref{d:tri}, whose additional condition is that the flag is false.}}
\leanline{\lean{existsisTriangle}}

\statement{arith}{\phantomsection\label{t:4}\textbf{Theorem 4} (the chain of lifts is injective).\enspace \emph{Let
$P$ be parameters, $D$ admissible data for $P$ and $\vartheta$ an $m$-fold twist of $D$. Then the composite
$\Theta^{\mathrm{lev}}_{D,\vartheta}$ of the $m$ theta lifts from $\U(\sY_0(D))$ to the definite group
$\U(\sY_m(D))$ is injective on the forms of principal level on $\U(\sY_0(D))$.}}
\leanline{\lean{iotaAtLevelinjective}}

\statement{arith}{\phantomsection\label{t:5}\textbf{Theorem 5} (the norm of the lift from $\U(1)$ is an
$L$-value).\enspace \emph{Let $P$ be parameters, $D$ admissible data, $\pp$ a pair of $D$ with character
$\omega_{\pp}$ and point $e_{\pp}$, and $\bn$ a finite index set such that $\pp$ has lift rank $\bn$. Then there is
a complex number $A\ne0$ such that for every tame character $\vartheta$ of $D$}
\[\intAlgPart{\pp,\bn}{\vartheta}=A\cdot L^{\Sigma\cup\bar\Sigma}\bigl(e_{\pp},((\omega_{\pp}\vartheta)^{-1})^{\mathrm{bc}}\bigr),
\qquad\Sigma=\Sigma_D\cup Q(\vartheta),\]
\emph{the value at $e_{\pp}$ of the Hecke $L$-function of the base change of $(\omega_{\pp}\vartheta)^{-1}$ without
the Euler factors at the chosen places $\Sigma_D$, at the places of $Q(\vartheta)$ and at their conjugates.}

\vspace{4pt}\phantomsection\label{t:6}\textbf{Theorem 6} (one measure for all twists of a modulus).\enspace
\emph{Let $P$ be parameters, $D$ admissible data, $\pp$ a pair of $D$ and $\bn$ a finite index set such that $\pp$
has lift rank $\bn$, and let $\vartheta$ be a tame character of $D$. Then
$\IsIntegralMeasure{\vartheta}{\vartheta'\mapsto\intAlgPart{\pp,\bn}{\vartheta'}}$ holds: the function
$\vartheta'\mapsto\intAlgPart{\pp,\bn}{\vartheta'}$ is an integral measure for $\vartheta$
(Definition~\ref{d:intmeas}).}}
\leanline{\lean{intAlgParteqmulL}\sep \lean{isIntegralMeasureintAlgPart}}

\sect{The setting}{s:setting}
\begin{conventions}\phantomsection\label{c:conventions}
The natural numbers are $0,1,2,\dots$. The letters of the handout keep their meaning: $P$ are parameters, $D$
admissible data, $m$ the number of lifts, $\vartheta$ a tame character or a twist of $D$, $\tau$ an infinite
place, $\sigma$ an automorphism of $\CC$, $\ell$ a prime, $\iota$ an isomorphism of $\CC$ with an algebraic closure
of the $\ell$-adic numbers and $j$ a ring homomorphism from $\CC$ to the completion of that closure. Objects local
to the theory of theta series are set in other alphabets, as in the first note: finite index sets are bold
($\ba$, $\bb$, $\bn$, $\bN$, $\bT$, $\bp$, $\br$, \dots) and $|\bn|$ is the number of elements; sans-serif letters
are a commutative ring $\sR$, hermitian spaces, matrices, points of groups and linear maps; fraktur letters are
real data, places, moduli ($\fm$), splittings, auxiliary pairs and $L$-value data; upright letters are constants.
The implicit arguments of the declarations of the unit are of the following kinds, and each statement introduces
the ones it uses: a finite commutative ring; finite index sets, one of them empty or non-empty where a statement
says so; a CM field $E$ with a CM type and, where a statement says so, Weil data; signs and hermitian spaces;
splittings, which are base-adapted where a statement says so; weights, characters, moduli and level places; real
data, finite data and functions; types of forms, forms and kernels; matrices and points; natural numbers and
integers; constants; a prime with an isomorphism or an embedding of $\CC$; inner product spaces, with non-zero
weights where a statement says so; parameters and admissible data with their tame characters, twists and pairs.
\end{conventions}

\subhead{What is recalled from the first note.} Hermitian spaces $\sW$ over a CM field $E$ with index set $\bn$,
signs, Gram matrix and adelic frame; $\negc(\sW)$, the number of negative signs; the Weil representation of a
pair, real data (rapid when all weighted norms are finite), finite data, the theta kernel
$\theta_{\sW,\sV}(\omega,\omega';f)$ and the theta integral $\thI{\sW}{\sV}{\omega}{\omega'}{f}$ over $[\U(\sV)]$,
for the measure fixed in the first note (the Haar measure that gives the stabilizer of the integral lattice the
mass $1$ at the finite places, and the measure of the formal text at the infinite places); Weil data, which
include the hypothesis that $E$ contains a primitive root of unity of order $5040$;
split tame characters $\fc$ with character $\omega_{\fc}$, places $\Sigma(\fc)$ and local character $\chi_{\fc}$;
the theta integrals $\thL{\sW}{\sV}{\fc}{\fv}{\dL}$ and $\thS{\sW}{\sV'}{\fc}{\fv}{\dS}$ on the large and on the
small side; for an identification $\se\colon\bb\sqcup\ba\simeq\bn$ and $\sk\in\GL_{\bn}(\sR)$, the matrices
$\colk{\se}{\sk}$ (columns of $\sk$ indexed by $\bb$) and $\rowk{\se}{\sk}$ (rows of $\sk^{-1}$ indexed by $\ba$);
partners for $\chi$ over a subring; the standard data $\str{\se}{\chi}{\sz}$, $\stc{\chi}{\sz}$; splittings and
real slices; the products $\LP{\fc}{e}$; the equal-rank ratio $\eqr{\sV}{\fc}{\fv}$; the standard exchange identity
$\Exch{\sX}{\sV}{\sV'}{\mu}{\se}{\fv}{\fy}{\mathcal{N}}$; determinant vectors; and $L$-value data with the lift of the trivial character.

\subhead{Moduli, lifts and descents.} A \kw{modulus} $\fm$ of $E$ is a product of powers of finite
places of which none is the conjugate of another; $\Sigma(\fm)$ is its set of places and $\cO_E/\fm$ its residue
ring. A character $\omega$ of $\U(1)(\A)$ is \kw{tame for} $\fm$ \kw{with weights} $\mu$ \kw{and character}
$\chi$ if $\fm$ has a place, $\omega$ has weights $\mu$, its tame part is non-trivial exactly at the places of
$\fm$, and $\chi$ is the character of $\cO_E/\fm$ made from the local tame parts. For a function $\dL$ on matrices
over $\cO_E/\fm$ and a real datum $\fx$, $\md{\fx}{\dL}{\fm}$ is the datum with real part $\fx$, the function
$\dL$ of the reduction modulo $\fm$ at the places of $\fm$, and the vacuum elsewhere. A \kw{type} of forms on a
hermitian space is a discrete series type at the infinite places; the \kw{theta lift} and the \kw{descent}
attached to a pair of spaces, characters, a real datum and a function modulo $\fm$ are the integrals of a form
against the theta kernel of the pair, from the first space to the second and back. An \kw{edge family}, for a
family of characters indexed by $\cI$, two types and a direction, is a real datum that lifts forms of one type
to the other in both directions, uniformly in the family; a real datum that does so for one choice of characters
is an \kw{edge}, and with a finite datum it gives an \kw{arrow}.

\subhead{The data.} For parameters $P$ and admissible data $D$: $E(D)$ is the CM field, $\Sigma_D$ the set of the
chosen places (above $g(D)$), $f(D)$ the number of infinite places. A \kw{tame character} $\vartheta$ of $D$ has a
finite set $Q(\vartheta)$ of places, its modulus; $\Np Q(\vartheta)$ is the product of their norms; $\vartheta$
is a \kw{twist} if its tame part is non-trivial at every place of $Q(\vartheta)$ and an \kw{$m$-fold twist} if
this holds for its first $m$ powers. $\fm_D(\vartheta,e)$ is the modulus that is the product of the chosen places
to the power $e$ times $Q(\vartheta)$. A \kw{pair} $\pp$ of $D$ consists of a character $\omega_{\pp}$ of
$\U(1)(\A)$ of weights $-\mu_{\pp}(\tau)$, where the $\mu_{\pp}(\tau)$ are natural numbers, whose tame part is
trivial off the chosen places, and of a natural number $e_{\pp}$, the \kw{point}, with $1\le e_{\pp}\le t(D)$ and
$e_{\pp}\le\mu_{\pp}(\tau)\le4b$ for every infinite place $\tau$ ($b$ the third scale of the data). Families are
indexed by the pairs $(\ell,\iota)$ of a prime $\ell$ and an isomorphism $\iota$ of $\CC$ with an algebraic
closure of the $\ell$-adic numbers, all primes at once: a family $\theta=(\theta_{\ell,\iota})$ of $m$-fold twists
may be \kw{fit} (Chapter~VIII), and a family of complex numbers $(y_{\ell,\iota})$ may have \kw{height at most} a
real number (Chapter~IV; the bound couples the primes, so the height of a family indexed by the $\iota$ of one
prime is not defined). The spaces $\sY_i(D)$, $0\le i\le m$, are
those of the chain of lifts, $\sY_m(D)$ being definite; the two \kw{grids} of spaces, named by a flag $o$, with
their index sets, parameters, characters, harmonics and moduli, are defined in the third note.

\statement{inksoft}{\textbf{Inputs.} The proofs below take the following for granted. Numbers are the units of the
reading; references are those of the handout's catalogue, and none of their places is confirmed.

\vspace{3pt}\textbf{Input 1} (the first note of this unit).\enspace The exchange identity at the top of the tower
(its Theorem~1); the equivariant lift of the difference of the two theta integrals; a
stable space that contains a multiple of a standard datum contains every locally constant partner; partners from a function of a block; a datum with prescribed slices; products of
data; the constancy of the equal-rank theta integral; the cone integral is onto the determinant vectors; the norm
of the lift of the trivial character (its Theorem~5).

\vspace{3pt}\textbf{Input 2} (the third note of this unit).\enspace The parameters of the grid correspond along
its edges; a character of the grid is tame for the modulus of the triangle; the index sets, signs, parameters and
harmonics of the grids.

\vspace{3pt}\textbf{Input 3} (units 131 and 98, the tower and Landherr's theorem).\enspace The \emph{tower
property}: the constant term of a theta integral on the large side along an isotropic line is a multiple
$\rc_{\mathrm{tow}}$ of the theta integral of the smaller pair with the cone integral of the real datum; the
complex number $\rc_{\mathrm{tow}}\ne0$ depends on the split tame character $\fc$, and
$|\rc_{\mathrm{tow}}|^2\,\LP{\fc}{e}=\rC_{\mathrm{tow}}\,\LP{\fc}{e+1}$ with a real number $\rC_{\mathrm{tow}}$ that does
not depend on $\fc$ (after S.~Rallis, Compositio Math.\ 51 (1984); an analogue);
\emph{Landherr's theorem}: two spaces with the same signs differ by a rational change of coordinates and a move of
the adelic point (W.~Landherr, Abh.\ Math.\ Sem.\ Hamburg 11 (1936); M.~Kneser, TIFR lectures 47 (1969), §5.5; a
special case); $\LP{\fc}{e}>0$.}

\statement{inksoft}{\textbf{Inputs (continued).}

\vspace{3pt}\textbf{Input 4} (units 26, 27 and 25, inner products of lifts).\enspace The inner product formula for theta
lifts (the \emph{Rallis inner product formula}) in its Euler form, for an edge and a sharp finite datum, with
injectivity of the lift, as stated in Input~6 of the first note (relation to the literature unresolved:
S.~Rallis, Compositio Math.\ 52 (1984), §1 and Thm~1.6, treats another dual pair and is an analogue only; the
unitary inner product in this rank range is J.-S.~Li, J.~reine angew.\ Math.\ 428 (1992), §2, Thm~2.1, eq.~(26),
p.~184; the constant uniform in the characters and the sharp data, the evaluation as a sum of coset terms and the
injectivity go beyond both); the Petersson product of two descents; an edge
gives arrows; the adjoint of a lift along an arrow is the conjugate of a descent; the theta kernel of the
opposite space is a multiple of the conjugate kernel; the decay of the theta kernel in the height.

\vspace{3pt}\textbf{Input 5} (unit 103, Gauss sums).\enspace The \emph{matrix Gauss sum} for primitive characters of
a finite ring, as stated in Input~1 of the first note (relation to the literature unresolved: T.~Kondo, J.~Math.\
Soc.\ Japan 15 (1963), pp.~244--245, has the law over a finite field at invertible frequencies; the statement
over a finite ring and at all frequencies is wider than the source); Jacobi's identity for complementary minors
(unit 112). A multiplicative character $\chi$ of a finite commutative ring $\sR$ is \kw{primitive} if for every
$\mathsf{a}\ne0$ in $\sR$ there is $\mathsf{c}$ with $\chi(1+\mathsf{c}\mathsf{a})\notin\{0,1\}$, and an additive
character is \kw{primitive} if it is non-trivial on every non-zero ideal (first note).

\vspace{3pt}\textbf{Input 6} (units 96 and 104, Hecke characters).\enspace Partial Hecke $L$-values of characters of
the weights that occur do not vanish at the points that occur: for a character $\omega$ of $\U(1)(\A)$ with
non-zero weights and a finite set $\Sigma$, $L^{\Sigma\cup\bar\Sigma}(s,\omega^{\mathrm{bc}})\ne0$ for real $s\ge1$
(D.~Ramakrishnan and R.~Valenza, \emph{Fourier Analysis on Number Fields}, 1999, Thm~7-28; J.~Tate, 1950, in
Cassels--Fröhlich, ch.~XV, §4.4; a special case: the real axis, with finitely many Euler factors removed); a
character of given weights is tame for a modulus
that contains the places where its tame part is non-trivial; the residue ring of $\fm_D(\vartheta,e)$.

\vspace{3pt}\textbf{Input 7} (units 134 Forms, 138 Selmer, 141 Seesaw, 150 Curve, 133 Discrete series, 152
Data).\enspace Auxiliary pairs and their data; standard data and admissible combinations of translates of the
integral vacuum, with the divisibility of their jets; the conjugation of jets by an automorphism of $\CC$; the
Harish-Chandra parameter of a weight by steps; fit families of twists and the height of their ratios of
$L$-values (Chapter~VIII); the served primes of the data (unit 150); the column and row index sets of the data.

\vspace{3pt}\textbf{Input 8} (unit 135 Heights, Chapter~IV; unit 25, lifts).\enspace Heights of families indexed
by the pairs $(\ell,\iota)$: heights add in products, and a bound for the height may be enlarged. \emph{Edge
families from corresponding parameters}: when the parameters of two spaces correspond, the joint harmonic is an
edge, uniformly in a family of characters; this is a result of the background library and not of the argument
(relation to the literature unresolved: A.~Ichino, arXiv:2002.09148, Thm~4.1 and Lemma~7.8, gives the criterion
for the joint types; the formal statement selects witnesses for a two-sided analytic statement that goes beyond
it).}

\sect{The square law over a finite ring}{s:square}
\begin{definition}[the square law and its data]\label{d:square}
Let $\sR$ be a finite commutative ring and $\ba$, $\bb$, $\bn$ finite index sets.
\begin{enumerate}[label=\textup{(\alph*)},leftmargin=2.3em,itemsep=1pt]
\item Let $\chi$ be a multiplicative character of $\sR$, $\se\colon\bb\sqcup\ba\simeq\bn$ an identification,
$\rc\in\CC$, and $\cL$, $\cL'$ complex functions on $\Mat{\ba}{\bn}(\sR)$ and on $\Mat{\bn}{\bb}(\sR)$. The
statement $\SquareLaw{\chi}{\se}{\rc}{\cL}{\cL'}$ holds if $\cL(\sk\sx)=\chi(\det\sk)\cL(\sx)$ for
$\sk\in\GL_{\ba}(\sR)$, $\cL'(\sy\sk)=\chi(\det\sk)\cL'(\sy)$ for $\sk\in\GL_{\bb}(\sR)$, and
\[\sum_{\sy}\dS(\sy)\,\cL'(\sy)=\rc\sum_{\sx}\dL(\sx)\,\cL(\sx)\]
for all partners $\dL$, $\dS$ for $\chi$ over $\sR$ with respect to $\se$.
\item For a commutative ring $\sR$ and an $(\ba\sqcup\bb)\times\ba$ matrix $\sx$ over it,
$\inclusionDatum{\sR}{\sx}$ is $1$ if $\sx$ is the inclusion matrix $\binom{1}{0}$ and $0$ otherwise.
\item For a complex function $G$ on $\Mat{\bn}{\ba}(\sR)$ and an $(\bn\sqcup\bb)\times(\ba\sqcup\bb)$ matrix $\sy$,
$\nextDatum{G}{\sy}$ is $G(-\sy_{11})$ divided by the number of $\bn\times\bb$ matrices over $\sR$ if $\sy$ has the
block form $\bigl(\begin{smallmatrix}\sy_{11}&*\\0&1\end{smallmatrix}\bigr)$, and $0$ otherwise.
\item The \kw{swap} is the bijection $\ba\sqcup(\bn\sqcup\bb)\simeq\bn\sqcup(\ba\sqcup\bb)$ that exchanges the
first two summands.
\item For $\se\colon\bb\sqcup\ba\simeq\bn$, $\blockInlGL{\se}$ is the homomorphism
$\sg\mapsto\operatorname{diag}(\sg,1)$ from $\GL_{\bb}(\sR)$ to $\GL_{\bn}(\sR)$.
\item For $\chi$, $\se$ as in \textup{(a)} and matrices $\sx\in\Mat{\ba}{\bn}(\sR)$, $\sy\in\Mat{\bn}{\bb}(\sR)$,
the \kw{graph kernel} $\graphKernel{\chi}{\se}{\sx}{\sy}$ is $\chi(\det\sk)/\#\GL_{\bb}(\sR)$ if
$(\sx,\sy)=(\rowk{\se}{\sk},\colk{\se}{\sk})$ for some $\sk\in\GL_{\bn}(\sR)$, and $0$ otherwise.
\end{enumerate}
\end{definition}
\leanline{\lean{SquareLaw}\sep \lean{inclusionDatum}\sep \lean{nextDatum}\sep \lean{swapFirstTwo}\sep
\lean{blockInlGL}\sep \lean{graphKernel}}

\begin{lemma}[elementary facts]\label{l:elem}
Let $\sR$ be a finite commutative ring (for \textup{(c)}, \textup{(d)} any commutative ring; for \textup{(a)},
\textup{(b)} any finite set) and $\ba$, $\bb$, $\bN$, $\bn$ finite index sets.
\begin{enumerate}[label=\textup{(\alph*)},leftmargin=2.3em,itemsep=1pt]
\item For a complex function $G$ on $(\bN\sqcup\bn)\times\ba$ matrices, the sum of $G$ over all matrices is the
double sum over the two row blocks.
\item For a complex function $G$ on $\ba\times(\bN\sqcup\bn)$ matrices, the sum of $G$ over all matrices is the
double sum over the two column blocks.
\item For $\sff\in\GL_{\ba}(\sR)$ and $\sg\in\GL_{\bb}(\sR)$ the map $\sx\mapsto\sff^{-1}\sx\,\sg$ on $\ba\times\bb$
matrices is injective.
\item The inclusion matrix of $\ba$ into $\ba\sqcup\bb$ is the block matrix $\binom10$.
\item For a multiplicative character $\chi$ of $\sR$, an identification $\se\colon\bb\sqcup\ba\simeq\bn$ and
$\sk\in\GL_{\bn}(\sR)$: $\graphKernel{\chi}{\se}{\rowk{\se}{\sk}}{\colk{\se}{\sk}}=\chi(\det\sk)/\#\GL_{\bb}(\sR)$; in
particular the value in Definition~\ref{d:square}\,\textup{(f)} does not depend on the choice of $\sk$.
\item For $\chi$, $\se$ as in \textup{(e)} and matrices $\sx$, $\sy$ with $\graphKernel{\chi}{\se}{\sx}{\sy}\ne0$
there is $\sk\in\GL_{\bn}(\sR)$ with $(\sx,\sy)=(\rowk{\se}{\sk},\colk{\se}{\sk})$.
\item For $\chi$, $\se$ as in \textup{(e)}, matrices $\sx$, $\sy$ and $\sg\in\GL_{\bb}(\sR)$, the graph kernel
satisfies $\graphKernel{\chi}{\se}{\sx}{\sy\sg}=\chi(\det\sg)\,\graphKernel{\chi}{\se}{\sx}{\sy}$.
\end{enumerate}
\end{lemma}
\begin{proof}
(a)--(d) are reindexings and definitions. (e): two elements of $\GL_{\bn}(\sR)$ with the same blocks
$\rowk{\se}{\sk}$, $\colk{\se}{\sk}$ have the same determinant, by Jacobi's identity (Input~5). (f) is the
definition. (g): if $(\sx,\sy)$ are the blocks of $\sk$, then $(\sx,\sy\sg)$ are the blocks of
$\sk\cdot\blockInlGL{\se}(\sg)$, of determinant $\det\sk\det\sg$.
\end{proof}
\leanline{\lean{sumeqsumfromRows}\sep \lean{sumeqsumfromCols}\sep \lean{injectiveglinvmulmul}\sep
\lean{inclMatrixinl}\sep \lean{graphKernelinvLastRowsfirstCols}\sep \lean{existseqofgraphKernelnezero}\sep
\lean{graphKernelmulright}}

\begin{lemma}[the square law at the inclusion]\label{l:next}
Let $\sR$ be a finite commutative ring, $\ba$, $\bb$, $\bn$ finite index sets, $\chi$ a multiplicative character
of $\sR$, $\rc\in\CC$, and $\cL$, $\cL'$ functions with $\SquareLaw{\chi}{\mathrm{swap}}{\rc}{\cL}{\cL'}$ for the
swap $\ba\sqcup(\bn\sqcup\bb)\simeq\bn\sqcup(\ba\sqcup\bb)$. Then for every complex function $G$ on
$\Mat{\bn}{\ba}(\sR)$, with $\sx_0$ the inclusion matrix $\binom10$,
\[\#\GL_{\ba}(\sR)\sum_{\sy}G(\sy)\,\cL'\tbinom{\sy}{\sx_0}=\rc\cdot\#\GL_{\bn\sqcup\bb}(\sR)\sum_{\sx}\nextDatum{G}{\sx}\,\cL(\sx_0\mid\sx).\]
\end{lemma}
\begin{proof}
Apply the partner identity of the square law to the block partners of the first note (Input~1) made from the
point function of $\binom{\sx}{0}$, and move the blocks to the standard ones with the two equivariances of $\cL$
and $\cL'$; the orbits are counted by the orders of the two groups, Lemma~\ref{l:elem}\,(a)--(d) reindexing the
sums.
\end{proof}
\leanline{\lean{sumeqnextDatum}}

\begin{lemma}[equivariance of the graph kernel]\label{l:greq}
Let $\sR$ be a finite commutative ring, $\ba$, $\bb$, $\bn$ finite index sets, $\chi$ a multiplicative character
of $\sR$ and $\se\colon\bb\sqcup\ba\simeq\bn$.
\begin{enumerate}[label=\textup{(\alph*)},leftmargin=2.3em,itemsep=1pt]
\item For $\sg\in\GL_{\bn}(\sR)$, $\sff\in\GL_{\ba}(\sR)$, $\sq\in\GL_{\bb}(\sR)$ and matrices $\sx$, $\sy$:
$\graphKernel{\chi}{\se}{\sx}{\sy}=\chi(\det\sg)\chi(\det\sff)^{-1}\chi(\det\sq)^{-1}\,\graphKernel{\chi}{\se}{\sff^{-1}\sx\sg}{\sg^{-1}\sy\sq}$.
\item Let $\dS$ be a complex function on $\Mat{\bn}{\bb}(\sR)$ with $\dS(\sy\sg)=\chi^{-1}(\det\sg)\dS(\sy)$ for
all $\sy$ and $\sg\in\GL_{\bb}(\sR)$. Then for every $\sk\in\GL_{\bn}(\sR)$:
$\sum_{\sy}\graphKernel{\chi}{\se}{\rowk{\se}{\sk}}{\sy}\,\dS(\sy)=\chi(\det\sk)\,\dS(\colk{\se}{\sk})$.
\end{enumerate}
\end{lemma}
\begin{proof}
(a) We may assume the kernel is non-zero at one of the two pairs; if $(\sx,\sy)$ are the blocks of $\sk$, then
$(\sff^{-1}\sx\sg,\sg^{-1}\sy\sq)$ are the blocks of $\sg^{-1}\sk\operatorname{diag}(\sq,\sff)$. (b) The non-zero
terms are the $\sy=\colk{\se}{\sk}\sg$, $\sg\in\GL_{\bb}(\sR)$, and each contributes
$\chi(\det\sk)\dS(\colk{\se}{\sk})/\#\GL_{\bb}(\sR)$ by Lemma~\ref{l:elem}\,(g).
\end{proof}
\leanline{\lean{graphKernelequivariant}\sep \lean{sumgraphKernelmul}}

\begin{lemma}[the graph kernel reproduces the small side]\label{l:pk}
Let $\sR$ be a finite commutative ring, $\ba$, $\bb$, $\bN$, $\bn$ finite index sets, $\chi$ a multiplicative
character of $\sR$, $\se\colon\bb\sqcup\ba\simeq\bN\sqcup\bn$ an identification, and suppose that $(\sR,\sR)$ has
the elementary divisor property (is a valuation pair). There is a function $\mathrm{d}_0$ on
$(\bN\sqcup\bn)\times\bb$ matrices such that:
\begin{enumerate}[label=\textup{(\arabic*)},leftmargin=2.3em,itemsep=1pt]
\item if $\mathrm{d}_0(\sy)\,\graphKernel{\chi}{\se}{\sx}{\sy}\ne0$, then this number is a root of unity, a left
inverse of the first row block of $\sy$ gives one of the second column block of $\sx$, and a left inverse of the
second row block of $\sy$ gives one of the first column block of $\sx$;
\item the average over $\sk\in\GL_{\bb}(\sR)$ of $\mathrm{d}_0(\sy\sk)\,\graphKernel{\chi}{\se}{\sx}{\sy}$ is
$\graphKernel{\chi}{\se}{\sx}{\sy}$;
\item for all $\rc$, $\cL$, $\cL'$ with $\SquareLaw{\chi}{\se}{\rc}{\cL}{\cL'}$ and every $\sy$ one of whose row
blocks has a left inverse: $\cL'(\sy)=\rc\sum_{\sx}\graphKernel{\chi}{\se}{\sx}{\sy}\,\cL(\sx)$.
\end{enumerate}
\end{lemma}
\begin{proof}
Let $\mathrm{d}_0$ be $\#\GL_{\bb}(\sR)$ at one chosen point of each orbit $\sy\GL_{\bb}(\sR)$ and $0$ elsewhere;
(2) follows. The values of the kernel are values of $\chi$ divided by the order of the group, whence the roots of
unity in (1), and the statements on left inverses are read off the blocks of an invertible matrix. For (3),
write $\sy=\colk{\se}{\mathsf{j}}$ with $\mathsf{j}$ invertible (this uses the left inverse and the elementary
divisor property), and apply the partner identity of the square law to the partners $(\dL,\dS)$ with
\[\dS=\graphKernel{\chi^{-1}}{\se}{\rowk{\se}{\mathsf{j}}}{\cdot},\qquad
\dL(\sx)=\sum_{\sy'}\graphKernel{\chi}{\se}{\sx}{\sy'}\,\dS(\sy');\]
they are partners, and the two sums are evaluated by Lemma~\ref{l:greq} and Lemma~\ref{l:elem}.
\end{proof}
\leanline{\lean{existspartnerkernel}}

\begin{lemma}[the square law after a Fourier transform]\label{l:fourier}
Let $\sR$ be a finite commutative ring, $\bN$, $\bT$ finite index sets, $\chi$ a multiplicative character of
$\sR$ that is primitive and $\mathrm{e}_{\sR}$ an additive character of $\sR$ that is primitive. There is a complex
number $\rc_{\mathrm{G}}$ with $|\rc_{\mathrm{G}}|^2\cdot(\#\sR)^{|\bT|(|\bT|+2|\bN|)}=1$ such that for all $\rc$, $\cL$, $\cL'$ and
every function $G$ on $\Mat{\bN}{\bT}(\sR)$: if $\SquareLaw{\chi}{\mathrm{swap}}{\rc}{\cL}{\cL'}$ holds for the swap
with $\ba=\bN\sqcup\bT$, $\bb=\bT$, $\bn=\bN$, then, with $\sx_0$ the inclusion matrix and $\widehat{\ }$ the
Fourier transform for $\mathrm{e}_{\sR}$,
\[\sum_{\sy}\widehat{\overline{G}}(\sy)\,\cL'\tbinom{\sy}{\sx_0}=\rc_{\mathrm{G}}\,\rc\sum_{\sx}\widehat{\overline{\nextDatum{G}{\cdot}}}(\sx)\,\cL(\sx_0\mid\sx).\]
\end{lemma}
\begin{proof}
Apply Lemma~\ref{l:next} to the Fourier transform of $\overline{G}$ and compute the Fourier transform of the next
datum with the matrix Gauss sum (Input~5); $\rc_{\mathrm{G}}$ is the conjugate Gauss sum divided by the number of
matrices.
\end{proof}
\leanline{\lean{sumfouriereqnextDatum}}

\sect{From the top of the tower to the exchange theorem}{s:exchange}
\begin{lemma}[the opposite space]\label{l:neg}
Let $E$ be a CM field with a CM type, $\bn$ a finite index set, $\sW$ a hermitian space with index set $\bn$ and
$\sS$ a hermitian space with the same index set and the opposite signs, and suppose that the identification
$\sg\mapsto\sg^{-}$ of the adelic unitary groups of opposite signs carries the rational points of $\sW$ onto those
of $\sS$: $\sg^{-}$ is rational for $\sS$ if and only if $\sg$ is rational for $\sW$. Then:
\textup{(1)} for every parameter $l$ and every function $F$ on $\U(\sS)(\A)$, the function $\sg\mapsto F(\sg^{-})$ is
a form on $\sW$ of the type of $l$ if and only if $F$ is a form on $\sS$ of the type of $l$ for the opposite
signs; \textup{(2)} for every complex function $F$ on $\U(\sS)(\A)$, the integral of $F$ over $[\U(\sS)]$ equals
the integral of $\sg\mapsto F(\sg^{-})$ over $[\U(\sW)]$.
\end{lemma}
\begin{proof}
The identification is an isomorphism of topological groups that preserves Haar measures and maximal compact
subgroups; unfolding the definitions of forms and of quotient measures, both statements are transport along it.
\end{proof}
\leanline{\lean{isFormnegiffandintegralneg}}

\begin{lemma}[moving the spaces]\label{l:move}
Let $E$ be a CM field with a CM type and Weil data; let $\sW$, $\sW_1$ be hermitian spaces with index set $\bn$
and the same signs, and $\sV$, $\sV_1$ hermitian spaces with index set $\ba$ and the same signs. Let
$\mathsf{c}\in\GL_{\bn}(E)$ and let $\sk$ be an adelic point of the unitary group of the signs with
$\mathrm{x}_{\sW}=\mathsf{c}^{-1}\,\mathrm{x}_{\sW_1}\sk$ for the adelic frames, and let $\fc$ be a split tame
character. Then there is a complex number $\mathrm{u}$ of absolute value $1$ such that for all real data $\fv$,
functions $\dL$ and points $\sg$: $\thL{\sW}{\sV}{\fc}{\fv}{\dL}(\sg)=\mathrm{u}\cdot\thL{\sW_1}{\sV_1}{\fc}{\fv}{\dL}(\sk\sg)$.
\end{lemma}
\begin{proof}
By Landherr's theorem (Input~3) $\sV_1$ is moved to $\sV$ by a rational change of coordinates and an adelic
point. The two theta integrals are then integrals over the same orbit; the theta functionals are compared on the
vacuum and on products of theta functions of lines, and the constant is the value of a character at a quotient of
determinants.
\end{proof}
\leanline{\lean{thetaLargemove}}

\begin{lemma}[adding hyperbolic planes; slices at zero]\label{l:planes}
Let $E$ be a CM field with a CM type and Weil data.
\begin{enumerate}[label=\textup{(\alph*)},leftmargin=2.3em,itemsep=1pt]
\item Let $\bn$, $\br$, $\bp$ be finite index sets, $\sW$ a hermitian space with index set $\bn$ and
$\mathsf{t}$ signs on $\br$ with $|\br|=|\bn|+2|\bp|$ and with $|\bp|+\negc(\sW)$ negative signs. Then there are a
hermitian space $\sW'$ with index set $\br$ and signs $\mathsf{t}$ and a base-adapted splitting of $\sW'$ over
$\sW$ along $\bp$.
\item Let $\sW$, $\sW'$, $\sU$ be hermitian spaces, $\fS$ a base-adapted splitting of $\sW'$ over $\sW$ along
$\bp$, and $\fS'$ a base-adapted splitting of $\sW'$ over $\sW$ along $\br$ whose isotropic frame is that of $\fS$
reindexed by a bijection and whose embedding of $\sW$ is that of $\fS$. Then the real slices at $0$ for the
target $\sU$ agree: $\mathrm{sl}_{\fS}(0;\cdot)=\mathrm{sl}_{\fS'}(0;\cdot)$.
\item Let $\sW$, $\sW'$, $\sU$, $\sV'$ be hermitian spaces, $\fS$ a base-adapted splitting of $\sW'$ over $\sW$
along $\bp$ and $\fS''$ a base-adapted splitting of $\sW$ over $\sV'$ along $\br$. For every real datum $\fy$ of
the pair $(\sW',\sU)$: the slice at $0$ for the composite splitting is
$\mathrm{sl}_{\fS''}(0;\mathrm{sl}_{\fS}(0;\fy))$.
\item Let $\bn$ be an empty finite index set, $\sW$ a hermitian space with index set $\bn$, $\mu$ weights, $\fc$ a
split tame character of weight $\mu$ and $\dL$ a complex function on the (empty) matrices. Then the vector named
by $\dL$ at the places of $\fc$ is $\dL(0)$ times the vacuum.
\end{enumerate}
\end{lemma}
\begin{proof}
(a) Reindex the sum of $\sW$ and $|\bp|$ hyperbolic planes by a bijection that matches the signs; its standard
splitting is base-adapted. (b) holds by definition. (c) Compare integrated zero terms: the zero term of a
composite splitting is the composite of the zero terms. (d) On empty matrices $\dL$ is $\dL(0)$ times the
indicator of the integral box, and naming by that indicator is a projection that fixes the vacuum.
\end{proof}
\leanline{\lean{existsisBaseAdapted}\sep \lean{slicezerocongr}\sep \lean{slicezerocomp}\sep
\lean{namevacElemofisEmpty}}

\begin{proposition}[descent: removing hyperbolic planes]\label{p:descent}
Let $E$ be a CM field with a CM type and Weil data; let $\bn$, $\bn'$, $\br$, $\br'$, $\bp$, $\bb$ be finite index
sets with $\bp$ non-empty; let $\sW$ (index set $\bn$), $\sW'$ (index set $\bn'$), $\sU$ (index set $\bb$),
$\sV'$ (index set $\br'$) be hermitian spaces and $\fS$ a base-adapted splitting of $\sW'$ over $\sW$ along $\bp$.
Let $\mu$ be weights, $\mathcal{N}$ a real function on the split tame characters of weight $\mu$, $e$ a natural
number and $i(\tau)$ natural numbers at the infinite places $\tau$, $\sV$ a hermitian space with index set $\br$
such that $\negc(\sV')=|\bp|+\negc(\sV)$, and
$\se$, $\se'$ identifications for the standard row data of $(\sW,\sV)$ and $(\sW',\sV')$; suppose that $|\bn|$ is
even. Then there is a map $\mathsf{f}$ from the real data of $(\sW',\sV')$ to those of $(\sW,\sV)$ such that
\begin{enumerate}[label=\textup{(\arabic*)},leftmargin=2.3em,itemsep=1pt]
\item for all real data $\fx$, $\fy$: if
\[\Exch{\sW'}{\sV'}{\sU}{\mu}{\se'}{\fx}{\fy}{\mathcal{N}}\]
holds, then so does
\[\Exch{\sW}{\sV}{\sU}{\mu}{\se}{\mathsf{f}(\fx)}{\mathrm{sl}_{\fS}(0;\fy)}{\mathcal{N}};\]
\item if $|\bn|=2(e+\negc(\sW))$, $|\bn'|=2(e+\negc(\sW'))$ and $\mu(\tau)=e+i(\tau)$ for all $\tau$, then every
determinant vector of $(\sW,\sV)$ with the parameters $e$ and $(i(\tau))_\tau$ (first note: the natural number and
the family of natural numbers at the infinite places) is the image under $\mathsf{f}$ of a determinant vector of
$(\sW',\sV')$ with the same parameters.
\end{enumerate}
\end{proposition}
\begin{proof}
Strong induction on $|\bp|$. \emph{One plane.} Let $\mathsf{f}$ be the cone integral of the tower property
(Input~3). The constant term along the isotropic line of the large theta integral of $(\sW',\sV')$ is
$\rc_{\mathrm{tow}}$ times the large theta integral of $(\sW,\sV)$ with the datum $\mathsf{f}(\fx)$, and
$|\rc_{\mathrm{tow}}|^2\LP{\fc}{e'}=\rC_{\mathrm{tow}}\LP{\fc}{e'+1}$ for $e'=|\br|-|\bn|/2$: this shifts the argument
of the product of $L$-values by one, which is the shift between the two pairs in the definition of the standard
exchange identity. On the small side the constant term of the theta integral of $(\sW',\sU)$ along the line is
the theta integral of $(\sW,\sU)$ with the slice at $0$, only the zero term surviving, times the vacuum factor
$\operatorname{vac}_{\fS}(\sU)$ of the mixed model and a character value. Dividing: if $\rc$ and $\rC$ are the
constants of the identity for $(\sW',\sV',\sU)$, the constant for $(\sW,\sV,\sU)$ is
$\rc\cdot\operatorname{vac}_{\fS}(\sU)\cdot\mathrm{u}/\rc_{\mathrm{tow}}$ with $|\mathrm{u}|=1$ (Lemma~\ref{l:move}). It
depends on $\fc$ through $\rc_{\mathrm{tow}}$; but its squared absolute value is governed by $\rC_{\mathrm{tow}}$,
which is free of $\fc$, and the real constant $\rC\,\rC_{\mathrm{tow}}/|\operatorname{vac}_{\fS}(\sU)|^2$ serves for
$(\sW,\sV,\sU)$. The standard data of an integral matrix restrict to standard data. Assertion (2) for one plane is the surjectivity
of the cone integral on determinant vectors (Input~1). \emph{Several planes.} Remove one plane, then the others,
and compose: Lemma~\ref{l:planes}\,(b), (c) identify the slices at $0$ of the composite, and Lemma~\ref{l:planes}\,(a)
provides the intermediate space.
\end{proof}
\leanline{\lean{existsstdExchangeremovingplanes}}

\begin{proposition}[the exchange identity with an implicit constant]\label{p:partners}
Let $E$ be a CM field with a CM type and Weil data; let $\bn$, $\bm$, $\br$, $\bp$, $\bb$ be finite index sets with
$\bp$ non-empty; let $\sW$ (index set $\bn$), $\sX$ (index set $\bm$), $\sU$ (index set $\bb$) be hermitian spaces;
let $\mu$ be weights, $\sX^{+}$ the sum of $\sX$ and one hyperbolic plane, $\fS$ a base-adapted splitting of
$\sX^{+}$ over $\sW$ along $\bp$, $\sV$ a hermitian space with index set $\br$, and $\se\colon\bb\sqcup\br\simeq\bn$
an identification. Let $\mathfrak{z}$ and $\fy$ be rapid real data, for $(\sX^{+},\sX^{+})$ and for $(\sU,\sW)$.
Suppose $1+\negc(\sX)=|\bp|+\negc(\sV)+\negc(\sU)$ and that $|\bn|$ is even. Then there are a real number $\rC$
and a rapid real datum $\fv$ for $(\sW,\sV)$ such that every split tame character $\fc$ of weight $\mu$, with
$\omega^{\circ}=\omega_{\fc}^{-1}\upsilon^{|\bn|/2}$, has a complex number $\rc$ with
\[\rC\,|\rc|^2=\LP{\fc}{|\br|-|\bn|/2}\cdot\eqr{\sX^{+}}{\fc}{\mathfrak{z}}\]
and, for all partners $\dL$, $\dS$ for $\chi_{\fc}$ over the integers of $\sF_{\fc}$ with respect to $\se$ with
$\dS$ locally constant, and every point $\sg$:
\[\thI{\sW}{\sV}{\omega^{\circ}}{1}{\fv\otimes[\dL]}(\sg)=\rc\int_{[\U(\sU)]}\theta_{\sU,\sW}\bigl(\omega^{\circ},\omega^{\circ};\fy\otimes[\dS]\bigr)(\sh,\sg)\,d\sh,\]
where $\sF_{\fc}$ is the product of the completions of $E$ at the places of $\fc$ and $[\,\cdot\,]$ the local
datum at these places.
\end{proposition}
\begin{proof}
Exchange the two members of the pair $(\sU,\sW)$: by Lemma~\ref{l:neg} the right-hand side is a theta integral on
the small side for $\sW$ and the space opposite to $\sU$. Build a space at the top of the tower over $\sW$ with
Lemma~\ref{l:planes}\,(a), and matched real data there: a datum with the prescribed slice $\fy$ at $0$ (Input~1),
and then the exchange identity at the top (Input~1), whose normalising function is the equal-rank ratio of the
datum $\mathfrak{z}$. Descend with Proposition~\ref{p:descent},
removing the planes of $\fS$: this gives the standard exchange identity for $(\sW,\sV,\sU)$, that is, the
identity for the standard data of all integral matrices, with $\rC|\rc|^2$ as stated. To pass to all partners,
use the equivariant lift of the difference of the two integrals (Input~1): its kernel is a stable space of
Schwartz--Bruhat functions that contains the standard data, so it contains every locally constant partner
(Input~1).
\end{proof}
\leanline{\lean{exchangeofpartners}}

\begin{proposition}[the equal-rank constant does not depend on the character]\label{p:equal}
Let $E$ be a CM field with a CM type and Weil data, $\bn$ a finite index set, $\sW$ a hermitian space with index
set $\bn$, $\mu$ integer weights and $e$ a natural number with $|\bn|=2e+2\negc(\sW)$ and $e\le\mu(\tau)$ for
every infinite place $\tau$. Let $\sW^{+}$ be the sum of $\sW$ and one hyperbolic plane. Then there are a rapid
real datum $\fv$ for $(\sW^{+},\sW^{+})$ and a real number $\rC\ne0$ such that $\eqr{\sW^{+}}{\fc}{\fv}=\rC$ for
every split tame character $\fc$ of weight $\mu$.
\end{proposition}
\begin{proof}
Let $\mathsf{J}$ be the standard definite space of rank $2e$, and $\sX_1$ a space with the signs of $\sW$ that
has a base-adapted splitting over $\mathsf{J}$ (Lemma~\ref{l:planes}\,(a)): it is $\mathsf{J}$ plus $\negc(\sW)$
hyperbolic planes. Its sum $\sX_1^{+}$ with one more plane has the signs of $\sW^{+}$, so Lemma~\ref{l:move}
applies, and it suffices to treat $\sX_1^{+}$, with a determinant vector as real datum. Apply
Proposition~\ref{p:partners} twice, with $\sX=\sX_1$ and $\sW=\mathsf{J}$, to $\mathsf{J}$ and the zero space in the
two orders: the product of the two identities, evaluated at the standard datum with
Lemma~\ref{l:planes}\,(d), shows that the square of the ratio is a product $B_1\,\rC_1\,|\langle\text{vacuum at }0\rangle|^2$
of numbers that do not depend on $\fc$, the $L$-values cancelling because the two exchange constants involve
$\LP{\fc}{e'}$ and $\LP{\fc}{-e'}$. If the ratio vanished for all $\fc$, the descended identity at the standard
datum would give $0=1$ at a point where the determinant vector does not vanish
(Proposition~\ref{p:descent}\,(2) keeps determinant vectors). So the constant is non-zero.
\end{proof}
\leanline{\lean{existsequalRankThetaRatiovalue}}

\begin{proof}[Proof of Theorem 1]
Put $e_1=|\bn|/2-\negc(\sW)$; by the hypotheses $e_1\ge1$ is a natural number and $e_1\le\mu(\tau)$.
Proposition~\ref{p:equal} gives a rapid $\mathfrak{z}$ and $\rC_0\ne0$ with $\eqr{\sW^{+}}{\fc}{\mathfrak{z}}=\rC_0$
for all $\fc$. Proposition~\ref{p:partners}, for $(\sV',\sW,\sV)$, the splitting of $\sW^{+}$ over $\sW$ that
removes its plane, and the data $\mathfrak{z}$, $\fv$, gives $\rC_1$ and a rapid $\fx$; its sign condition is
$\negc(\sW)=\negc(\sV)+\negc(\sV')$. Take $\rC=\rC_0/\rC_1$ and $\fx$ (if $\rC_1=0$ the quotient is read as $0$;
this happens only when there is no tame triple at all, see below). Let $\omega$ be tame for $\fm$ with weights
$\mu$ and character $\chi$. The weights are positive by the bound on $\mu$, the set $\Sigma(\fm)$ is non-empty and
contains no conjugate of its members, and the tame part of $\omega$ is non-trivial exactly on it: this is a split
tame character $\fc$ with $\Sigma(\fc)=\Sigma(\fm)$. For it Proposition~\ref{p:partners} gives $\rc$ with
$\rC_1|\rc|^2=\LP{\fc}{|\ba|-|\bn|/2}\,\rC_0$; since $|\bn|=|\ba|+|\bb|=2|\bb|+2e$ the argument is $e\ge0$ and
$\LP{\fc}{e}=\LPP{\Sigma(\fm)}{\omega}{e}>0$. The right-hand side is not zero; in particular $\rC_1\ne0$ and
$\rc\ne0$, and dividing by $\rC_1$ gives $|\rc|^2=\rC\,\LPP{\Sigma(\fm)}{\omega}{e}$. Finally
let $\pi_0$ be the surjective ring homomorphism from the integers of $\sF_{\fc}$ onto $\cO_E/\fm$ that agrees with
the reduction of adelic coordinates. A function on matrices over $\cO_E/\fm$ pulls back along $\pi_0$, entry by
entry, to a function on the integral matrices over $\sF_{\fc}$, and is extended by zero off the integral matrices.
Partners pull back to locally constant partners over the integers of $\sF_{\fc}$, and the data modulo $\fm$ are the
local data of the pulled-back functions. The identity of Proposition~\ref{p:partners} becomes the asserted one.
\end{proof}

\sect{One square}{s:one}
\begin{lemma}[three tools]\label{l:tools}
Let $E$ be a CM field with a CM type; let $\sA$, $\sU_1$, $\sU_2$ be hermitian spaces with index sets $\bn'$,
$\bn_1$, $\bn_2$.
\begin{enumerate}[label=\textup{(\alph*)},leftmargin=2.3em,itemsep=1pt]
\item Let $\rc_i$, $G_i$ be complex numbers and $F_i$ functions on $([\U(\sU_1)]\times[\U(\sU_2)])\times[\U(\sA)]$,
indexed by a finite set. If every $F_i$ with $\rc_i\ne0$ is integrable with integral $G_i$, then the iterated
integral of $\sum_i\rc_iF_i$ is $\sum_i\rc_iG_i$.
\item Let $y$ be a form on $\sA$, $F$ a form on $\sU_1$, $F'$ a form on $\sU_2$, let $\mathsf{u}$, $\mathsf{d}$,
$\mathsf{l}$ be linear maps and $\rc$, $\rc'$, $\rC$ complex numbers. If $y=\rc\,\mathsf{u}(F)$,
$F'=\rc'\,\mathsf{d}(y)$ and $\rc'\rc=\rC$, then $\mathsf{l}(F')=\rC\,\mathsf{l}(\mathsf{d}(\mathsf{u}(F)))$. (The formal
statement writes $F$ and $F'$ as their right translates by the unit element.)
\item Suppose Weil data are given. Let $y$ and $y'$ be functions on $\U(\sU_1)(\A)$ and $\U(\sU_2)(\A)$ that are
automorphic and rapidly decreasing, $\omega_1$, $\omega_2$ characters of $\U(1)(\A)$, and $\mathcal{K}_0(\fv,\sg,\sz)$ a
function of a vector $\fv$ of the Fock model, a point $\sg$ of $\U(\sU_1\oplus\sU_2)$ and a point $\sz$ of
$\U(\sA)$, such that $|\mathcal{K}_0(\fv,\sg,\sz)|$ is at most a constant times a weighted norm of $\fv$ times a power of the
height of $\sg$ whenever that weighted norm is finite, and such that each $\mathcal{K}_0(\fv,\cdot,\cdot)$ is continuous.
Then there are a weighted norm and a constant such that for every $\fv$ of finite norm the function
$y(\sq_1)\,y'(\sq_2)\,(\omega_1(\det\sq_1)\omega_2(\det\sq_2))^{-1}\mathcal{K}_0(\fv,\sq_1\oplus\sq_2,\sz)$ is integrable over
$[\U(\sU_1)]\times[\U(\sU_2)]\times[\U(\sA)]$ and bounded by the constant times the norm of $\fv$.
\end{enumerate}
\end{lemma}
\begin{proof}
(a) Fubini and linearity. (b) Substitute. (c) The rapid decay of $y$, $y'$ absorbs the polynomial growth of $\mathcal{K}_0$
in the height; the function is then bounded and continuous on a space of finite measure.
\end{proof}
\leanline{\lean{integralintegralsumeq}\sep \lean{applyeqmulofformRepone}\sep \lean{existsintegrablekernelbound}}

\begin{proof}[Proof of Theorem 2]
Apply Theorem~1 to $\sW=\sU_1\oplus\sU_2$, $\sV=\sA$, $\sV'=\sA'$; its hypotheses are the numerical hypotheses of
Theorem~2. It gives $\rC$ and a rapid datum, and for each $i$ a constant $\rc\ne0$ with
$|\rc|^2=\rC\,\LPP{\Sigma(\fm_i)}{\omega_i}{e}$. The first two conditions of the square law, the equivariance of
$\sx\mapsto\mathsf{l}(\mathsf{u}_i(\sx)F)$ and $\sy\mapsto\mathsf{l}(\mathsf{v}_i(\sy)F)$ under the general linear groups of the blocks, come
from the action of the Weil representation on the indicator data. For the partner identity write both
composites as iterated integrals of $F$ against theta kernels: the composite through $\sA$ is the integral of $F$
against the theta integral over $[\U(\sA)]$ of the pair $(\sU_1\oplus\sU_2,\sA)$, restricted to
$\U(\sU_1)\times\U(\sU_2)$, by the adjointness of lift and descent (Input~4); the composite through $\sA'$
likewise for the pair $(\sA',\sU_1\oplus\sU_2)$. The integrals converge and may be interchanged with the finite
sums over matrices by Lemma~\ref{l:tools}, the kernel bound coming from the decay of the theta kernel (Input~4).
Summing against partners $\dL$, $\dS$ turns the indicator data into the data of $\dL$ and $\dS$, and the exchange
identity of Theorem~1 identifies the two sides up to $\rc$. The real data of the edges are polynomial vectors,
while Theorem~1 provides only some rapid datum on one side: the integrals for polynomial data lie in a fixed
line, which is closed, and approximating the datum of Theorem~1 by polynomial vectors gives the number
$\varkappa$ that compares the two; it is not shown to be non-zero. Hence the constant $\varkappa/\rc$.
\end{proof}

\sect{Triangles}{s:tri}
\begin{definition}[triangles and weighted chains]\label{d:tri}
Let $\cI$ be an index set, $\cF_{i,n}$ ($i,n\ge0$) complex inner product spaces and $\cV$ a complex inner product
space. The elements of $\cI$ are written $\vartheta$, because in the only application (Theorem~3) $\cI$ is the
set of the $m$-fold twists of the data; in this definition and in Lemma~\ref{l:chain} $\cI$ is an arbitrary set.
Let $m$ be a natural number; $\mathrm{w}$ a family of natural numbers $\mathrm{w}_\vartheta(n)$
($\vartheta\in\cI$), the \kw{weights}; $\cM$ a relation between natural numbers $n$ and families
$\rc=(\rc_\vartheta)_{\vartheta\in\cI}$ of complex numbers, read ``$\rc$ has grade at most $n$''; $\mathsf{a}$ and
$\mathsf{d}$ families of linear maps $\mathsf{a}_{\vartheta,i,n}\colon\cF_{i,n}\to\cF_{i,n+1}$ and
$\mathsf{d}_{\vartheta,i,n}\colon\cF_{i,n+1}\to\cF_{i+1,n}$; $\mathsf{j}_i\colon\cV\to\cF_{i,0}$ linear isometries;
and $\mathrm{G}$ a proposition.
\begin{enumerate}[label=\textup{(\alph*)},leftmargin=2.3em,itemsep=1pt]
\item $\IsTriangle{m}{\mathrm{w}}{\cM}{\mathsf{a}}{\mathsf{d}}{\mathsf{j}}{\mathrm{G}}$ holds if: \textup{(1)} the
constant family $1$ has grade at most $0$; \textup{(2)} if $\rc$ has grade at most $n$, $\rc'$ grade at most $n'$
and $n+n'\le n''$, then $\rc\rc'$ has grade at most $n''$; \textup{(3)} for $i+(n+1)<m$ there is a family
$\mathrm{b}$ with
$\mathrm{w}_\vartheta(n)\ip{y''}{\mathsf{a}_{\vartheta,i+1,n}\mathsf{d}_{\vartheta,i,n}y}=\overline{\mathrm{b}_{\vartheta}}^{\,-1}\mathrm{w}_\vartheta(n+2)\ip{y''}{\mathsf{d}_{\vartheta,i,n+1}\mathsf{a}_{\vartheta,i,n+1}y}$
for all $\vartheta$, $y$, $y''$, and such that, if $\mathrm{G}$ holds, $\mathrm{b}$ has grade at most $1$ and there is
a family $\mathrm{k}$ of grade at most $0$ such that for all $\vartheta$ and $y''$ some $y'$ satisfies
$\ip{\mathsf{a}_{\vartheta,i,n+1}y}{y''}=\ip{y}{y'}$ for all $y$ and
$\ip{\mathsf{d}_{\vartheta,i,n}y'}{y}=\mathrm{k}_{\vartheta}\ip{\mathsf{d}_{\vartheta,i,n+1}y''}{\mathsf{a}_{\vartheta,i+1,n}y}$
for all $y$; \textup{(4)} for $i<m$, all $\vartheta$ and $y\in\cV$:
$\ip{y''}{\mathsf{a}_{\vartheta,i,0}\mathsf{j}_iy}=\ip{\mathsf{d}_{\vartheta,i,0}y''}{\mathsf{j}_{i+1}y}$ for all $y''$, and
$\mathsf{d}_{\vartheta,i,0}\mathsf{a}_{\vartheta,i,0}\mathsf{j}_iy=\mathsf{j}_{i+1}y'$ for some $y'\in\cV$ that
vanishes only if $y$ does.
\item $\weightedChain{\mathrm{w}}{\mathsf{a}}{\vartheta}{i}$ is the family, indexed by $n$, of the composites
$\cF_{i,0}\to\cF_{i,n}$ of the first $n$ maps $\mathsf{a}_{\vartheta,i,n'}$, the one of index $n'$ multiplied by
$\mathrm{w}_\vartheta(n')$; its $n$-th member is written $\weightedChain{\mathrm{w}}{\mathsf{a}}{\vartheta}{i}_n$.
\end{enumerate}
\end{definition}
\leanline{\lean{IsTriangle}\sep \lean{weightedChain}}

\begin{lemma}[weighted chains in a triangle]\label{l:chain}
Let $\cI$, $\cF_{i,n}$, $\cV$ be as in Definition~\ref{d:tri}, and suppose that
$\IsTriangle{m}{\mathrm{w}}{\cM}{\mathsf{a}}{\mathsf{d}}{\mathsf{j}}{\mathrm{G}}$ holds and that all weights are
non-zero.
\begin{enumerate}[label=\textup{(\alph*)},leftmargin=2.3em,itemsep=1pt]
\item Let $i$, $n$ be natural numbers with $i+n<m$. There are families $\mathrm{b}$, $\mathrm{k}$ such that for all
$\vartheta$ and $y$
\[\weightedChain{\mathrm{w}}{\mathsf{a}}{\vartheta}{i+1}_n\bigl(\mathsf{d}_{\vartheta,i,0}\mathsf{a}_{\vartheta,i,0}y\bigr)
=\frac{\mathrm{w}_\vartheta(n+1)}{\mathrm{w}_\vartheta(1)\mathrm{w}_\vartheta(0)\overline{\mathrm{b}_{\vartheta}}}\;
\mathsf{d}_{\vartheta,i,n}\bigl(\weightedChain{\mathrm{w}}{\mathsf{a}}{\vartheta}{i}_{n+1}y\bigr),\]
and such that, if $\mathrm{G}$ holds, $\mathrm{k}\mathrm{b}$ has grade at most $n$ and
\[\ip{y''}{\mathsf{a}_{\vartheta,i,n}\weightedChain{\mathrm{w}}{\mathsf{a}}{\vartheta}{i}_n\mathsf{j}_iy}
=\mathrm{k}_{\vartheta}\ip{\mathsf{d}_{\vartheta,i,n}y''}{\weightedChain{\mathrm{w}}{\mathsf{a}}{\vartheta}{i+1}_n\mathsf{j}_{i+1}y}\]
for all $\vartheta$, $y$, $y''$.
\item For natural numbers $n$, $i$ with $i+n\le m$, every $\vartheta$ and every $y\in\cV$: if
$\weightedChain{\mathrm{w}}{\mathsf{a}}{\vartheta}{i}_n(\mathsf{j}_iy)=0$ then $y=0$.
\item Suppose $\mathrm{G}$ holds. Let $\cV'$ be a complex vector space, $\mathsf{s}\colon\cV'\to\cV$ a linear map
and $\mathsf{S}_{\vartheta,i}$ endomorphisms of $\cV'$ with
$\mathrm{w}_\vartheta(0)\ip{\mathsf{a}_{\vartheta,i,0}\mathsf{j}_iy'}{\mathsf{a}_{\vartheta,i,0}\mathsf{j}_i\mathsf{s}y}=\ip{y'}{\mathsf{s}\,\mathsf{S}_{\vartheta,i}y}$
for $i<m$ and all $\vartheta$, $y'$, $y$. Then for all natural numbers $n$, $i$ with $i+n=m$ there is a family
$\mathrm{k}$ of grade at most $n\cdot n$ such that every $\vartheta$ has an endomorphism $\mathsf{T}$ of $\cV'$ with
$\det\mathsf{T}=\prod_{i\le i'<m}\det\mathsf{S}_{\vartheta,i'}$ and, for all $y'$, $y$,
\[\ip{\weightedChain{\mathrm{w}}{\mathsf{a}}{\vartheta}{i}_n\mathsf{j}_iy'}{\weightedChain{\mathrm{w}}{\mathsf{a}}{\vartheta}{i}_n\mathsf{j}_i\mathsf{s}y}
=\mathrm{k}_{\vartheta}\,\frac{\mathrm{w}_\vartheta(0)\,\mathrm{w}_\vartheta(1)^{n}}{\mathrm{w}_\vartheta(n)}\,\ip{y'}{\mathsf{s}\,\mathsf{T}y}.\]
\end{enumerate}
\end{lemma}
\begin{proof}
(a) Induction on $n$. For $n=0$ the identity is trivial and the adjoint law is condition (4) of a triangle. The
step moves one more square past the chain by condition (3): the commutation with the scalar
$\overline{\mathrm{b}}^{-1}\mathrm{w}(n+2)/\mathrm{w}(n)$ gives the first identity, the weights telescoping, and the
adjoint law of (3) under $\mathrm{G}$ gives the second, the grades adding by (2). (b) Induction on $n$: by (a),
$\mathsf{d}_{\vartheta,i,n-1}$ of the chain of length $n$ of row $i$ is a non-zero multiple of the chain of length
$n-1$ of row $i+1$ applied to $\mathsf{j}_{i+1}y'$, with $y'=0$ only if $y=0$ by (4). (c) Induction on $n$: move the
first arrow of row $i$ down by (a), use the formula for row $i+1$, and compose the endomorphisms; the grade
$n\cdot n$ is the sum of the grades collected.
\end{proof}
\leanline{\lean{weightedChaindescent}\sep \lean{eqzeroofweightedChain}\sep \lean{innerweightedChaineq}}

\sect{Costs in squares}{s:cost}
\begin{definition}[costs]\label{d:cost}
Let $P$ be parameters and $D$ admissible data, with first scale $x$ and $t(D)$ as in the handout. For a natural
number $n$ and a complex function $\rc$ on the $m$-fold twists of $D$, the statement $\CostsSquares{D}{n}{\rc}$
(``$\rc$ costs $n$ squares'') holds if for every fit family $\theta=(\theta_{\ell,\iota})$ of $m$-fold twists of
$D$, one twist for each prime $\ell$ and each isomorphism $\iota$ of $\CC$ with an algebraic closure of the
$\ell$-adic numbers, the family $(\ell,\iota)\mapsto|\rc(\theta_{\ell,\iota})/\rc(1)|^{-2}$ has height at most
$n\cdot(2\,t(D)\,x^{999})$; here, as in the formal text, the inverse of $0$ is $0$.
\end{definition}
\leanline{\lean{CostsSquares}}

\begin{lemma}[rules for costs]\label{l:cost}
Let $P$ be parameters and $D$ admissible data.
\begin{enumerate}[label=\textup{(\alph*)},leftmargin=2.3em,itemsep=1pt]
\item If $\CostsSquares{D}{n}{\rc}$ and $\CostsSquares{D}{n'}{\rc'}$ hold and $n+n'\le n''$, then
$\CostsSquares{D}{n''}{\rc\rc'}$ holds.
\item For a natural number $e$ and a tame character $\vartheta$: the residue ring of $\fm_D(\vartheta,e)$ has
$(\prod_{\fp\in\Sigma_D}\Np\fp)^{e}\cdot\Np Q(\vartheta)$ elements.
\item Let $e$ be a natural number, $\rC$ a real number and $i_0$ an integer. If
$|\rc(\vartheta)|^2=\rC\cdot\#(\cO_E/\fm_D(\vartheta,e))^{i_0}$ for every $m$-fold twist $\vartheta$, then
$\CostsSquares{D}{0}{\rc}$ holds.
\item $\CostsSquares{D}{0}{1}$ holds.
\item Let $\vartheta$ be a tame character, $\omega$ a character of $\U(1)(\A)$, $\mu$ weights and $\chi$ a
multiplicative character of $\cO_E/\fm_D(\vartheta,1)$. If $\omega$ is tame for $\fm_D(\vartheta,1)$ with weights
$\mu$ and character $\chi$, then $\chi$ is primitive.
\item Let $\pp$ be a pair of $D$ that is ramified exactly at the chosen places, $e\ne0$ a natural number and $n$
a lift index, $1\le n\le m$. Then for every $m$-fold twist $\vartheta$ the character
$(\omega_{\pp}\vartheta^{n})^{-1}$ is tame for $\fm_D(\vartheta,e)$, with the weights $\mu_{\pp}$ (the character
$\omega_{\pp}$ itself has the weights $-\mu_{\pp}$) and the character of the tame parts; and every $\rc$ with
$|\rc(\vartheta)|^2=\rC\cdot\LPP{\Sigma(\fm_D(\vartheta,e))}{(\omega_{\pp}\vartheta^{n})^{-1}}{e_{\pp}}$ for all
$\vartheta$ and some real $\rC$ satisfies $\CostsSquares{D}{1}{\rc}$.
\end{enumerate}
\end{lemma}
\begin{proof}
(a) Heights add in products, and a bound for a height may be enlarged (Input~8). (b) is a count (Input~6). (c)
For a fit family, $|\rc(\theta_{\ell,\iota})/\rc(1)|^{-2}$ is a power of $\Np Q(\theta_{\ell,\iota})$ by (b); the
height of the constant family is~$0$, and a family that differs from one of height at most $h_0$ by powers of the norms of the moduli of a fit
family has height at most any $h_1\ge h_0$ (Lemma~\ref{l:meas}\,(c)). (d) is (c) with exponent $0$. (e) A non-zero
element $\mathsf{a}$ of the residue ring has a non-zero residue at some place of the modulus; the tame part there
is non-trivial, so it takes a value $\ne1$ at some unit, and lifting gives $\mathsf{c}$ with
$\chi(1+\mathsf{c}\mathsf{a})\notin\{0,1\}$. (f) Tameness follows from the weights (Input~2, Input~6). By
Lemma~\ref{l:ratio}\,(b) the product of $L$-values is $\Np Q(\vartheta)^{e_{\pp}^2}$ times its value at the trivial
twist times the squares of the ratios of twisted to untwisted $L$-values at the points $1,\dots,e_{\pp}$; for a
fit family the inverse square of each of these $e_{\pp}$ ratios has height at most $2x^{999}$ (Input~7), so
together they have height at most $2e_{\pp}x^{999}\le2\,t(D)\,x^{999}$, since $e_{\pp}\le t(D)$ for a pair; and
the power of the norm is harmless as in (c).
\end{proof}
\leanline{\lean{mul}\sep \lean{card}\sep \lean{ofnormsqeqmulcardzpow}\sep \lean{one}\sep
\lean{isPrimitiveMulCharofisTameFor}\sep \lean{existsisTameForcostsSquaresone}}

\sect{The lift of the trivial character from $\U(1)$}{s:lvalue}
\begin{definition}[$L$-value data and the lift of the trivial character]\label{d:ldata}
Let $E$ be a CM field with a CM type and $\bn$ a finite index set.
\begin{enumerate}[label=\textup{(\alph*)},leftmargin=2.3em,itemsep=1pt]
\item If $\bn$ is non-empty, one index of $\bn$ is chosen once and for all (an embedding of the one-point set).
\item $\LValueData{E}{\bn}$ is the set of \kw{$L$-value data}: such a datum $\fl$ consists of level places of $E$
(chosen places and twist places), jet degrees $\mathrm{dg}_{\fl}(\tau)\in\mathbb{N}$ at the infinite places, and a
character $\omega_{\fl}$ of $\U(1)(\A)$ of weights $\tau\mapsto-(\mathrm{dg}_{\fl}(\tau)+\lfloor|\bn|/2\rfloor)$ whose
tame part is trivial off the places; moreover $|\bn|>1$.
\item $\auxPair{\fl}$ is the auxiliary pair with the places of $\fl$, the definite line $\sI_1$, the chosen index
of $\bn$, the weight with one row of length $\mathrm{dg}_{\fl}(\tau)$ at each $\tau$, and the character
$\omega_{\fl}$.
\item Suppose Weil data are given. For a scale $\rc$, a datum $f$ and finite adelic points $\sg$ of $\U(\sI_1)$ and
$\sh$ of $\U(\sI_{\bn})$, $\kernelValue{\fl}{\rc}{f}{\sg}{\sh}$ is the family of jets, for the one-row weight of
$\fl$ and with derivatives scaled by $\rc$, of the datum $f$ moved by $(\sg,\sh)$ under the Weil representation
of the pair of $\auxPair{\fl}$.
\item $\liftKernel{\fl}$ is the theta kernel from $\sI_1$ to $\sI_{\bn}$ for the weight and the characters of
$\auxPair{\fl}$, with scale $1$ and the full level on $\U(\sI_1)$.
\item $\kernelForm{\fl}{\sg}$ is the algebraic form $\sh\mapsto\kernelValue{\fl}{1}{\Upsilon}{\sg}{\sh}$ on
$\sI_{\bn}$, where $\Upsilon$ is the finite datum of $\auxPair{\fl}$; it has the level, character and weight of
$\auxPair{\fl}$.
\item $\chosenLevel{\fl}$ is the principal level on the finite adelic points of $\U(\sI_1)$ at the chosen places
of $\fl$.
\item $\trivialLift{\fl}$, the \kw{lift of the trivial character}, is the sum of the forms $\kernelForm{\fl}{\sg}$
over the classes $\sg$ of $\U(\sI_1)$ at the level $\chosenLevel{\fl}$.
\end{enumerate}
\end{definition}
\leanline{\lean{chosenIndexEmb}\sep \lean{LValueData}\sep \lean{auxPair}\sep \lean{kernelValue}\sep
\lean{liftKernel}\sep \lean{kernelForm}\sep \lean{chosenLevel}\sep \lean{trivialLift}\sep the characters of
$\auxPair{\fl}$ have the expected weights: \lean{ahasWeight}, \lean{bhasWeight}}

\begin{definition}[pairs of the data and the algebraic part]\label{d:pairs}
Let $P$ be parameters, $D$ admissible data, $\pp$ a pair of $D$ and $\bn$ a finite index set.
\begin{enumerate}[label=\textup{(\alph*)},leftmargin=2.3em,itemsep=1pt]
\item $\HasLiftRank{\pp}{\bn}$ (``$\pp$ has lift rank $\bn$'') holds if $|\bn|=2e_{\pp}$ and $e_{\pp}\ge2$.
\item $\IsFlat{\pp}{\bn}$ (``$\pp$ is flat'') holds if $\mu_{\pp}(\tau)=e_{\pp}$ at every infinite place $\tau$ and
$\pp$ has lift rank $\bn$.
\item $\denomBound{\pp}=g(D)^{2f(D)\cdot2e_{\pp}}\cdot(g(D)\,d)^{\sum_\tau(\mu_{\pp}(\tau)-e_{\pp})}$, a natural
number; $d$ is the prime of the parameters.
\item Suppose $\pp$ has lift rank $\bn$, and let $\vartheta$ be a tame character. The formal definitions that
follow take the proof of $\HasLiftRank{\pp}{\bn}$ as their first argument; we write the pair and the index set in
its place. $\lValueData{\pp,\bn}{\vartheta}$ is the $L$-value datum for $E(D)$ and $\bn$ with the level places of
$\vartheta$ (chosen places $\Sigma_D$, twist places $Q(\vartheta)$), the jet degrees $\mu_{\pp}(\tau)-e_{\pp}$ and
the character $\omega_{\pp}\vartheta$.
\item $\trivialLifta{\pp,\bn}{\vartheta}=\mathrm{vs}^{|\bn|}\,(\bar\pi_g\,\mathrm{per})^{\sum_\tau(\mu_{\pp}(\tau)-e_{\pp})}\cdot\trivialLift{\lValueData{\pp,\bn}{\vartheta}}$,
the scaled lift of the trivial character; $\mathrm{vs}$ is the vacuum scale and $\mathrm{per}$ the jet scale of the
data.
\item $\intAlgPart{\pp,\bn}{\vartheta}=\mathrm{M}_D\cdot\denomBound{\pp}^2\cdot\|\trivialLifta{\pp,\bn}{\vartheta}\|^2\big/\Np Q(\vartheta)^{|\bn|-1}$,
the \kw{algebraic part}; here $\mathrm{M}_D=\#\prod_{\fp\in\Sigma_D}\GL_{r\cdot t(D)}(\cO_E/\fp)$ is the order of the
product, over the chosen places, of the general linear groups of rank $r\cdot t(D)$ of the residue fields.
\item $\pointDatum{\pp,\bn}{\vartheta}$ is $\mathrm{vs}^{|\bn|}$ times the datum that is the level indicator of
$\auxPair{\lValueData{\pp,\bn}{\vartheta}}$ modulo its places, on the vacuum, and
\[\pointKernel{\pp,\bn}{\vartheta}=\kernelValue{\lValueData{\pp,\bn}{\vartheta}}{\bar\pi_g\,\mathrm{per}}{\pointDatum{\pp,\bn}{\vartheta}}{\cdot}{\cdot},\]
a function of finite points of $\U(\sI_1)$ and $\U(\sI_{\bn})$.
\item For an integer $i$, $\lowerPoint{\pp}{i}$ is the pair with the character of $\pp$ and the point
$\min(\max(1,i),e_{\pp})$.
\end{enumerate}
\end{definition}
\leanline{\lean{HasLiftRank}\sep \lean{IsFlat}\sep \lean{denomBound}\sep \lean{lValueData}\sep \lean{trivialLifta}\sep
\lean{intAlgPart}\sep \lean{pointDatum}\sep \lean{pointKernel}\sep \lean{lowerPoint}}

\begin{definition}[integral measures]\label{d:intmeas}
Let $P$ be parameters, $D$ admissible data, $\vartheta$ a tame character of $D$ and $G$ a complex function on
the tame characters of $D$. $\IsIntegralMeasure{\vartheta}{G}$ holds if there are finitely many points $\mathsf{u}_i$
of $\U(1)(\A)$ and complex numbers $\rc_i$ such that $|j(\rc_i)|\le1$ for every prime $\ell$ that does not divide
$\Np Q(\vartheta)$ and every ring homomorphism $j$ from $\CC$ to the completed algebraic closure $\CC_\ell$ of the
$\ell$-adic numbers, and such that $G(\vartheta')=\sum_i\rc_i\,\vartheta'(\mathsf{u}_i)$ for every tame character
$\vartheta'$ with $Q(\vartheta')=Q(\vartheta)$.
\end{definition}
\leanline{\lean{IsIntegralMeasure}}

\subhead{Small facts used in the proofs.} They are not stated as lemmas; each is a single step in the formal
text. For data $D$: the rank of $\sY_i(D)$ is at most that of $\sY_m(D)$ for $i\le m$; the index set of the spaces
$\sV$ of the grid has $r+t(D)$ elements; the number of row indices whose number lies in an interval $[i,i')$
with $i'\le r$ is $i'-i$; the parameter of the rows is strictly decreasing at each infinite place, and the
doubled parameter of $\sW_i(D)$ is ordered against the numbering of the basis; exactly $m-i$ basis vectors of
$\sW_i(D)$ are negative. For $L$-value data $\fl$: the weight of $\auxPair{\fl}$ is a lift weight from the line;
the Weyl representation of a one-row weight is trivial on the line; the class-sum lift of a form $y$ at
$(1,\sg')$ is $\ip{\kernelForm{\fl}{\sg'}}{y}$, and $\ip{\trivialLift{\fl}}{y}$ is the class sum of
$\mathsf{u}\mapsto\theta_{\auxPair{\fl}}(y)(1,\mathsf{u}\sg)$. For a pair of lift rank $\bn$: the ratio of
Lemma~\ref{l:ratio} is $\intAlgPart{\pp,\bn}{\vartheta}/\intAlgPart{\pp,\bn}{1}$; $\intAlgPart{\pp,\bn}{\vartheta}\ne0$;
and $\trivialLifta{\pp,\bn}{\vartheta}\ne0$ (all three from Theorem~5 and Input~6). The passages to the opposite
space of Definition~\ref{d:triforms} commute with right translation. A modulus $\fm$ of $E(D)$ whose ideal contains
a natural number $e_2>0$ has a residue ring with at most $e_2^{2f(D)}$ elements. There are a number and an integer
exponent such that for every $m$-fold twist the index of the last level of the chain is that number times a power
of $\#(\cO_E/\fm_D(\vartheta,0))$ times the $m$-th power of the Gram factor of the first auxiliary pair times
$\mathrm{ord}_{D,\mathrm{false}}(\vartheta,m)/(\mathrm{ord}_{D,\mathrm{false}}(\vartheta,0)\,\mathrm{ord}_{D,\mathrm{false}}(\vartheta,1)^m)$,
with the triangle orders of Definition~\ref{d:triforms}\,\textup{(f)}.
\leanline{\lean{cardYle}\sep \lean{cardV}\sep \lean{cardfilterbasisNumber}\sep \lean{varsigmastrictAnti}\sep
\lean{paramWltiff}\sep \lean{cardnegsignsW}\sep \lean{isLiftWeight}\sep \lean{conjWeylRepeqone}\sep
\lean{thetaeqinnerkernelForm}\sep \lean{innertrivialLift}\sep \lean{ratioeqintAlgPartdiv}\sep
\lean{intAlgPartnezero}\sep \lean{trivialLiftnezero}\sep \lean{triNegEquivformRep}\sep \lean{cardresRingle}\sep
\lean{stdLevelIndexlastResLeveleq}}

\begin{lemma}[parameters of the chain]\label{l:params}
Let $P$ be parameters and $D$ admissible data.
\begin{enumerate}[label=\textup{(\alph*)},leftmargin=2.3em,itemsep=1pt]
\item At each infinite place the doubled parameter of $\sY_m(D)$ has distinct entries and is twice the
Harish-Chandra parameter of the weight of the forms on $\sY_m(D)$ with its last entry.
\item For a tame character $\vartheta$ and a lift index $i$, $1\le i\le m$: the doubled parameter of $\sW_i(D)$
fits the $i$-th auxiliary pair of $(D,\vartheta)$.
\end{enumerate}
\end{lemma}
\begin{proof}
Both are verifications entry by entry. The Harish-Chandra parameter of a weight is computed by steps along its
columns (Input~7); the indices of $\sY_m(D)$, the block of the rows and $m$ blocks of size $t(D)$, are matched
with these steps. For (b) the same formula gives the parameter of the definite space of the auxiliary pair, and
the four conditions of fitting are checked with the small facts above.
\end{proof}
\leanline{\lean{isHCParamY}\sep \lean{fitsParamsauxPair}}

\begin{lemma}[inner products of algebraic forms]\label{l:petersson}
Let $E$ be a CM field with a CM type, $\bn$ a finite index set and $\sV$ a hermitian space with index set $\bn$
that is a standard space. Let a weight and its parameter $l$ be given such that $l$ is the Harish-Chandra
parameter of the weight, and let $F$, $F'$ be algebraic forms on $\sV$ of that weight, of a level and a character
at given level places. Then $\ip{F}{F'}$ is the index of the level in the stabilizer of the integral lattice,
times the number of elements of the Weyl index set of the weight, times the Petersson product of the functions on
$\U(\sV)(\A)$ attached to $F$ and $F'$.
\end{lemma}
\begin{proof}
Unfold the Petersson product and integrate the class function over the classes of the level.
\end{proof}
\leanline{\lean{innereqpetersson}}

\begin{lemma}[the kernel forms]\label{l:kform}
Let $E$ be a CM field with a CM type and Weil data, $\bn$ a finite index set and $\fl\in\LValueData{E}{\bn}$.
\begin{enumerate}[label=\textup{(\alph*)},leftmargin=2.3em,itemsep=1pt]
\item The map $\sg\mapsto\kernelForm{\fl}{\sg}$ is invariant on the left under the rational points of $\U(\sI_1)$
and on the right under $\chosenLevel{\fl}$.
\item For every finite point $\sg$ of $\U(\sI_1)$: $\trivialLift{\fl}$ is the sum of the forms
$\kernelForm{\fl}{\mathsf{u}\sg}$ over the classes $\mathsf{u}$.
\item For a real point $\sx$ of $\U(\sI_1)$, a vector $\mathsf{w}$ of the weight, a real point $\sy$ of $\U(\sI_{\bn})$
and finite points $\sg$, $\sh$: the theta functional of the Fock datum of $\mathsf{w}$ moved by $((\sx,\sg),(\sy,\sh))$
is the Gram form of $\bar{\mathsf{w}}$ with the image of $\kernelForm{\fl}{\sg}(\sh)$ under the conjugate Weyl
representation of $\sy^{-1}$.
\item Suppose that all jet degrees of $\fl$ vanish. Then for finite points $\sg$, $\sh$ and every coordinate $i$:
$\kernelForm{\fl}{\sg}(\sh)_i$ is the value at $((1,\sg),(1,\sh))$ of the theta kernel of $(\sI_1,\sI_{\bn})$ with
the characters $1$ and that of $\auxPair{\fl}$, the polynomial $1$ and the finite datum of the determinant
modulo the square of the chosen places times the twist places.
\end{enumerate}
\end{lemma}
\begin{proof}
(a) On the left this is the automorphy of the theta kernel; on the right, the finite datum $\Upsilon$ is fixed by
the elements of the maximal compact subgroup that reduce to $1$ at the chosen places. (b) A class sum of a
function invariant on both sides is unchanged by a translation of its argument, $\U(\sI_1)$ being commutative.
(c) Separate the Weil representation into its real and finite parts; the theta functional of the Fock datum is a
Gram form with the Schur jets. (d) Without derivatives the jets are the theta values, and the coefficient
polynomial is the constant~$1$.
\end{proof}
\leanline{\lean{kernelForminvariant}\sep \lean{trivialLifteqstdClassSumtranslate}\sep
\lean{thetapairWeilfockDatum}\sep \lean{kernelFormofjetDegeqzero}}

\begin{lemma}[the polynomials of a one-row weight]\label{l:poly}
Let $E$ be a CM field with a CM type, $\bn$ a finite index set and $\fl\in\LValueData{E}{\bn}$.
\begin{enumerate}[label=\textup{(\alph*)},leftmargin=2.3em,itemsep=1pt]
\item For a coordinate vector $\mathsf{w}$ of the weight, the polynomial of $\mathsf{w}$ is the harmonic polynomial of
$\auxPair{\fl}$ at $\mathsf{w}$ with its variables renamed injectively (each variable to the one with the same row
index and column $0$).
\item For such $\mathsf{w}$, the harmonic polynomial of $\auxPair{\fl}$ at $\mathsf{w}$ is the Schur coefficient polynomial of
$\mathsf{w}$ times the Weyl Gram matrix.
\item The hook number of the weight of $\auxPair{\fl}$ in rank $|\bn|$ is $\prod_\tau\mathrm{dg}_{\fl}(\tau)!$.
\end{enumerate}
\end{lemma}
\begin{proof}
(a) Both sides are the same intertwining map on Schur modules, checked on the Weyl evaluations. (b) Evaluate both
polynomials at every point, using (a). (c) The highest polynomial of a one-row weight is the monomial
with the exponents $\mathrm{dg}_{\fl}(\tau)$, whose Fock pairing with itself is the product of their factorials.
\end{proof}
\leanline{\lean{toPolyeqrenameharmonicPoly}\sep \lean{harmonicPolyeqschurCoeffPoly}\sep
\lean{hookNumbereqprodfactorial}}

\begin{lemma}[the lift as a sum of point kernels; integrality]\label{l:apply}
Let $P$ be parameters, $D$ admissible data, $\pp$ a pair of $D$ that has lift rank $\bn$, and $\vartheta$ a tame
character.
\begin{enumerate}[label=\textup{(\alph*)},leftmargin=2.3em,itemsep=1pt]
\item For every finite point $\sh$ of $\U(\sI_{\bn})$:
$\trivialLifta{\pp,\bn}{\vartheta}(\sh)=\sum_{(\mathsf{u},\mathsf{A})}\pointKernel{\pp,\bn}{\vartheta}(\mathsf{u}\,\tilde{\mathsf{A}},\sh)$,
the sum over representatives $\mathsf{u}$ of the classes of $\U(\sI_1)$ at the principal level at the chosen places
and over the elements $\mathsf{A}$ of the product of the unit groups of the residue fields at $Q(\vartheta)$,
with $\tilde{\mathsf{A}}$ a fixed lift.
\item Let $\ell$ be a prime with a homomorphism $j$ from $\CC$ to the completed algebraic closure of the
$\ell$-adic numbers, and $\sg$, $\sh$ finite points. If $\ell$ is prime to $\Np Q(\vartheta)$, then
$\denomBound{\pp}\cdot\pointKernel{\pp,\bn}{\vartheta}(\sg,\sh)$ is $j$-integral: its gauge at $\sh$ is at most $1$.
\end{enumerate}
\end{lemma}
\begin{proof}
(a) Unfold the lift as a class sum of kernel forms and write the finite datum $\Upsilon$ as the sum, over
$\mathsf{A}$, of the translates of the datum of the level indicator (Input~7). (b) The point datum, scaled, lies in
the standard data at $\ell$ (Lemma~\ref{l:std}), jets of standard data are divisible (Input~7), and the jet scale
and the boxes of the weight are cleared by $\denomBound{\pp}$.
\end{proof}
\leanline{\lean{trivialLiftapply}\sep \lean{weylGaugepointKernelleone}}

\begin{lemma}[integral measures and heights]\label{l:meas}
Let $P$ be parameters and $D$ admissible data.
\begin{enumerate}[label=\textup{(\alph*)},leftmargin=2.3em,itemsep=1pt]
\item For a tame character $\vartheta$ and a function $G$ with $\IsIntegralMeasure{\vartheta}{G}$:
$\IsIntegralMeasure{\vartheta}{\vartheta'\mapsto|G(\vartheta')|^2}$ holds.
\item For a prime $\ell$, an isomorphism $\iota$ of $\CC$ with an algebraic closure of the $\ell$-adic numbers and a
function $G$ with $\IsIntegralMeasure{1}{G}$: $|\iota(G(1))|\le1$.
\item Let $\theta=(\theta_{\ell,\iota})$ be a fit family of $m$-fold twists, $y$ and $y'$ families of complex
numbers indexed by the pairs $(\ell,\iota)$, and $h_0\le h_1$ real numbers. If $y$ has height at most $h_0$ and
every non-zero $y'_{\ell,\iota}$ equals $\Np Q(\theta_{\ell,\iota})^{i_0}\,y_{\ell,\iota}$ for an integer $i_0$,
which may depend on $(\ell,\iota)$, then $y'$ has height at most $h_1$.
\end{enumerate}
\end{lemma}
\begin{proof}
(a) $|G|^2=G\bar G$ is given by the points $\mathsf{u}_i\mathsf{u}_{i'}^{-1}$ and the coefficients $\rc_i\bar{\rc}_{i'}$.
(b) $G(1)$ is the sum of the coefficients, and the trivial character has empty modulus. (c) For a fit family the
prime $\ell$ is prime to $\Np Q(\theta_{\ell,\iota})$, so the factor is a unit at $\ell$: the witnesses of the
height of $y$ serve for $y'$ at every pair $(\ell,\iota)$, and the bound may be enlarged from $h_0$ to $h_1$
(Input~8).
\end{proof}
\leanline{\lean{normsq}\sep \lean{normapplyonele}\sep \lean{heightLeofeqzpowmul}}

\begin{proof}[Proof of Theorems 5 and 6]
\emph{Theorem 5.} By the first note (Input~1) the squared norm of $\trivialLift{\lValueData{\pp,\bn}{\vartheta}}$ is
a non-zero constant, depending only on the field and the rank, times explicit counts at the chosen and the twist places, the hook number,
the Weyl index, the value $L^{\Sigma\cup\bar\Sigma}(e_{\pp},((\omega_{\pp}\vartheta)^{-1})^{\mathrm{bc}})$ with
$\Sigma=\Sigma_D\cup Q(\vartheta)$, and $\Np Q(\vartheta)^{|\bn|-1}$; here $2e_{\pp}=|\bn|$ and $e_{\pp}>1$ by the
lift rank. The scaling factors of
Definition~\ref{d:pairs}\,(e), (f) do not depend on $\vartheta$, the hook number is given by
Lemma~\ref{l:poly}\,(c), and the power of $\Np Q(\vartheta)$ cancels against the normalisation of the algebraic
part. What remains is a constant $A\ne0$.

\emph{Theorem 6.} By Lemma~\ref{l:apply}\,(a) the lift is a finite sum of point kernels at the points
$\mathsf{u}\tilde{\mathsf{A}}$. Its squared norm expands into a double sum over pairs of these points; the
dependence on $\vartheta'$ is through the values of $\vartheta'$ at the quotients of the determinants of the
points, and the coefficients are $\mathrm{M}_D\,\Np Q(\vartheta)^{-(|\bn|-1)}$ times class sums of Gram forms of
the point kernels multiplied by the denominator bound. Let $\ell$ be a prime that does not divide
$\Np Q(\vartheta)$ and $j\colon\CC\to\CC_\ell$ a ring homomorphism. The Gram forms are $j$-integral by
Lemma~\ref{l:apply}\,(b). The inner product of algebraic forms on $\sI_{\bn}$ is a sum over the classes of the
level in which each class is weighted by the inverse of the order of its stabilizer; these orders divide
$\mathrm{M}_D$, so $\mathrm{M}_D$ times the class sum is $j$-integral, and this is why $\mathrm{M}_D$ stands in the
algebraic part. Finally $\Np Q(\vartheta)^{-(|\bn|-1)}$ is a unit at $\ell$. So $|j(\rc_i)|\le1$ for every
coefficient. The points and the coefficients do not depend on $\vartheta'$ as long as its modulus is that of
$\vartheta$.
\end{proof}

\begin{lemma}[ratios of $L$-values]\label{l:ratio}
Let $P$ be parameters, $D$ admissible data, $\pp$ a pair of $D$, $\vartheta$ a tame character and $\pm$ a sign.
Write, with $\Sigma=\Sigma_D\cup Q(\vartheta)$,
\[\operatorname{rat}(\pp,\vartheta,\pm)=\frac{\bigl|L^{\Sigma\cup\bar\Sigma}\bigl(e_{\pp},((\omega_{\pp}\vartheta^{\pm1})^{-1})^{\mathrm{bc}}\bigr)\bigr|}{\bigl|L^{\Sigma_D\cup\bar\Sigma_D}\bigl(e_{\pp},(\omega_{\pp}^{-1})^{\mathrm{bc}}\bigr)\bigr|}\]
for the ratio of the third note: the Euler factors are removed at the chosen places and their conjugates in both
values, and at $Q(\vartheta)$ and its conjugates in the numerator.
\begin{enumerate}[label=\textup{(\alph*)},leftmargin=2.3em,itemsep=1pt]
\item $\operatorname{rat}(\pp,\vartheta,\pm)$ is non-zero and real.
\item The products of $L$-values satisfy
\begin{multline*}
\LPP{\Sigma_D\cup Q(\vartheta)}{(\omega_{\pp}\vartheta^{\pm1})^{-1}}{e_{\pp}}\\
=\Np Q(\vartheta)^{e_{\pp}^2}\cdot\LPP{\Sigma_D}{\omega_{\pp}^{-1}}{e_{\pp}}\cdot\prod_{1\le i\le e_{\pp}}\operatorname{rat}(\lowerPoint{\pp}{i},\vartheta,\pm)^2 .
\end{multline*}
\end{enumerate}
\end{lemma}
\begin{proof}
(a) Both partial $L$-values are non-zero (Input~6), and a quotient of absolute values is real. (b) Unfold the
product over the union of the two disjoint sets of places and write each $|L|^2$ of the twisted character as the
square of the ratio times $|L|^2$ of the untwisted one.
\end{proof}
\leanline{\lean{rationezeroandreal}\sep \lean{lValueProdeq}}

\begin{proposition}[conjugates of the lift]\label{p:galois}
\textup{(a)} Let $P$ be parameters, $i_1$ a natural number (a level of the tower of fields, with field $E_{i_1}$),
$\bn$ a finite index set, $\Sigma_1$ level places of $E_{i_1}$ without twist places, and $\sigma$ an automorphism
of $\CC$ over the quadratic field of the parameters. There are a complex number $\mathrm{t}_1$, whose absolute value
is the inverse square root of the period of the parameters at that level, a complex number $\mathrm{l}_1\ne0$ and
a point $\mathsf{u}_1$ such that for all $L$-value data $\fl$ on $\bn$ with the places $\Sigma_1$ and every
$L$-value datum $\fl'$ with the same places whose jet degrees are those of $\fl$ permuted by the action of $\sigma$
on the infinite places and whose character is the $\sigma$-conjugate of $\omega_{\fl}$ on the points with trivial
infinite component: if $|\bn|$ is even, then for every finite point $\sh$ and every coordinate, the
$\sigma$-conjugate of the coordinate of $\mathrm{per}^{\sum_\tau\mathrm{dg}_{\fl}(\tau)}\trivialLift{\fl}(\sh)$ is
\[\mathrm{t}_1^{\sum_\tau\mathrm{dg}_{\fl}(\tau)}\cdot\mathrm{l}_1\cdot\bigl(\omega_{\fl'}\,\upsilon^{\lfloor|\bn|/2\rfloor}\bigr)(\mathsf{u}_1)
\cdot\frac{\sigma\bigl(\Delta_{\mathrm f}(\sh^{\natural})\bigr)}{\Delta_{\mathrm f}(\sh^{\natural})}\]
times the permuted coordinate of $\trivialLift{\fl'}(\sh)$. Here $\omega_{\fl'}\upsilon^{\lfloor|\bn|/2\rfloor}$ is the
second character of the theta kernel of $\auxPair{\fl'}$ (the first is $1$), evaluated at $\mathsf{u}_1$;
$\sh^{\natural}$ is the finite point $\sx\mapsto\sh\,\sx$ of the unitary group of the $\bn\times1$ matrices, the
finite component of the image of $(1,(1,\sh))$ under the homomorphism of the pair $(\sI_1,\sI_{\bn})$ of
$\auxPair{\fl}$; and for a finite point $\sg$ of a unitary group, $\Delta_{\mathrm f}(\sg)\ne0$ is its multiplier in
the Weil representation of the background library (the multiplier of the invertible matrix of $\sg$), so that the
last factor is the \kw{multiplier ratio} of $\sigma$ at $\sh^{\natural}$.

\textup{(b)} Let $P$ be parameters, $D$ admissible data, $\bn$ a finite index set and $\sigma$ an automorphism of
$\CC$ over the quadratic field of the parameters. There is a real number $\rc$ such that for every pair $\pp$ of
$D$ that has lift rank $\bn$ there is a real number $\mathrm{r}_0$ with $-4f(D)\le\mathrm{r}_0$, with
$\mathrm{r}_0\le4f(D)$ if $\pp$ is flat, and with
\[\log\bigl|\sigma\bigl(\|\trivialLifta{\pp,\bn}{1}\|^2\bigr)\bigr|=\rc+\mathrm{r}_0+\sum_\tau\operatorname{jet}_D(\mu_{\pp}(\tau)-e_{\pp}),\]
where $\operatorname{jet}_D(i)=\log i!-i\log(\text{period}/g(D))$ is the jet size of the data.
\end{proposition}
\begin{proof}
(a) A $\sigma$-semilinear conjugation of data conjugates the jets of the kernel data place by place (Input~7);
apply it term by term to the class sums that define the two lifts, the second translated by a scalar. The
conjugation commutes with the finite Weil representation up to the multiplier ratio, which gives the last factor;
it sends the finite datum of $\auxPair{\fl}$ on the vacuum to $\mathrm{l}_1$ times the conjugate datum translated by
a scalar point, and moving this scalar through the class sum gives the value of the character at $\mathsf{u}_1$:
the square of a multiplier ratio is a quotient of two values of $\upsilon$, and since $|\bn|$ is even the multiplier
ratio of the scalar on the $\bn\times1$ matrices is a power of such a quotient. (b)
Apply (a) to $\sigma$ and to $\sigma$ conjugated by the complex conjugation; the second application gives the
complex conjugate of the $\sigma$-conjugate of the complex conjugate of each coordinate. The two multiplier ratios
at $\sh^{\natural}$ are $\sigma(\Delta_{\mathrm f}(\sh^{\natural}))/\Delta_{\mathrm f}(\sh^{\natural})$ and
$\overline{\sigma(\overline{\Delta_{\mathrm f}(\sh^{\natural})})}/\Delta_{\mathrm f}(\sh^{\natural})$, and the first
times the complex conjugate of the second is
$\sigma(|\Delta_{\mathrm f}(\sh^{\natural})|^2)/|\Delta_{\mathrm f}(\sh^{\natural})|^2=1$, because
$|\Delta_{\mathrm f}(\sh^{\natural})|^2$ is a natural number (the index of a lattice; background library). So the
multiplier ratios cancel in the inner product, and the $\sigma$-conjugate of the squared norm of the scaled lift
is the squared norm of the lift for the conjugate character times the constant of (a) for $\sigma$ and the
complex conjugate of the constant of (a) for the other automorphism; both powers of $\mathrm{t}_1$ have the
absolute value of $|\mathrm{t}_1|^{\sum_\tau\mathrm{dg}_{\fl}(\tau)}$, the two character values have absolute
value $1$, and the two numbers $\mathrm{l}_1$ depend only on $\bn$, $\sigma$ and the places. Insert the formula of
the first note for the norm of the lift for the conjugate character (Input~1), and collect the error $\mathrm{r}_0$: it is the logarithm of the number of elements of the Weyl
index set plus the logarithm of the absolute value of the partial $L$-value; it is at least $-4f(D)$ always, and at
most $4f(D)$ when the pair is flat, where the Weyl index set has one element (second reading). Then collect the factorials of the hook number (Lemma~\ref{l:poly}\,(c)) and the scales into the jet
size.
\end{proof}
\leanline{\lean{maptrivialLift}\sep \lean{lognormmapnormsqtrivialLift}}

\sect{The triangle of the data}{s:tridata}
\begin{definition}[forms of the triangle]\label{d:triforms}
\textup{(a)} Let $E$ be a CM field with a CM type and Weil data, $\sS$ a hermitian space with index set $\bN$ and
$l$ a parameter. $\formNegEquiv{\sS}{l}$ is the linear isometry from the forms on $\sS$ of the type of $l$ to the
forms of the type of $l$ on the opposite space of $\sS$, given by the identification of the two unitary groups
(Lemma~\ref{l:neg}).

\textup{(b)} Let $P$ be parameters, $D$ admissible data, $o$ a flag and $i$, $n$ natural numbers.
$\TriForm{D}{o}{i}{n}$ is the space of forms on the space of the grid $o$ at the node that the triangle places at
$(i,n)$, of the type with the doubled parameter of the grid at that node.

\textup{(c)} For an $m$-fold twist $\vartheta$ of $D$ and $y\in\TriForm{D}{o}{i}{n}$,
$\triLift{D}{o}{\vartheta}{i}{n}{y}\in\TriForm{D}{o}{i}{n+1}$ is the theta lift of $y$ with the characters of the
grid of index $n+i$, the joint harmonic of the two nodes, and the indicator
$\inclusionDatum{\cO_E/\fm^{\triangle}_{o}(\vartheta)}{\cdot}$ of the standard inclusion modulo the modulus
$\fm^{\triangle}_{o}(\vartheta)$ of the triangle. It is linear in $y$.

\textup{(d)} $\triNegEquiv{D}{o}{i}$ is the linear isometric isomorphism from $\TriForm{D}{o}{0}{0}$ to
$\TriForm{D}{o}{i}{0}$ composed of $i$ passages to the opposite space, as in \textup{(a)}.

\textup{(e)} $\triBase{D}{o}$ is the subspace of $\TriForm{D}{o}{0}{0}$ of the forms fixed by the level: the level
of the chosen places if the flag is false, the principal level of the generator of the grid if it is true.

\textup{(f)} For an $m$-fold twist $\vartheta$ and a natural number $n$, the \kw{triangle order}
$\mathrm{ord}_{D,o}(\vartheta,n)$ is the order of the group of invertible matrices, in the rank of the node $(0,n)$
of the grid $o$, over the residue ring of $\fm^{\triangle}_{o}(\vartheta)$. (The formal definition belongs to
another unit and is recalled here.)
\end{definition}
\leanline{\lean{formNegEquiv}\sep \lean{TriForm}\sep \lean{triLift}\sep \lean{triNegEquiv}\sep \lean{triBase}}

\begin{lemma}[the first arrows]\label{l:first}
\textup{(a)} Let $E$ be a CM field with a CM type and Weil data, $\sS$ and $\sT$ hermitian spaces with index sets
$\bN$ and $\bN'$, $\cI$ an index set and $\omega^{\mathrm a}_i$ a family of characters. Suppose given an edge
family for $\omega^{\mathrm a}$ from the type of a parameter $l$ on $\sS$ to a type on $\sT$, in the direction
false; let $i\in\cI$ and let $\dS$ be a complex function on the matrices modulo a modulus $\fm$. Put
$\omega^{\mathrm b}_i=(\omega^{\mathrm a}_i)^{-1}\upsilon^{\mathrm{rank}}$. If there is an edge family for
$\omega^{\mathrm b}$ from the type of $l$ on the opposite space to the same type on $\sT$, in the direction true,
then there is a function $\dS'$ on the matrices modulo $\fm$ such that for all forms $y$ on $\sS$ and $F$ on $\sT$
the inner product of
$F$ with the lift of $y$ (for $\omega^{\mathrm a}_i$ and $\dS$) equals the inner product of the descent of $F$ to
the opposite space (for $\omega^{\mathrm b}_i$ and $\dS'$) with $\formNegEquiv{\sS}{l}(y)$.

\textup{(b)} Let $P$ be parameters, $D$ admissible data, $o$ a flag and $i$ a natural number with $i<m$. There is
a complex number $\rC$ such that for every $m$-fold twist $\vartheta$ and every $y\in\triBase{D}{o}$, writing
$\mathsf{a}(y')=\triLift{D}{o}{\vartheta}{i}{0}{\triNegEquiv{D}{o}{i}(y')}$ for all $y'$ in the ambient space
$\TriForm{D}{o}{0}{0}$: \textup{(1)} every $y''$ in the ambient space with
$\ip{y''}{y'}=\ip{\mathsf{a}(y)}{\mathsf{a}(y')}$ for all $y'$ in the ambient space lies in $\triBase{D}{o}$, and
$y''=0$ only if $y=0$; \textup{(2)} for every $y'$ in the ambient space $\TriForm{D}{o}{0}{0}$,
$\mathrm{ord}_{D,o}(\vartheta,0)\cdot\ip{\mathsf{a}(y')}{\mathsf{a}(y)}$ is the sum of an unconditionally convergent
series of $\rC$ times the coset terms of the inner product formula for the character of the grid, off the places
of $\fm^{\triangle}_{o}(\vartheta)$ and their conjugates, divided by the number of matrices modulo
$\fm^{\triangle}_{o}(\vartheta)$ of the size of the two nodes.
\end{lemma}
\begin{proof}
(a) The theta kernel of the opposite space is a constant multiple of the conjugate kernel, and the Petersson
product of a descent with a form is an integral against the kernel (Input~4); $\dS'$ is the corresponding multiple
of the conjugate datum. (b) The parameters of the two nodes correspond (Input~2), so the joint harmonic is an
edge; the indicator of the inclusion is a sharp datum; the Euler form of the inner product formula (Input~4) gives
the series and the injectivity, and the invariance under the level is seen by translating with lifts of the
elements of the level.
\end{proof}
\leanline{\lean{existsadjointfirstarrow}\sep \lean{hasSuminnertriLift}}

\begin{proof}[Proof of Theorem 3]
Conditions (1) and (2) of a triangle are Lemma~\ref{l:cost}\,(d) and (a). The descents are theta descents whose
point data are, for the first column, the adjoint data of the first arrows given by Lemma~\ref{l:first}\,(a) (the
edge families come from the correspondence of the parameters of the grid, Input~2), and then, column by column,
the next data of Definition~\ref{d:square}\,(c). \emph{Condition (4)} is Lemma~\ref{l:first}: part (a) is the
adjointness on the images of the base, and part (b)\,(1) shows that the descent of the lift of a base form is the
image of a base form that vanishes only with the given one. \emph{Condition (3).} For a square of the triangle,
Theorem~2 applies to the four spaces of its corners, with the family of all $m$-fold twists, the characters of
the grid and the moduli of the triangle; its hypotheses on signs, ranks and weights are those of the grid, and the
tameness is Input~2. It gives a square law over the residue ring of the modulus with constant $\varkappa/\rc$,
and Lemma~\ref{l:next}, applied with the inclusion datum of the lift and the next datum of the descent, turns it
into the commutation of the lift and the descent up to the scalar; the weights of the triangle, the triangle
orders of Definition~\ref{d:triforms}\,(f), are the orders of the groups that appear in Lemma~\ref{l:next}. When the flag is false, the family
$\rc$ costs one square by Lemma~\ref{l:cost}\,(f), since $|\rc|^2$ is a constant times a product of $L$-values of
the twisted character; and the adjoint law follows from the Fourier form of the square law
(Lemma~\ref{l:fourier}, the character being primitive by Lemma~\ref{l:cost}\,(e)) and from the Petersson products
of descents (Input~4), with a family of grade $0$ by Lemma~\ref{l:cost}\,(c).
\end{proof}

\begin{lemma}[the norm of the chain]\label{l:chainnorm}
Let $P$ be parameters, $D$ admissible data, $\vartheta$ an $m$-fold twist and $y'$, $y$ forms of principal level on
$\U(\sY_0(D))$. Let $\mathsf{u}$ be the $m$-th member of
$\weightedChain{\mathrm{ord}_{D,\mathrm{false}}}{\triLift{D}{\mathrm{false}}{\cdot}{\cdot}{\cdot}{\cdot}}{\vartheta}{0}$,
the chain of the $m$ weighted lifts of the triangle from $\sY_0(D)$ to $\sY_m(D)$. Then
$\ip{\Theta^{\mathrm{lev}}_{D,\vartheta}y'}{\Theta^{\mathrm{lev}}_{D,\vartheta}y}$ is the index of the last level
times the number of elements of the Weyl index set of the weight on $\sY_m(D)$ times
$\ip{\mathsf{u}(y')}{\mathsf{u}(y)}$.
\end{lemma}
\begin{proof}
Lemma~\ref{l:petersson} for the standard space $\sY_m(D)$, with Lemma~\ref{l:params}\,(a); the functions attached
to the algebraic forms $\Theta^{\mathrm{lev}}_{D,\vartheta}y$ are the chain of lifts.
\end{proof}
\leanline{\lean{inneriotaAtLevel}}

\begin{proof}[Proof of Theorem 4]
If $\Theta^{\mathrm{lev}}_{D,\vartheta}y=0$ then $\ip{\mathsf{u}(y)}{\mathsf{u}(y)}=0$ by Lemma~\ref{l:chainnorm}, so
the weighted chain of the triangle of Theorem~3 (flag false) vanishes at $y$; by Lemma~\ref{l:chain}\,(b), with
the non-zero triangle orders, $y=0$.
\end{proof}

\sect{Data for the crude rows}{s:crude}
The third note and the unit of the rows use the following vocabulary and three facts on standard data.

\begin{definition}[chains of column lifts; marks]\label{d:crude}
Let $P$ be parameters and $D$ admissible data, with the index sets $\mathsf{X}_i(D)$ of the columns and the weight
of the columns.
\begin{enumerate}[label=\textup{(\alph*)},leftmargin=2.3em,itemsep=1pt]
\item For a natural number $i$, $\homX{i}$ is the hermitian space of homomorphisms between the standard definite
spaces on $\mathsf{X}_i(D)$ and $\mathsf{X}_{i+1}(D)$.
\item For a tame character $\vartheta$ and $i$, $\liftKernela{\vartheta}{i}$ is the theta kernel of the lift
number $i+1$, from forms on $\mathsf{X}_i(D)$ to forms on $\mathsf{X}_{i+1}(D)$ at the levels of the columns, for
the character of that lift twisted by $\vartheta^{i+1}$, with scale $\bar\pi_g\,\mathrm{per}$.
\item A \kw{translation tuple} is a tuple $(\sx_i)$ of finite adelic matrices with rows $\mathsf{X}_{i+1}(D)$ and
columns $\mathsf{X}_i(D)$. For such a tuple $\sx$ and $i$, $\translateDatum{\sx}{i}$ is the normalised vacuum of
$\homX{i}$ translated by the finite Heisenberg operator of $\sx_i$.
\item For a tame character $\vartheta$, a finitely supported complex function $\dS$ on translation tuples, a
family $\dL$ of data and $i$: $\chainComb{\vartheta}{\dS}{\dL}{i}$ is the sum, over the translation tuples
$\sx$, of $\dS(\sx)$ times the chain of the first $i$ lifts with the kernels $\liftKernela{\vartheta}{\cdot}$ and
the data $\dL(\sx,\cdot)$.
\item $\Bmu{\vartheta}{\dS}$ is a map from the columns to algebraic forms on $\sY_m(D)$: restrict the map
$\chainComb{\vartheta}{\dS}{\translateDatum{\cdot}{\cdot}}{m}$ to the columns and compose it with the restriction
to $\sY_m(D)$.
\item Let $\ell$ be a prime with a homomorphism $j$ from $\CC$ to the completed algebraic closure of the
$\ell$-adic numbers. For an infinite place $\sigma_0$ of $E(D)$, $\inducedPlace{\sigma_0}$ is the prime of $E(D)$
above $\ell$ induced by $j$ composed with the embedding of $\sigma_0$. For a finite set $\Sigma_0$ of finite
places: $\IsMarked{\Sigma_0}$ holds if $\Sigma_0$ contains the conjugate of none of its members, $\ell$ is served
if $\Sigma_0$ is non-empty, every member of $\Sigma_0$ is $\inducedPlace{\tau}$ for a top $\tau$, and
$\ell<x_1$ is not split in $K$. Here $\ell$ is \kw{served} in the sense of unit 150 (Input~7): $\chi_K(\ell)=-1$,
$\ell\not\equiv\pm1$ modulo $m$, and the order of the class modulo $N_0$ of the infinite place named by $-\ell$ is
odd and at most the least natural number $n$ with $m\le\ell^{n}$. Further, $\IsFullMark{\Sigma_0}$ holds if $\IsMarked{\Sigma_0}$ holds and for each top
$\tau$, $\Sigma_0$ contains $\inducedPlace{\tau}$ or its conjugate; $\evenPlaces{\Sigma_0}$ is the set of the
infinite places $\sigma_0$ with $\inducedPlace{\sigma_0}\in\Sigma_0$; and for a finite place $\fp$,
$\markExponent{\fp}=a(D)\cdot\#\{\tau\text{ a top}:\inducedPlace{\tau}=\fp\}$.
\item The \kw{integral columns} are the columns whose values have gauge at most $1$ at every point, a module over
the ring of $j$-integral numbers.
\item For $i$ and an infinite place $\tau$, $\lowColumns{i}{\tau}$ is the number of columns of the weight at
$\tau$ of height less than $|\mathsf{X}_{i+1}(D)|$ and different from $1$, and $\lowBoxes{i}{\tau}$ the number of
boxes in these columns.
\end{enumerate}
\end{definition}
\leanline{\lean{homX}\sep \lean{liftKernela}\sep \lean{Translations}\sep \lean{translateDatum}\sep
\lean{chainComb}\sep \lean{B}\sep \lean{inducedPlace}\sep \lean{IsMarked}\sep \lean{IsFullMark}\sep
\lean{evenPlaces}\sep \lean{markExponent}\sep \lean{integralColumns}\sep \lean{lowColumns}\sep \lean{lowBoxes}}

\begin{lemma}[three facts]\label{l:std}
\textup{(a)} Let $P$ be parameters, $D$ admissible data and $i$ a natural number. Then the weight of the columns
goes up from $\mathsf{X}_i(D)$ to $\mathsf{X}_{i+1}(D)$: $|\mathsf{X}_i(D)|<|\mathsf{X}_{i+1}(D)|$, the weight has
no column of height strictly between the two, and the last entry at $\tau$ is
$\lfloor|\mathsf{X}_{i+1}(D)|/2\rfloor$ minus the nominal weight of index $i+1$ plus the number of columns of
height $|\mathsf{X}_i(D)|$.

\textup{(b)} Let $P$ be parameters and $D$ admissible data. The ideal of $\cO_E$ generated by $\bar\pi_g$ is the
product of the two chosen primes that contain it.

\textup{(c)} Let $P$ be parameters, $i_1$ a level of the tower of fields with field $E_{i_1}$, $\bn$, $\bn'$
finite index sets, $\fm$ a modulus of $E_{i_1}$ and $\dS$ a function on $\bn'\times\bn$ matrices over the residue
ring of $\fm$. \textup{(i)} Let $\mathrm{N}_1$, $\mathrm{q}_1$ be natural numbers and $\ell$ a prime with a homomorphism $j$. If
$\mathrm{q}_1\mathrm{N}_1$ lies in the ideal of $\fm$, every value of $\dS$ is integral away from $\mathrm{q}_1$, and $\ell$ is prime to
$\mathrm{q}_1$, then $\mathrm{N}_1^{2f_1|\bn'\times\bn|}\,\mathrm{vs}^{|\bn'\times\bn|}$ times the datum of $\dS$ on the vacuum is a
standard datum at $\ell$ of denominator $\mathrm{N}_1$ ($f_1$ the number of infinite places of $E_{i_1}$). \textup{(ii)} Let $\mathrm{M}_1$ be a
real number with $|\sigma(\dS(\sx))|\le \mathrm{M}_1$ for every automorphism $\sigma$ of $\CC$ and every $\sx$. Then
$\mathrm{vs}^{|\bn'\times\bn|}$ times the datum of $\dS$ on the vacuum is an admissible combination of the finite
Heisenberg translates of the integral vacuum, the sum of the absolute values of its coefficients being at most
the number of matrices times $\mathrm{M}_1$ under every automorphism.
\end{lemma}
\begin{proof}
(a) Arithmetic on the heights of the columns of the weight. (b) Compare the two ideals through the valuations at
the chosen places. (c) The datum of $\dS$ on the vacuum is the sum of the translates of the vacuum with
coefficients the values of $\dS$ (Input~7); after the scaling the coefficients are $\ell$-integral, which is (i),
and they are bounded by $\mathrm{M}_1$ under every automorphism, which is (ii).
\end{proof}
\leanline{\lean{isLiftWeightkappa}\sep \lean{spanpiGBar}\sep \lean{smulvacuumDatummemstandardData}\sep
\lean{isAdmissibleCombvacuumDatum}}

\sect{Notes on the formal text}{s:notes}
\begin{enumerate}[label=\textup{\arabic*.},leftmargin=1.6em,itemsep=5pt]
\item \emph{The exchange theorem.} The formal statement writes the two sides with the data modulo $\fm$; the
integral on the right is over $[\U(\sV')]$ of the theta kernel of the pair $(\sV',\sW)$, so the roles of the two
members of the pair are exchanged with respect to the first note, and Lemma~\ref{l:neg} is what relates the two.
The real datum $\fv$ on the small side is given and the datum $\fx$ on the large side is found. The statement
does not say that $\rC>0$, but it forces it as soon as one tame triple $(\fm,\omega,\chi)$ exists, because
$|\rc|^2=\rC\cdot\LPP{\Sigma(\fm)}{\omega}{e}$ with $\rc\ne0$ and a positive product; the proof gives
$\rC=\rC_0/\rC_1$.
\item \emph{One square.} The number $\varkappa$ is not shown to be non-zero, as the authors' comment says. The
hypothesis on the products of characters is stated for all pairs of ranks that add up to $|\bn_1|+|\bn_2|$.
\item \emph{Three census entries that are not arguments.} For Lemma~\ref{l:neg} the tool that reads the
statements takes the arrow inside the quantification over functions $F$ in the conclusion for a premise, and for
the lift of Definition~\ref{d:triforms}\,(c) the arrow of the type ``linear map'' likewise. In the Lean text
Lemma~\ref{l:neg} has three explicit arguments (the two spaces and the hypothesis on rational points, which has a
default proof), and the lift takes $D$, the flag, $\vartheta$, $i$, $n$ and returns a linear map; here the form
to which the map is applied fills the sixth slot of the notation. In the structured proof of
Lemma~\ref{l:first}\,(b) one step of kind ``case'' covers a placeholder of the Lean syntax and carries no
mathematical content.
\item \emph{Arguments that only Lean can settle.} For a number of declarations of this note the dossier lists the
arguments by Lean's rules applied by hand; no compiler was run. Among the statements above they are Theorems 1,
2, 3 and~5, Lemma~\ref{l:elem}\,(a)--(c), (e)--(g), Lemma~\ref{l:next}, Lemma~\ref{l:greq}, Lemma~\ref{l:pk},
Lemma~\ref{l:planes}\,(a), Lemma~\ref{l:cost}\,(e), (f), Lemma~\ref{l:meas}\,(a), (c), Lemma~\ref{l:ratio}\,(a),
Lemma~\ref{l:first}, Lemma~\ref{l:chainnorm}, Lemma~\ref{l:std}\,(c)\,(ii), the definitions of costs, of integral
measures, of the point datum and the point kernel, of the lift kernel of the columns and of the marks, and two
of the small facts. The doubt concerns implicit and instance arguments: which index sets, signs, primes with
their embeddings and instances a declaration takes. In particular, for the definitions of
Definition~\ref{d:crude}\,(f) that use the embedding of a prime it is not certain from the text whether Lean keeps
the primality of $\ell$ as a separate argument.
\item \emph{Hypotheses with default proofs.} The sign condition and the evenness in
Proposition~\ref{p:partners}, and the hypothesis on rational points in Lemma~\ref{l:neg}, are arguments with a
default tactic; they are hypotheses all the same and are stated as such.
\item \emph{Statements that begin with a local definition.} Theorems 1 and~2 and the small fact on the index of
the last level begin, after their hypotheses, with local definitions (of $\omega^{\circ}$, of the two composites,
of the weights); the dossier restores them, and they were read in the Lean text for Theorems 1 and~2.
\item \emph{Small facts.} The one-step results are listed after Definition~\ref{d:intmeas} without a
statement of their own; the build of this note records them as absorbed. Their formal statements are in the
concordance.
\item \emph{The second reading.} A second reader compared the statements of the unit with the Lean text. Her
points on this note are followed here: the $L$-function in Theorem~1 is that of the base change, without the Euler
factors at the places of the modulus and at their conjugates, while the norms run over the places of the modulus
only; in the proof of Theorem~1 the ring homomorphism goes from the integers of the local product onto the residue
ring of the modulus, and functions are pulled back along it (the dossier had the opposite direction); in
Lemma~\ref{l:first}\,(b) the second form ranges over the ambient space, not over the base; and in
Proposition~\ref{p:galois}\,(b) the error is the logarithm of a Weyl count plus that of an $L$-value.
\item \emph{What is taken for granted.} Everything listed as Inputs; in particular the results of the first
note, the tower property, Landherr's theorem, the inner product formula in its Euler form and the height bounds
for fit families of Chapter~VIII. For the inner product formula, the matrix Gauss sum and the edge families the
second reader found the relation to the cited literature unresolved, and the Inputs say so; none of them is
called a special case of a published theorem here. Proposition~\ref{p:galois} is used by the unit of the rows and
not in this note.
\item \emph{Families over all primes.} In the formal text a fit family is one map from the pairs $(\ell,\iota)$ to
the $m$-fold twists, and a height is a condition on a family indexed by all these pairs: its bound adds the
contributions of the primes. An earlier edition of this note quantified Definition~\ref{d:cost} over one prime at
a time, which is not what the formal text says and is not well formed.
\item \emph{Twists, not tame characters, in the triangle.} The triangle orders, the characters of the grid and the
moduli of the triangle are defined in the formal text for $m$-fold twists only, and Theorem~3 is a statement
about families indexed by them. Theorems 5 and~6 and the definitions of §\ref{s:lvalue} are stated for all tame
characters, as in the formal text.
\end{enumerate}

\subhead{How this note was written.} From the dossier of the unit (its second part). The program check of the
whole dossier covered \UnitSteps{} leaf steps of \UnitProofs{} proofs and found every statement at its lines. I
read all statements of this part in the dossier and the summaries of all its structured proofs, and step by step
the proof of Theorem~1. In the Lean text I read, for the first edition, the statements of Theorems 1, 2 and~3, of
Lemma~\ref{l:neg}, of Lemma~\ref{l:first}\,(b), and the definitions of the square law, of triangles and weighted
chains, of the forms of the triangle and of tameness for a modulus; and for this revision, after a referee's
report found errors in statements that had been written from the dossier alone, the statements of Theorems 5
and~6, of Lemma~\ref{l:first}\,(a), Lemma~\ref{l:cost}, Lemma~\ref{l:meas}\,(c), Lemma~\ref{l:pk} and of
Propositions \ref{p:descent}, \ref{p:partners} and~\ref{p:equal}, the proofs of Proposition~\ref{p:descent} (its
constants), of Proposition~\ref{p:equal} (its spaces) and of Theorem~6 (its integrality), and the definitions of
costs, of heights and fit families, of pairs, of the algebraic part, of integral measures, of the triangle orders
and of marks. The revision corrected: Definition~\ref{d:cost} and Lemma~\ref{l:meas}\,(c) (families over all
pairs $(\ell,\iota)$), Theorem~3 with Definition~\ref{d:triforms} and Lemma~\ref{l:first} ($m$-fold twists),
Proposition~\ref{p:descent} (a family $i(\tau)$, and the constant), the definition of a pair, and the definition
of the triangle orders, which was missing. The other statements are written from the dossier's readings and were
not compared with the Lean text by me; the proofs are outlines that name every input, shorter than the formal
ones, and those of Theorem~2, Theorem~3, Proposition~\ref{p:descent}, Proposition~\ref{p:partners},
Proposition~\ref{p:equal} and Proposition~\ref{p:galois} follow the summaries of structured proofs whose long
automation steps the readers mark as coarse. In Proposition~\ref{p:galois}\,(b) the source of the two-sided bound
for the logarithm of the $L$-value is taken from the dossier's summary and was not traced. Nothing was compiled.

\head{Appendix A. Concordance with the Lean artifact}
{\fontsize{8.6}{11.6}\selectfont\renewcommand{\leanfont}{\ttfamily\fontsize{7.0}{9.6}\selectfont\color{leancol}}
\setlength{\LTleft}{0pt}\setlength{\LTright}{\fill}\setlength{\tabcolsep}{0pt}
\begin{longtable}{@{}p{30mm}p{56mm}p{43mm}r@{}}
\textbf{In this note} & \textbf{Lean declaration} & \textbf{File} & \textbf{Lines}\\\midrule\endhead
\ConcordanceRows
\end{longtable}}

\subhead{Three statements as they are written in the artifact.}
\quoted{exchangetheorem}
\quoted{SquareLaw}
\quoted{existsisTriangle}
Legend: {\leanfont W}, {\leanfont V}, {\leanfont V'} are $\sW$, $\sV$, $\sV'$; {\leanfont σ} is $\se$; {\leanfont ν}
is $\omega$; {\leanfont γ} is $\omega^{\circ}$; {\leanfont 𝔪} is $\fm$; {\leanfont modDatum} is the datum
$\md{\fx}{\dL}{\fm}$; {\leanfont lValueProd} is $\LPP{\Sigma}{\omega}{e}$; {\leanfont IsTameFor} is tameness for a
modulus; {\leanfont z} is the flag $o$; {\leanfont triangleOrder} are the weights $\mathrm{ord}_{D,o}$
(Definition~\ref{d:triforms}\,(f)); the formal text writes {\leanfont Φ}, {\leanfont Ψ} for the partners $\dL$, $\dS$.

\head{Notation}
\begingroup\fontsize{8.6}{10.9}\selectfont\raggedright
\setlength{\columnsep}{9mm}\setlength{\parskip}{0pt}
\begin{multicols}{2}
\nt{$E$, $\cO_E$, $\A$, $\upsilon$; $E(D)$, $K$}{a CM field, its integers, its adeles, the splitting character of Weil data; the field of the data and its quadratic subfield}
\nt{$P$, $D$, $m$, $r$, $t(D)$, $a(D)$, $g(D)$, $f(D)$, $d$, $x$, $x_1$}{parameters, admissible data, the number of lifts, the rank of the rows, the step of the ranks, the order of derivatives, the second level, the number of infinite places, the prime of the parameters, the first two scales}
\nt{$\vartheta$, $\vartheta'$, $Q(\vartheta)$, $\Np Q(\vartheta)$; $\theta$, $\theta_{\ell,\iota}$}{tame characters and twists of $D$, their set of places and its norm; a family of twists indexed by the pairs $(\ell,\iota)$}
\nt{$\tau$; $\sigma$; $\ell$, $\iota$, $j$, $\CC_\ell$; $\pi_g$}{an infinite place; an automorphism of $\CC$; a prime, an isomorphism of $\CC$ with an algebraic closure of the $\ell$-adic numbers, a ring homomorphism from $\CC$ to the completion $\CC_\ell$ of that closure; the integer of norm $g(D)$}
\nt{$\chi_K$, $N_0$, $b$; $f_1$; $y_{\ell,\iota}$}{the quadratic character of $K$, the first parameter, the third scale (handout); the number of infinite places of $E_{i_1}$; a family of complex numbers}
\nt{$\ba,\bb,\bn,\bn',\bn_1,\bn_2,\bN,\bN',\bT,\bp,\br,\br',\bm$; $|\bn|$}{finite index sets; the number of elements}
\nt{$\sR$; $\chi$, $\mathrm{e}_{\sR}$}{a finite commutative ring; a multiplicative character, an additive character}
\nt{$\se$, $\se'$; $\colk{\se}{\sk}$, $\rowk{\se}{\sk}$}{identifications of index sets; the columns of $\sk$ indexed by $\bb$, the rows of $\sk^{-1}$ indexed by $\ba$}
\nt{$\dL$, $\dS$, $\dS'$; $\str{\se}{\chi}{\sz}$, $\stc{\chi}{\sz}$}{functions on the large and on the small side, partners (the formal text writes $\Phi$, $\Psi$); the standard data}
\nt{$\cL$, $\cL'$; $G$; $\rc$, $\rc'$, $\rC$, $\rC_0$, $\rC_1$, $\varkappa$, $A$, $B_1$}{functions on matrices; a function; constants}
\nt{$\SquareLaw{\chi}{\se}{\rc}{\cL}{\cL'}$}{the square law (Definition~\ref{d:square})}
\nt{$\inclusionDatum{\sR}{\sx}$, $\nextDatum{G}{\sy}$, $\blockInlGL{\se}$, $\graphKernel{\chi}{\se}{\sx}{\sy}$}{the inclusion datum, the next datum, the block embedding, the graph kernel (Definition~\ref{d:square})}
\nt{$\sg,\sh,\sk,\sx,\sy,\sz,\sff,\sq$, $\mathsf{c},\mathsf{j},\mathsf{u},\mathsf{v},\mathsf{l},\mathsf{s},\mathsf{w}$, $\sx_0$, $\mathsf{A}$}{matrices, points of groups, vectors and linear maps: $\sx_0$ the inclusion matrix, $\mathsf{l}$ a linear functional, $\mathsf{w}$ a vector of a weight, $\mathsf{u}$ a class or a composite map}
\nt{$\sW,\sW',\sW^{+},\sV,\sV'$, $\sX,\sX^{+},\sX_1,\sX_1^{+}$, $\sU,\sU_1,\sU_2,\sA,\sA',\sS,\sT,\mathsf{J}$; $\sI_{\bn}$}{hermitian spaces ($^{+}$: with one more hyperbolic plane); the standard definite space}
\nt{$\negc(\sW)$, $\mathrm{x}_{\sW}$, $\U(\sW)$, $[\U(\sW)]$, $d\sh$}{the number of negative signs, the adelic frame, the unitary group, its adelic quotient, a Haar measure}
\nt{$\omega$, $\omega'$, $\omega^{\circ}$, $\omega_i$, $\omega^{\mathrm a}_i$, $\omega^{\mathrm b}_i$; $\mu$}{characters of $\U(1)(\A)$; weights}
\nt{$\fm$, $\fm_i$, $\Sigma(\fm)$; $\fm_D(\vartheta,e)$, $\fm^{\triangle}_o(\vartheta)$; $\Sigma_D$, $\Sigma_0$, $\Sigma_1$}{moduli and their places; the moduli of the data and of the triangle; the chosen places, sets of places}
\nt{$\fc$, $\omega_{\fc}$, $\chi_{\fc}$, $\Sigma(\fc)$, $\sF_{\fc}$}{a split tame character with its character, local character, places and local product}
\nt{$\fv,\fx,\fy,\mathfrak{z}$; $f$; $\md{\fx}{\dL}{\fm}$}{real data; a datum; the datum with real part $\fx$ and finite part $\dL$ modulo $\fm$}
\nt{$\theta_{\sW,\sV}$, $\thI{\sW}{\sV}{\omega}{\omega'}{f}$, $\thL{\sW}{\sV}{\fc}{\fv}{\dL}$, $\thS{\sW}{\sV}{\fc}{\fv}{\dS}$}{the theta kernel, the theta integral, the theta integrals on the large and small side}
\nt{$\LPP{\Sigma}{\omega}{e}$, $\LP{\fc}{e}$; $e$, $e_1$, $e_2$, $e'$}{products of partial Hecke $L$-values; natural numbers or integers}
\nt{$L^{\Sigma\cup\bar\Sigma}(s,\omega^{\mathrm{bc}})$; $\omega^{\mathrm{bc}}$; $\bar\Sigma$; $s$}{the Hecke $L$-function of the base change without the Euler factors at $\Sigma$ and at the conjugate places; the base change $\mathsf{x}\mapsto\omega(\mathsf{x}/\bar{\mathsf{x}})$; the set of the conjugates of the places of $\Sigma$; a real number $\ge1$}
\nt{$\eqr{\sV}{\fc}{\fv}$, $\Exch{\sX}{\sV}{\sV'}{\mu}{\se}{\fv}{\fy}{\mathcal{N}}$}{the equal-rank ratio and the standard exchange identity of the first note}
\nt{$\fS$, $\fS'$, $\fS''$; $\mathrm{sl}_{\fS}(0;\fy)$; $\mathsf{f}$; $\rc_{\mathrm{tow}}$, $\rC_{\mathrm{tow}}$, $\rc_{\mathrm{G}}$; $\operatorname{vac}_{\fS}(\sU)$, $\mathrm{u}$}{splittings; the real slice at $0$; the descent map on real data; the complex and the real constant of the tower property, the constant of the Gauss sum; the vacuum factor of the mixed model, a complex number of absolute value $1$}
\nt{$\cI$; $\mathrm{t}$, $\mathrm{t}_1$, $\mathrm{t}_2$, $\mathrm{t}'$, $\mathrm{d}$; $\mathsf{u}_i(\sx)$, $\mathsf{v}_i(\sy)$}{an index set; types of forms and a direction; the two composites of one square}
\nt{$F$, $F'$, $y$, $y'$, $y''$; $\ip{y}{y'}$, $\pet{y}{y'}$; $l$}{forms and elements of inner product spaces; inner products; a parameter of a type}
\nt{$\cF_{i,n}$, $\cV$, $\cV'$; $\mathrm{w}$, $\cM$, $\mathsf{a}$, $\mathsf{d}$, $\mathsf{j}$, $\mathrm{G}$; $\mathrm{b}$, $\mathrm{k}$; $\mathsf{S}$, $\mathsf{T}$}{the spaces, weights, grading, arrows, isometries and condition of a triangle; an index; families of constants; endomorphisms}
\nt{$\IsTriangle{m}{\mathrm{w}}{\cM}{\mathsf{a}}{\mathsf{d}}{\mathsf{j}}{\mathrm{G}}$, $\weightedChain{\mathrm{w}}{\mathsf{a}}{\vartheta}{i}$}{triangles and weighted chains (Definition~\ref{d:tri})}
\nt{$\CostsSquares{D}{n}{\rc}$; $h_0$, $h_1$; $i$, $i'$, $i_0$, $i_1$, $n$, $n'$, $n''$}{costs (Definition~\ref{d:cost}); height bounds; natural numbers and integers}
\nt{$\fl$, $\fl'$, $\mathrm{dg}_{\fl}$, $\omega_{\fl}$; $\LValueData{E}{\bn}$, $\auxPair{\fl}$, $\kernelValue{\fl}{\rc}{f}{\sg}{\sh}$, $\liftKernel{\fl}$, $\kernelForm{\fl}{\sg}$, $\chosenLevel{\fl}$, $\trivialLift{\fl}$}{$L$-value data and the objects of Definition~\ref{d:ldata}}
\nt{$\pp$, $\omega_{\pp}$, $e_{\pp}$, $\mu_{\pp}$; $\HasLiftRank{\pp}{\bn}$, $\IsFlat{\pp}{\bn}$, $\denomBound{\pp}$, $\lValueData{\pp,\bn}{\vartheta}$, $\trivialLifta{\pp,\bn}{\vartheta}$, $\intAlgPart{\pp,\bn}{\vartheta}$, $\pointDatum{\pp,\bn}{\vartheta}$, $\pointKernel{\pp,\bn}{\vartheta}$, $\lowerPoint{\pp}{i}$}{a pair of the data and the objects of Definition~\ref{d:pairs}}
\nt{$\mathrm{vs}$, $\mathrm{per}$, $\mathrm{M}_D$, $\mathrm{ord}_{D,o}(\vartheta,n)$, $\operatorname{jet}_D$, $\operatorname{rat}(\pp,\vartheta,\pm)$}{the vacuum scale, the jet scale, the order of $\prod_{\fp\in\Sigma_D}\GL_{r\cdot t(D)}(\cO_E/\fp)$, the triangle order (the order of the invertible matrices of the rank of the node $(0,n)$ of the grid $o$ over the residue ring of $\fm^{\triangle}_o(\vartheta)$, Definition~\ref{d:triforms}), the jet size, the ratio of $L$-values}
\nt{$\IsIntegralMeasure{\vartheta}{G}$; $\mathsf{u}_i$, $\rc_i$}{integral measures (Definition~\ref{d:intmeas}); their points and coefficients}
\nt{$o$; $\formNegEquiv{\sS}{l}$, $\TriForm{D}{o}{i}{n}$, $\triLift{D}{o}{\vartheta}{i}{n}{y}$, $\triNegEquiv{D}{o}{i}$, $\triBase{D}{o}$}{the flag of a grid; the objects of Definition~\ref{d:triforms}}
\nt{$\sY_i(D)$, $\sW_i(D)$, $\Theta^{\mathrm{lev}}_{D,\vartheta}$; $\mathsf{X}_i(D)$}{the spaces of the chain, the composite of the $m$ lifts at principal level; the index sets of the columns}
\nt{$\homX{i}$, $\liftKernela{\vartheta}{i}$, $\translateDatum{\sx}{i}$, $\chainComb{\vartheta}{\dS}{\dL}{i}$, $\Bmu{\vartheta}{\dS}$}{chains of column lifts (Definition~\ref{d:crude})}
\nt{$\inducedPlace{\sigma_0}$, $\IsMarked{\Sigma_0}$, $\IsFullMark{\Sigma_0}$, $\evenPlaces{\Sigma_0}$, $\markExponent{\fp}$, $\lowColumns{i}{\tau}$, $\lowBoxes{i}{\tau}$; $\sigma_0$, $\fp$}{marks and counts of columns (Definition~\ref{d:crude}); an infinite place, a finite place}
\nt{$\Delta_{\mathrm f}(\sg)$, $\sh^{\natural}$}{the multiplier of a finite point of a unitary group in the Weil representation; the finite point $\sx\mapsto\sh\,\sx$ on the $\bn\times1$ matrices (Proposition~\ref{p:galois})}
\nt{$\mathcal{K}_0$; $\mathrm{d}_0$; $\mathrm{t}_1$, $\mathrm{l}_1$, $\mathsf{u}_1$, $\mathrm{r}_0$; $\mathrm{M}_1$, $\mathrm{N}_1$, $\mathrm{q}_1$; $\mathcal{N}$}{a kernel; a function on matrices; constants of Proposition~\ref{p:galois}; numbers of Lemma~\ref{l:std}; a function on split tame characters}

\end{multicols}
\endgroup

\end{document}
